\documentclass[11pt]{article}

\usepackage[utf8]{inputenc}
\usepackage[T1]{fontenc}
\usepackage[margin=1in]{geometry}
\usepackage{amsmath,amssymb}
\usepackage{graphicx}
\usepackage{longtable,booktabs}
\usepackage{microtype}
% hyperref goes last; it makes \cite clickable through to the bibliography.
\usepackage[colorlinks=true,linkcolor=blue,citecolor=teal,urlcolor=magenta]{hyperref}

% pandoc emits \tightlist but expects the preamble to define it.
\providecommand{\tightlist}{%
  \setlength{\itemsep}{0pt}\setlength{\parskip}{0pt}}

% If the compile times out on Overleaf, comment out \tableofcontents below first.
\title{A Comprehensive Survey of In-Context Learning: Foundations, Mechanisms, Methodologies, Multimodal Expansion, and Applications}
\author{Generated by AutoSurvey}
\date{}

\begin{document}
\maketitle
\tableofcontents
\newpage

\section{Foundations and Theoretical Understanding of In-Context Learning}

\subsection{Formalizing In-Context Learning as Algorithm Learning}

In-context learning (ICL) has emerged as a defining capability of large language models, enabling them to adapt to new tasks by conditioning on input--output demonstrations presented directly in the prompt---without any update to their parameters. This phenomenon is often described as a model performing ``learning at inference time'' \cite{ref1,ref2}. A formal definition frames ICL as the following procedure: a model receives a prompt consisting of a set of example pairs \((x_1, y_1), (x_2, y_2), \dots, (x_n, y_n)\) together with a query input \(x_q\), and it then produces a prediction \(\hat{y}_q\). Crucially, the model's weights remain frozen throughout the process; all adaptation occurs via the model's internal computations triggered by the context. This stands in stark contrast to classical fine-tuning, where gradient-based updates modify the model's parameters. By decoupling learning from weight updates, ICL effectively treats the forward pass of a transformer as an algorithm that infers a hypothesis function from the provided demonstrations, making it a form of meta-learning at scale---a perspective that naturally sets the stage for the Bayesian interpretations explored in the next section.

The meta-learning perspective provides a unified framework for understanding ICL. In meta-learning, or learning-to-learn, a model is trained on a distribution of tasks such that it learns a procedure capable of adapting to new tasks from limited examples. When applied to transformers, the pretraining process often exposes the model to a vast collection of sequences that resemble structured demonstrations, thereby encouraging the model to develop an internal mechanism for extracting and applying task-specific patterns. The work in \cite{ref3} explicitly meta-trains transformers and other black-box models from scratch to become general-purpose in-context learners, demonstrating that the model's forward pass converges to an algorithm that can handle a wide range of problems without any explicit loss function or optimizer specification. This meta-learning interpretation casts the pretraining of language models as a way to amortize the cost of finding a good learning algorithm: the trained transformer embodies a reusable learning rule that can be invoked on the fly by presenting in-context examples. This amortized inference directly prefigures the Bayesian view, where the pretrained model encodes a prior over tasks and the forward pass approximates posterior computation.

Within this meta-learning framework, ICL can be formalized as an \textbf{algorithm learning problem}, where the transformer implicitly constructs a hypothesis function at inference time using the contextual demonstrations. The paper \cite{ref4} frames ICL precisely in this manner, treating the transformer as an algorithm that maps an input sequence of examples and a query to a prediction. Under this abstraction, the model's behavior is governed by the algorithm it has learned during meta-training, and generalization guarantees can be derived by analyzing the stability and sample complexity of that implicit algorithm. The same work connects this formulation to multitask learning, showing that the excess risk in ICL can be bounded by the stability of the algorithm implemented by the transformer, thereby grounding ICL in classical statistical learning theory. The algorithm learning view also clarifies that ICL is not merely a heuristic but a principled computation that, under suitable conditions, can approximate optimal estimators such as ordinary least squares \cite{ref1}.

A growing body of evidence indicates that the implicit algorithms learned by transformers are not arbitrary but often correspond to well-understood optimization and statistical procedures. In the context of linear regression, \cite{ref5} proves by construction that transformers can implement gradient descent and closed-form ridge regression, and it shows empirically that trained in-context learners closely match the predictions of these standard algorithms. More surprisingly, \cite{ref6} demonstrates that transformers can implement higher-order optimization methods such as Iterative Newton's Method, with each transformer layer computing a few Newton iterations, thereby achieving exponentially faster convergence than gradient descent. These findings suggest that the transformer architecture is not just a pattern matcher but a powerful algorithmic engine capable of rediscovering and executing sophisticated learning rules---rules that, as the next section will show, can also be interpreted as implicit Bayesian inference procedures.

The algorithm learning perspective is further enriched by the observation that a single transformer can perform \textbf{in-context algorithm selection}---that is, it can adaptively choose among different base learning algorithms depending on the input sequence. \cite{ref7} provides both theoretical constructions and empirical evidence that transformers can implement a statistician-like process: they can test the nature of the data within the context and then apply the most appropriate algorithm, such as ridge regression for some sequences and Lasso for others, without external prompting. This capability aligns with the idea that the forward pass of a transformer performs a form of meta-reasoning, selecting the right tool for the current task based on the provided demonstrations. Similarly, \cite{ref8} shows that when presented with a teaching sequence that uniquely identifies a function class, transformers can learn two distinct algorithms and dynamically pick the more sample-efficient one, reinforcing the notion of transformers as flexible, adaptive algorithm learners. This adaptive selection mirrors the Bayesian ideal of model averaging, where the posterior implicitly weights different hypotheses according to their consistency with the observed data.

The structural dissection of pretrained transformers further supports the algorithm learning formalization. \cite{ref9} studies synthetic problems where the label depends on the input through a nonlinear representation composed with a linear function. The optimal ICL algorithm in this setting first transforms the inputs via the representation function and then applies linear ICL on top. Trained transformers were found to internally decompose the task: lower layers learned to represent the data while upper layers implemented a linear ICL mechanism. Such a layered implementation mirrors the architecture of a modular algorithm, where early layers preprocess the data and later layers perform the main inference---exactly what one would expect from a model that has meta-learned a specific algorithm. This decomposition into representation learning and inference stages also echoes the hierarchical structure of Bayesian models, where latent variables are first inferred and then used for prediction.

Taken together, these theoretical and empirical contributions converge on a clear formalization: ICL is the process of extracting a hypothesis function from in-context examples through a learned algorithm, with transformers serving as implicit algorithm learners whose internal computations emulate established optimization, statistical, and representational procedures. This perspective not only provides a rigorous foundation for ICL but also connects it to broader themes in meta-learning, multitask generalization, and the emergence of systematic reasoning. It paves the way for analyses of ICL through the lenses of Bayesian inference, PAC learning, and gradient dynamics, which will be explored in subsequent sections. By viewing transformers as learning algorithms that are instantiated at inference time, researchers can more precisely characterize the capabilities, limitations, and inductive biases of ICL, and the algorithm learning view naturally transitions into the Bayesian framework: the learned algorithm can be understood as an efficient approximation to posterior inference, where the forward pass computes an implicit posterior over tasks and the final prediction marginalizes over it---a connection that the following subsection develops in depth.

\subsection{Bayesian Perspectives on In-Context Learning}

Building directly on the algorithm learning framework established in the previous section, the lens of Bayesian inference provides a powerful and coherent framework for understanding in-context learning (ICL), reframing the process not as a sequence of heuristic token associations but as a principled probabilistic computation. Where the algorithm learning perspective treats the transformer as implicitly executing a learned learning rule---such as gradient descent or ridge regression---the Bayesian view reveals that these learned algorithms themselves can be understood as efficient approximations to posterior inference. Under this perspective, the pretrained language model acts as an implicit Bayesian engine that, given a prompt containing a few demonstrations, constructs an approximate posterior distribution over latent tasks or functions and uses it to generate predictions. This subsection examines three interconnected Bayesian interpretations---implicit Bayesian model averaging, latent variable inference, and posterior predictive generation---and discusses how the prior knowledge distilled from pretraining data governs the entire mechanism, thereby providing a probabilistic foundation for the optimization-based accounts explored in the next section.

A foundational Bayesian viewpoint is that ICL implicitly performs Bayesian model averaging (BMA). In this framing, each demonstration pair \((x_i, y_i)\) is generated from an unknown latent concept or function \(f\) that belongs to a family of tasks learned during pretraining. The prompt serves as a set of observations that update the model's belief from a prior distribution over tasks to a posterior, and the prediction for a new query \(x_{\text{query}}\) is the expectation of \(f(x_{\text{query}})\) under this posterior---precisely the kind of inference that the meta-learned algorithms described earlier are designed to approximate. A landmark study \cite{ref10} formally proves that, without any parameter updates, the forward pass of a transformer implicitly approximates the Bayesian model averaging algorithm. The authors show that the attention mechanism parameterizes a posterior inference procedure that, for a latent variable model of the data, computes a combination of the hypotheses consistent with the in-context examples. This perspective yields rigorous performance guarantees: under an online learning setup, perfectly pretrained ICL achieves an \(\mathcal{O}(1/T)\) regret bound, where \(T\) is the number of demonstrations. Moreover, the same work establishes that the approximation error of the transformer decays exponentially with depth while the generalization error decreases sublinearly with the amount of pretraining tokens, providing a unified picture of how the architecture and data scale enable the emergence of a near-optimal Bayesian predictor---effectively characterizing the statistical efficiency of the implicit algorithms that the transformer has learned to execute.

The BMA interpretation directly links ICL to posterior predictive distributions, offering a probabilistic counterpart to the algorithm learning view's emphasis on explicit hypothesis construction. In a Bayesian model, the predictive distribution for a new input given a context \(\mathcal{C}\) is \(p(y \mid x, \mathcal{C}) = \int p(y \mid x, \theta) p(\theta \mid \mathcal{C}) d\theta\), where \(\theta\) represents latent task parameters or the underlying function itself. Large language models (LLMs) trained on vast corpora that implicitly contain many tasks can learn to approximate this integral over the induced posterior. This view is supported by experiments on hierarchical meta-ICL setups \cite{ref2}, where transformers are trained on sequences of input--output pairs drawn from a union of multiple function families. The study finds that when exact Bayesian inference is tractable, high-capacity transformers closely mimic the Bayesian predictor, and even when the model deviates from the exact Bayesian posterior---for instance, when generalizing to unseen function classes---the deviations follow systematic patterns that can be understood as a product of the model's prior inductive biases. Thus, ICL is elegantly characterized as a computation of the posterior predictive distribution conditioned on the prompt, with the pretrained weights encoding a sophisticated prior over both the underlying generative process and the structural relationships among tasks. This Bayesian characterization naturally complements the finding that transformers can implement algorithms like gradient descent or Newton's method: these optimization procedures themselves can be derived as maximum a posteriori or posterior mean estimators under appropriate priors, forging a direct link to the implicit gradient descent perspective developed in the following section.

Complementing the BMA viewpoint, another influential interpretation casts LLMs as latent variable models that perform implicit inference over a compact task representation, paralleling the algorithm selection capabilities discussed in the previous section. In this framework, each prompt induces a latent variable \(z\) that captures the essence of the task, e.g., the intended input--output mapping, the semantic category, or the syntactic regularity. Demonstrations serve as evidence to update the posterior \(p(z \mid \text{prompt})\), and the model's next-token prediction effectively marginalizes over \(z\). This perspective is explicitly adopted in \cite{ref11}, where the authors propose that LLMs naturally infer a latent task variable and show that the optimal selection of demonstrations can be framed as choosing examples that maximize the likelihood of this latent variable. Their algorithm, which selects demonstrations using a smaller LM and then transfers the selection to larger models, yields consistent performance gains across multiple text classification and math reasoning benchmarks, underscoring the empirical validity of the latent-variable Bayesian account. The same logic explains why certain demonstrations are more informative: they reduce the entropy of the posterior over the task, making the model's predictions sharper and more aligned with the true target distribution. This latent-variable perspective mirrors the in-context algorithm selection phenomenon, where transformers adaptively choose among different base learning algorithms depending on the input sequence; from the Bayesian standpoint, algorithm selection is a natural consequence of posterior inference over a discrete set of generative models.

The role of priors---distilled from the pretraining corpus---is central to all Bayesian interpretations and provides the inductive bias that shapes both the implicit algorithms and the optimization trajectories discussed in neighboring sections. The prior distribution over tasks, functions, or latent variables is not explicitly specified but emerges from the massive multitask mixture seen during pretraining. When the data contains diverse and compositional patterns, the resulting prior endows the model with the ability to quickly recognize and generalize within a rich set of task families. \cite{ref2} demonstrates that transformers trained on multiple function classes develop a prior that aligns with the Bayesian ideal when test tasks fall within the training manifold, but they also exhibit the ability to extrapolate to unseen classes, a behavior that departs from strict Bayesian coherence and hints at a learned, structural prior that supports meta-generalization. Meanwhile, \cite{ref10} connects the prior directly to the pretraining data distribution: the ICL estimator's performance and its regret bound are intimately tied to the alignment between the prior induced by the pretraining mixture and the downstream task distribution. This prior-centric view also illuminates why the optimization processes that transformers learn in-context---whether gradient descent, ridge regression, or Newton's method---carry specific inductive biases: the choice of implicit algorithm is itself shaped by the pretraining distribution, with the Bayesian prior effectively selecting which optimization procedure best approximates the posterior for the expected task family.

This Bayesian unification also sheds light on many empirical idiosyncrasies of ICL. The sensitivity to demonstration ordering, label imbalance, and prompt formatting can be understood as artifacts of a posterior that has not converged to the true task due to finite, biased evidence, while the effectiveness of well-chosen examples stems from their ability to quickly concentrate the posterior around the correct latent variable. The prior's influence further explains the ``pull'' of semantic priors: when the pretraining data heavily favors certain input--label associations, the posterior may be dominated by the prior rather than the in-context evidence, especially in subjective or ambiguous tasks. In summary, the Bayesian prism reveals ICL as an emergent inference machine, where the transformer's forward pass computes an approximate posterior over latent tasks and yields predictions that are, in many regimes, simultaneously sample-efficient and theoretically grounded. These interpretations not only deepen our understanding of why ICL works but also provide guidance for improving it---through better example selection, prior calibration, and the design of pretraining data that shapes a more effective inductive bias. Most importantly, the Bayesian framework serves as the conceptual bridge between the algorithm learning perspective, which characterizes what algorithms transformers implement, and the implicit optimization perspective, which explains how those algorithms arise from the mechanics of the forward pass: the learned optimization procedures are precisely the computational mechanisms that realize approximate Bayesian inference at scale.

\subsection{Implicit Gradient Descent and Optimization Duality}

Complementing the Bayesian perspective that casts ICL as an approximate posterior computation, a parallel and equally compelling theoretical interpretation frames in-context learning (ICL) as an implicit optimization process. This subsection examines the formal equivalence between the Transformer forward pass and gradient descent, demonstrating that ICL can be understood as performing meta-optimization---computing parameter updates on the fly without altering stored weights. This perspective not only demystifies the surprising efficacy of ICL but also provides a blueprint for designing novel attention mechanisms inspired by optimization theory.

The foundational insight, articulated by \cite{ref12}, is that language models can be interpreted as meta-optimizers, and ICL as a process of implicit fine-tuning. The authors establish a theoretical duality: the operation of a Transformer attention head, under specific construction of its key, query, and value projections, is mathematically equivalent to a step of gradient descent on a linear layer. In this framework, a GPT-style model first derives ``meta-gradients'' from the provided demonstration examples in the context. These meta-gradients are subsequently applied to the original, ``base'' model's parameters to construct an ICL model tailored to the specific task hinted at by the demonstrations. Crucially, this entire process transpires without ever explicitly updating the model's stored weights. Empirical evidence robustly supports this theory, showing that ICL and explicit fine-tuning exhibit remarkably similar behaviors across various phenomena, including sensitivity to label quality and the effect of the number of training examples. This duality is not merely a theoretical curiosity; it has direct practical implications. Drawing an analogy with the momentum method in classical optimization, the authors designed a momentum-based attention mechanism. By incorporating a persistent, decaying trace of past meta-gradient computations into the attention operation, this variant yielded improved performance over standard attention, reinforcing the validity of the optimization perspective and its potential for guiding future architecture design \cite{ref12}.

Expanding this viewpoint, \cite{ref13} provides a complementary perspective by interpreting the ICL inference process as a gradient descent step realized within a contrastive learning framework. By leveraging kernel methods, this work bridges the gap between gradient descent and the self-attention mechanism under the standard softmax attention setting, moving beyond the simplified linear attention regime often assumed in prior theoretical analyses. The analysis reveals that the model constructs a key-value associative memory that functions like a contrastive learning objective without explicit negative samples. The attention score between a query and a set of keys inherently performs a comparative, discriminative operation that implicitly contrasts the relevant ``positive'' demonstrations against the broader distribution of context representations. This contrastive pattern suggests that the forward pass is effectively optimizing an inner-loop objective to align the test query with the most pertinent demonstration representations, pulling it away from less relevant ones. Recasting ICL through this lens not only deepens its mechanistic understanding but also opens avenues for improving self-attention layers by importing successful techniques from the mature field of contrastive learning, such as novel similarity metrics or hard-negative sampling strategies \cite{ref13}.

The precise nature of the optimization algorithm being implemented in-context has been a subject of intensive study, particularly for linear models. \cite{ref14} and \cite{ref15} provide rigorous analyses of this regime. The former demonstrates that when a single linear self-attention layer is trained with gradient flow on random instances of linear regression tasks, it converges to a global minimum. At this minimum, the Transformer's in-context predictions closely mimic those of the optimal ordinary least squares estimator for the given prompt distribution. This work not only characterizes the dynamics of learning this explicit optimization algorithm but also investigates its robustness, revealing that while the trained Transformer can tolerate certain distribution shifts, it is brittle under covariate shifts---a finding that motivates research into more robust ICL architectures \cite{ref14}. Meanwhile, \cite{ref15} proves that any linear Transformer maintains an implicit linear model and can be seen as executing a variant of preconditioned gradient descent. In a compelling demonstration of emergent algorithmic discovery, when trained on data corrupted with varying noise levels, linear transformers did not merely implement a standard off-the-shelf algorithm; they discovered a novel, highly effective optimization strategy. Reverse-engineering this learned algorithm revealed it to be a sophisticated method incorporating momentum and adaptive rescaling tailored to the noise levels, exceeding the performance of many reasonable hand-designed baselines. This finding elegantly illustrates that even relatively constrained architectures can autonomously invent powerful optimization procedures to solve ICL tasks \cite{ref15}.

While linear cases provide clean tractability, the power of ICL in practice stems from nonlinear learning. The study in \cite{ref16} tackles this frontier by analyzing Transformers with a fully connected layer (MLP) preceding a linear attention layer. This MLP acts as a shared, learned nonlinear feature map, drastically enhancing ICL capability. Analyzing this system through mean-field theory and a two-timescale limit, the authors proved that despite the highly nonconvex infinite-dimensional loss landscape, the optimization dynamics are remarkably benign, with Wasserstein gradient flow almost always avoiding saddle points. This work provides the first rigorous saddle-point analysis for mean-field dynamics, establishing that Transformers can efficiently learn complex nonlinear representations that then enable powerful, algorithm-like ICL in the subsequent attention layer \cite{ref16}.

The convergence behavior of the attention layer itself is deeply tied to optimization principles. \cite{ref17} investigates the implicit bias of gradient descent in training a self-attention layer for binary classification. The study formally proves that, under certain data conditions, the key-query weight matrix converges globally to a hard-margin SVM solution, providing finite-time convergence rates and quantifying the sparsification dynamics of the attention map. This ``implicit bias'' towards maximum-margin solutions directly connects the learning of attention parameters to the structural properties of the ICL algorithms they later mimic. Critically, the analysis demonstrates that adaptive step-size rules, such as normalized gradient descent, can accelerate this convergence, further entrenching the analogy between training attention and optimizing learning algorithms \cite{ref17}. In a related vein, \cite{ref18} delves into the staged-wise learning dynamics of a one-layer softmax attention Transformer for linear function classes. It shows how the model navigates distinct training phases characterized by competing attention weights. For data with imbalanced features, the model learns to predict query tokens of dominant features first before transitioning through four distinct phases to master under-represented features, a process governed by the evolving competition between in-context and in-weights learning mechanisms \cite{ref18}.

Furthermore, the idea that Transformers learn to represent and manipulate causal structures through gradient descent is profoundly related to their optimization capabilities. \cite{ref19} proves that on an ICL task requiring latent causal inference, gradient descent on a simplified two-layer Transformer encodes the latent causal graph directly into its first attention layer. The key insight is that the gradient of the attention matrix explicitly encodes the mutual information between tokens, causing the largest entries to correspond to true causal edges. As a special case for in-context Markov chains, this dynamic directly yields the formation of induction heads, the elementary circuit for ICL. This work elegantly unifies the learning of latent structure with the optimization of an inner algorithm that can leverage that structure for prediction \cite{ref19}.

Finally, the adaptation of attention to the properties of the task itself is a hallmark of its optimization-derived flexibility. \cite{ref20} demonstrates this vividly by showing that in an ICL setting for regression, a softmax attention unit learns a spatial ``window'' to implement a nearest-neighbors predictor. The width of this window adaptively scales with the Lipschitzness of the pretraining task functions and the level of label noise. Critically, this context-dependent adaptivity is a unique property of the softmax activation and cannot be achieved by the often-theoretically-studied linear attention. This study underscores that the softmax operation is not just a normalization choice; it is fundamental to the Transformer's ability to dynamically craft its local optimization algorithm, adjusting its effective learning rate or receptive field based on the inferred complexity of the task at hand \cite{ref20}.

In summary, the implicit gradient descent interpretation casts ICL as a learned optimization process, wherein each forward pass executes a few steps of a task-adapted learning algorithm. This duality not only explains why ICL often mirrors explicit fine-tuning but also opens the door to rigorous sample complexity and generalization analyses, a natural bridge to the PAC-learning frameworks discussed next.

\subsection{PAC and PAC-Bayesian Learnability Frameworks}

Building on the optimization perspective that casts in-context learning as a learned algorithm executed during the forward pass, we now turn to the formal learning-theoretic guarantees that underpin such meta-learning. Probably Approximately Correct (PAC) learning provides a foundational lens for quantifying the sample efficiency and generalization of learning algorithms. In the context of in-context learning (ICL), where a frozen model adapts to a task solely from a prompt containing input--output demonstrations, PAC and PAC-Bayesian frameworks must grapple with a two-tiered generalization problem: the meta-level generalization across tasks observed during pretraining, and the base-level generalization within a single task from the limited in-context examples. This subsection surveys the emerging theoretical treatments that establish PAC-style guarantees for ICL, predominantly using multitask PAC-Bayes bounds, task-environment decompositions, and finite-sample complexity analyses under latent task distributions.

A pioneering formalisation tailored specifically to ICL is introduced in \cite{ref21}. The work proposes a two-phase PAC framework that mirrors the real-world ICL pipeline: an initial pretraining phase fits a function to a distribution of tasks, and a subsequent in-context phase keeps that function frozen while concatenating training examples of a downstream task in its input. The pretraining distribution is modelled as a mixture of latent tasks, a common assumption for natural language data. Under mild assumptions, the framework proves that when the number of pretraining tasks is sufficiently large, the in-context learner can efficiently learn any task from the mixture using a finite number of in-context examples, despite the model weights being unchanged and the input format diverging from the pretraining distribution. The key insight is that in-context learning essentially performs task identification: the model, having observed a rich set of tasks during pretraining, can locate which latent task the prompt belongs to, rather than learning a hypothesis from scratch. This aligns with empirical findings that ICL leverages priors acquired from large-scale data. The paper provides finite sample complexity results that bound both the number of pretraining tasks and the number of in-context demonstrations required to achieve a target error. Such guarantees represent the first PAC-style learnability results for ICL, offering a rigorous underpinning for why few-shot prompts work.

The generalization gap in ICL naturally decomposes into an environment-level gap (related to the finite number of pretraining tasks) and a task-level gap (due to the finite number of in-context examples per task). This decomposition is elegantly captured by PAC-Bayes bounds for meta-learning, as developed in \cite{ref22}. The framework derives new PAC-Bayes bounds that upper-bound arbitrary convex functions linking expected and empirical losses at both the environment and per-task levels. By modelling the meta-learner's posterior over hypotheses as a Gibbs distribution conditioned on the task environment, the bounds separate the meta-generalization error into two terms: one that penalises the discrepancy between the true environment risk and the empirical risk over observed tasks, and another that captures the average per-task generalisation gap. In ICL, the pretrained transformer acts as a meta-learner whose prior is shaped by the pretraining data; the in-context prompt effectively instantiates a task-specific posterior. The PAC-Bayes bounds then quantify how many pretraining tasks and how many in-context shots are needed to ensure that the expected error on a new, unseen task remains small. The numerical experiments in \cite{ref22} demonstrate that their bounds tighten as the base learner adapts more quickly, which resonates with the rapid adaptation characteristic of ICL.

Other works have refined the PAC-Bayes analysis for meta-learning with even tighter guarantees. In \cite{ref23}, the authors combine uniform stability bounds at the base level with PAC-Bayes bounds at the meta level. Because the base learner in ICL (the forward pass through a transformer) operates with high algorithmic stability -- the output distribution does not change drastically when a single in-context example is perturbed -- the resulting compound bound can be significantly tighter than vanilla PAC-Bayes bounds. The result is a bound that explicitly rewards fast adaptation: when the model can achieve low training error with only a few gradient steps (or, analogously, a few in-context examples), the generalisation guarantee improves. This provides theoretical justification for why models that exhibit strong ICL capability also demonstrate high base-level stability across demonstrations.

The assumption of a fixed prior over hypotheses in classical PAC-Bayes can be too conservative when the meta-learner has access to task-related side information. \cite{ref24} addresses this by deriving three novel generalisation error bounds using the relative entropy PAC-Bayes bound and extending them with data-dependent priors based on empirical risk minimisation. In the ICL setting, the prompt itself encodes information about the task, effectively serving as a data-dependent prior that can dramatically shrink the hypothesis search space. The bounds with data-dependent priors exhibit a rapid convergence rate, mirroring the empirical observation that ICL often requires only a handful of demonstrations to achieve near-optimal accuracy. The work confirms that learning a good data-dependent prior from the pretraining tasks is crucial for sample-efficient in-context adaptation.

The PAC-Bayes paradigm further extends to providing robust finite-sample guarantees for few-shot learning with unseen tasks and samples, as shown in \cite{ref25}. This approach models the task-specific posterior implicitly through a generative process, allowing for a richer representation than diagonal Gaussian assumptions. For ICL, where the transformer's internal activations implicitly define a task-specific weight configuration, this implicit posterior perspective directly matches the operational mechanism: the in-context prompt induces a posterior distribution over latent task representations that conditions the subsequent predictions. The resulting PAC-Bayes bounds upper-bound the error on any, even unseen, tasks and samples, providing rigorous generalisation guarantees that complement the empirical success of ICL in low-shot regimes.

Collectively, these theoretical advances establish that in-context learning is not merely an empirical artifact but a learnable paradigm with quantifiable sample complexity. The multitask PAC-Bayes analysis clarifies the roles of pretraining task diversity, number of tasks, and shot count in controlling the generalisation error. It also highlights a fundamental trade-off: while increasing the number of in-context examples can reduce the task-level gap, it must be balanced against the environment-level risk stemming from limited pretraining tasks. The finite-sample bounds from \cite{ref21} concretely demonstrate that, under a latent task mixture model, the number of pretraining tasks needs to scale polynomially with the complexity of the task family, while the required in-context demonstrations scale with the intrinsic dimension of the task identification problem. These insights not only explain why ICL emerges only after pretraining on large, diverse corpora, but also offer principled guidelines for designing more data-efficient in-context learners.

Nevertheless, current frameworks remain largely limited to simplified settings: they typically assume independent and identically distributed in-context examples, do not yet account for the full combinatorial complexity of natural language prompts, and do not capture the interplay between in-context and in-weights learning dynamics observed in practice. Extending PAC-Bayes bounds to cover adaptive demonstration selection, chain-of-thought reasoning, and long-context dependencies constitutes an important open direction. Crucially, the PAC guarantees surveyed here consistently highlight that the number of pretraining tasks, their diversity, and the stability of the in-context learner are pivotal determinants of the transfer risk. This realization sets the stage for the following subsection, which delves into a quantitative analysis of multitask generalization, examining precisely how task similarity, task complexity, and the size of the pretraining task set govern the emergence and quality of in-context learning.

\subsection{Multitask Generalization and Task Complexity}

Building on the PAC and PAC-Bayes guarantees discussed in the previous subsection, we now turn to the factors that govern how a model generalizes across the tasks encountered during pretraining. A central promise of in-context learning (ICL) is the ability to adapt to novel tasks from a few demonstrations without parameter updates, which inherently relies on the multitask nature of pretraining: the model must extract a reusable learning algorithm from a distribution of tasks. Understanding the transfer risk therefore requires analyzing the interplay between the number of pretraining tasks, their similarity, the complexity of each task, and the stability of the implicit algorithm implemented by the transformer. The theoretical frameworks surveyed here formalize these dependencies, providing generalization bounds that clarify precisely when and why ICL succeeds.

A foundational abstraction is to view ICL as an algorithm learning problem, where the transformer constructs a hypothesis at inference time from the prompt. \cite{ref4} formalizes this view and derives generalization bounds for ICL when the prompt consists of i.i.d. input--label pairs or a trajectory from a dynamical system. Crucially, the excess risk of the learned hypothesis is related to the stability of the algorithm implemented by the transformer---that is, how much the output hypothesis changes under small perturbations of the prompt. When the attention architecture satisfies a specific stability condition, the in-context learner achieves low transfer risk on unseen tasks. Moreover, the work identifies a clear inductive bias: transfer risk is governed by task complexity and the number of multitask training tasks in a highly predictable manner, forging a direct link between the diversity of the pretraining tasks and the emergent generalization.

A critical quantitative factor is the number of pretraining tasks. For the illustrative setting of linear regression with a Gaussian prior, \cite{ref26} establishes a statistical task complexity bound for a single-layer linear attention model. The result shows that only a modest number of independent tasks suffices for pretraining, and the resulting model closely approximates the Bayes optimal algorithm---optimally tuned ridge regression---yielding near-Bayes optimal risk on unseen tasks. This highlights that the pretraining task distribution need not be exhaustively large; rather, sufficient diversity to cover the underlying task structure is what matters. The bound explicitly characterises the trade-off among the number of pretraining tasks, context length, and problem dimension, offering a precise condition for the emergence of ICL.

Task similarity is another essential dimension that modulates generalization. In meta-learning, tasks are assumed to be drawn from a common environment, and their relatedness dictates the effectiveness of the learned inductive bias. Information-theoretic bounds in \cite{ref27} capture this via novel upper bounds on the meta-generalization gap that depend on the Kullback--Leibler and Jensen--Shannon divergences between per-task data distributions. These bounds shrink as tasks become more similar, and the gain from additional tasks is modulated by how closely related they are. For ICL, this implies that pretraining on highly dissimilar tasks may not yield the same inductive transfer as a set of correlated tasks, because the shared statistical structure is less compressible into a single in-context algorithm.

When source and target task environments differ, the problem becomes transfer meta-learning. \cite{ref28} provides a rigorous treatment of the transfer meta-generalization gap---the difference between meta-training loss and average loss on tasks from a new target environment. The bounds incorporate the meta-environment shift via the KL divergence between source and target data distributions, while PAC-Bayesian and single-draw variants use the log-likelihood ratio between source and target task distributions. These results directly apply to ICL when the pretraining task distribution is misaligned with the test-time distribution, quantifying the cost of environment mismatch. The associated Information Meta-Risk Minimization (IMRM) algorithm, which minimizes the PAC-Bayesian bound, is shown to outperform empirical risk minimization, underscoring the importance of explicitly controlling transfer risk.

The complexity of a task itself can be measured in a model-agnostic way. \cite{ref29} proposes an inductive bias complexity measure that quantifies the total information needed to generalize on a task minus the information provided by the data. This measure scales exponentially with the intrinsic dimensionality over which the model must generalize, so tasks demanding generalization over many dimensions are drastically harder. Applied to few-shot meta-learning, the measure can differentiate difficulty across ICL benchmarks. A higher task complexity demands a stronger inductive bias, which must be supplied by the pretraining data and the transformer's architecture, thus providing a quantitative lens to understand why some tasks are more amenable to ICL than others.

The interplay among task similarity, number of tasks, and stability is elegantly unified by PAC-Bayesian frameworks, as already introduced in \cite{ref22}. The decomposition of the meta-generalization gap into environment-level and task-level gaps directly translates to ICL: the environment-level gap reflects the finite number of pretraining tasks, while the task-level gap stems from the limited in-context examples. The stability condition from \cite{ref4} controls the task-level variance, and task diversity reduces the environment-level gap. Together, these insights reveal that ICL generalization improves when pretraining tasks are numerous, varied, yet drawn from a coherent distribution, and when the transformer's attention mechanism is stable enough to extract a consistent algorithm from them. This holistic picture aligns with the broader theory of multi-task and transfer learning, where generalization error decreases with the number of related tasks provided they share a common structure \cite{ref27,ref30}.

In summary, the theoretical underpinnings of multitask generalization in ICL rest on three pillars: (i) the number and diversity of pretraining tasks, which supply the empirical basis for the implicit algorithm; (ii) the similarity between tasks, which determines how efficiently shared structure is exploited; and (iii) the stability of the transformer's computation, which ensures that small prompt variations do not lead to large deviations in the inferred hypothesis. These PAC-Bayes, information-theoretic, and stability-based insights characterise the inductive biases that govern transfer risk. Crucially, they set the stage for the next subsection, where we will see how the composition of pretraining data---task diversity thresholds, burstiness, compositional structure---concretely implements the theoretical trade-offs identified here and drives the empirical emergence of ICL.

\subsection{Pretraining Data Composition and the Emergence of ICL}

The theoretical abstractions of task diversity, similarity, and stability developed in the previous subsection find their concrete realization in the composition and structure of pretraining data. Rather than being a mere byproduct of scaling model parameters, the ability to learn from demonstrations at inference time fundamentally depends on whether and how the training corpus embodies these principles. A growing body of work has identified specific data properties---task diversity, compositional structure, burstiness, long-range dependencies, and token frequency distributions---that translate the abstract trade-offs into empirical conditions that enable or inhibit ICL. This section synthesizes those findings, highlighting the concept of a task diversity threshold and the interplay of distributional characteristics that drive the emergence of this capability.

\textbf{Task Diversity and the Threshold for Emergence}\\
A central finding is that ICL on fundamentally new tasks requires a sufficient diversity of tasks during pretraining. In a controlled study of linear regression, \cite{ref31} demonstrated a clear task diversity threshold: below this threshold, a pretrained transformer behaves like a Bayesian estimator that uses the narrow pretraining task distribution as a prior, failing to generalize to unseen tasks. Above the threshold, the transformer's behavior shifts to match that of ridge regression, effectively implementing a Gaussian prior over all possible tasks, including those never encountered during pretraining. This deviation from the Bayes optimal estimator with the pretraining distribution as prior is crucial for solving out-of-distribution tasks. The threshold underscores that task diversity, alongside data scale, is a prerequisite for robust ICL. Complementing this empirical evidence, theoretical work on linear attention models for linear regression established that only a small number of independent tasks are needed for effective pretraining, and the learned model approximates the Bayes optimal ridge regression solution \cite{ref26}. In a PAC-based framework, \cite{ref21} showed that when the pretraining distribution is a mixture of latent tasks, in-context learning can efficiently identify and perform these tasks without weight updates, linking the diversity of latent tasks to the learnability of ICL. The diversity coefficient proposed in \cite{ref32} provides a quantitative measure of formal data diversity, confirming that high-quality LLM pretraining corpora are indeed diverse, further supporting the hypothesis that diversity is a driving force behind ICL.

\textbf{Compositional Structure and Parallel Patterns}\\
Natural language data contain rich compositional structures, and these structures appear to be a key ingredient for ICL. \cite{ref33} established that parallel structures---pairs of phrases following similar templates within the same context window---are strongly predictive of ICL ability. When such parallel structures were ablated from pretraining data, ICL accuracy dropped by 51\%, compared to only 2\% from random ablation, and this effect persisted even after excluding simple n-gram repetitions and long-range dependencies. The parallel patterns span diverse linguistic tasks and long distances, acting as a substrate for learning to copy and apply input-output mappings. From a theoretical perspective, \cite{ref34} argues that ICL emerges from the recombination of compositional operations in pretraining data. An information-theoretic bound shows that if the pretraining distribution contains sufficient compositional structure, next-token prediction naturally gives rise to ICL. The same work provided a controlled setup where trained transformers exhibited ICL for a range of tasks, with performance scaling with parameters and data, and probing revealed representations of compositional structure.

\textbf{Burstiness, Long-Range Dependencies, and Token Frequency}\\
Beyond task diversity, the statistical properties of token sequences profoundly influence the development of ICL. \cite{ref35} demonstrated that burstiness (the tendency of tokens to appear in clusters), large dictionaries, and skewed rank-frequency distributions control the trade-off between in-context learning and in-weights learning. In a minimal attention-only network, these properties drove the abrupt emergence of an induction head, the mechanistic core of ICL, which then competed with in-weights learning. The study proposed that the sequential learning of nested logits due to an intrinsic curriculum underpins the sharp phase transition to ICL. Similarly, \cite{ref36} identified that pretraining data subsets most supportive of ICL contain a higher proportion of long-tail tokens and are challenging examples where the information gain from long-range context is below average. This suggests that learning to effectively incorporate difficult long-range context---a skill enforced by bursty and long-tail data---encourages the model to rely on in-context demonstrations rather than memorized patterns.

\textbf{Pretraining Corpus Composition and Curation Effects}\\
The overall composition of the pretraining corpus, including domain mix and data curation strategies, also modulates ICL emergence. \cite{ref37} observed that in-context learning performance depends heavily on the corpus domain source, and that combining multiple corpora can enable ICL even when each corpus alone does not yield ICL. Moreover, low perplexity on a corpus does not guarantee strong ICL, indicating that language modeling loss and ICL ability are not always correlated. This highlights that the data properties driving ICL are distinct from those that minimize next-token prediction loss. Purposeful data construction can further enhance ICL. \cite{ref38} proposed concept-aware training (CoAT), which constructs training scenarios that reward analogical reasoning, leading to improved ICL performance even when training on only a few datasets. This demonstrates that the inductive biases of pretraining data can be deliberately shaped to foster stronger in-context learners.

In summary, the emergence of ICL is intricately linked to the composition of pretraining data---task diversity above a critical threshold, abundant compositional parallel structures, bursty token distributions, and long-range dependencies all contribute to the development of in-context learning circuits. Crucially, these same data properties that enable ICL also embed the semantic priors that will later interact with the in-context demonstrations, setting the stage for the next subsection's examination of the dynamic balance between prior knowledge and input-label mappings. These findings collectively suggest that careful data engineering---balancing diversity, structure, and distributional properties---can be as important as model scaling for unlocking robust ICL capabilities.

\subsection{Semantic Priors versus Input-Label Mappings}

While the previous section characterized the pretraining data properties---task diversity, compositional parallel structures, burstiness, and long-range dependencies---that govern whether in-context learning (ICL) can emerge at all, the inference-time behavior of an ICL-capable model unfolds as a dynamic negotiation between two forces. The model brings to each prompt a rich set of semantic priors distilled from its pretraining distribution, yet it is simultaneously confronted with novel input--label mappings encoded in the demonstrations themselves. A central tension in understanding ICL is how much predictions are driven by entrenched semantic knowledge versus how much the model genuinely learns and utilizes the in-context relationships. Untangling this interplay is critical for predicting when ICL will succeed and for designing demonstrations that steer the model toward desired behavior.

A large body of evidence indicates that LLMs often rely heavily on their semantic priors, sometimes at the expense of the explicit input--label mappings presented in the prompt. This phenomenon is vividly captured by the ``Demonstration Shortcut'' \cite{ref39}, where the model defaults to its prior knowledge about the task or the surface form of the labels instead of faithfully processing the demonstrated relationships. Experiments that deliberately flip the ground-truth labels in demonstrations provide a stark illustration: small language models largely ignore the flipped labels and continue to classify according to their pretrained semantic associations, while larger models can, after a certain scale threshold, override these priors to respect the contradictory demonstrations \cite{ref40}. This emergent ability to suppress semantic priors is itself scale-dependent, suggesting that sufficiently large models can re-weight the influence of prior knowledge in the face of strong contextual evidence.

However, even when large models are able to learn from flipped or semantically-unrelated labels, the pull of prior knowledge remains strong. The study of semantically-unrelated label ICL (SUL-ICL), where labels are replaced with arbitrary tokens (e.g., ``foo'' and ``bar'' instead of ``positive'' and ``negative''), shows that only models beyond a certain size can accurately learn the new mapping purely from in-context examples \cite{ref40}. Yet, even these large models can struggle to fully overcome prediction preferences inherited from pretraining, and ICL does not treat all in-context information equally \cite{ref41}. For instance, a model may learn to use the correct label tokens but still exhibit a bias toward the pretraining-conditioned likelihood of those labels, leading to performance saturation at suboptimal levels, especially in subjective tasks like emotion recognition where the mapping from text to labels is inherently ambiguous and highly variable across annotators \cite{ref42}. In these settings, the model's priors ``ossify'' predictions, and counter-intuitively, larger models exhibit an even stronger pull of prior knowledge, limiting the benefits of additional in-context demonstrations \cite{ref42}.

Quantifying the respective contributions of semantic priors and input-label mappings has led to a set of carefully designed metrics and experimental manipulations. The Label-Correctness Sensitivity and Ground-truth Label Effect Ratio (GLER) explicitly measure how much a model's performance changes when ground-truth labels are corrupted \cite{ref43}. These analyses reveal that the impact of correct input--label mappings is not uniform; it varies significantly depending on factors such as the verbosity of the prompt template and the model size, indicating that the ostensible robustness to label noise in some configurations may be due to the model relying more on semantic priors than on the demonstrated mappings. Similarly, saliency-based investigations that flip ground-truth labels in contrastive demonstrations find that the change in attention patterns is especially pronounced in larger LLMs \cite{ref44}, underscoring that while large models process labels more explicitly, their final decisions remain a complex blend of learned context and internal prior.

The tension between priors and input-label mappings is further illuminated by studies that decompose the overall ICL effect. Decomposing performance improvements into contributions from label space, format, and discrimination, it was observed that demonstrations have a marginal impact on provoking discriminative knowledge; instead, they exert their greatest influence by regulating the label space and response format \cite{ref45}. This suggests that in-context learners may primarily rely on their pre-existing semantic priors for the actual classification decision, while using the demonstrations to calibrate which label words to output. Such a view aligns with the finding that ICL predictions ``almost always depend on in-context labels'' but the models ``struggle to fully overcome prediction preferences acquired from pre-training data'' \cite{ref41}. Thus, ICL can be characterized as a process in which new input-label mappings are superimposed over a strong prior base, but the base often dominates.

The phenomenon of demonstration shortcut---reverting to semantic priors---has prompted the development of mitigation strategies. In-Context Calibration aims to reduce reliance on priors by adjusting the model's predictions based on the confidence it would have assigned before seeing the demonstrations, effectively subtracting out the prior bias \cite{ref39}. Constructing Comparable Demonstrations (CDs) that minimally edit texts to flip labels forces the model to attend to the inter-demonstration contrast rather than to spurious surface patterns, thereby highlighting the true input--label relationship and mitigating demonstration bias \cite{ref46}. Similarly, concept-aware training (CoAT), while primarily a pretraining intervention, equips models to better exploit analogical reasoning, thereby strengthening the ability to learn input-label mappings from demonstrations rather than falling back on shallow priors \cite{ref38}.

In summary, the interplay between semantic priors and input-label mappings lies at the heart of ICL behavior. The very pretraining data properties that enable ICL also embed deep semantic biases that shape inference. While small models are heavily governed by priors, scaling confers an emergent ability to learn from contradictory demonstrations, yet the pull of prior knowledge is never fully eliminated---especially in subjective or out-of-domain tasks. Instructional tuning further amplifies both the use of semantic priors and the capacity to learn new mappings, with a net advantage for prior reliance \cite{ref40}. Understanding this balance not only clarifies the theoretical underpinnings of ICL as an algorithm-learning mechanism but also directly informs practical strategies for prompt design, calibration, and data augmentation, all aiming to tip the scales toward faithful learning of the demonstrated input--label relationships.

\section{Internal Mechanisms and Interpretability}

\subsection{Induction Heads and Core ICL Circuits}

While the view of in-context learning (ICL) as implicit algorithm execution provides a high-level functional account of how transformers adapt to new tasks, a deeper mechanistic understanding requires identifying the specific circuits that implement this capability. Induction heads have emerged as a fundamental building block, serving as the concrete neural machinery that underpins the algorithmic behaviors observed in the previous sections. Defined by a simple match-and-copy operation, an induction head attends to a previous occurrence of the current token and then predicts the token that followed it in the context, enabling the model to complete sequences like \texttt{\cite{ref47}\cite{ref48}\ …\ \cite{ref47}\ \$\textbackslash{}rightarrow\$\ \cite{ref48}} \cite{ref49}. This operation is not merely a superficial pattern; it generalizes to a form of statistical induction, where the head aggregates bigram statistics from the context to infer the most likely next token, effectively implementing an in-context language model over local transition probabilities. In the study of Markov chain sequence modeling, transformers trained on sequences drawn from a random distribution of Markov chains develop what are termed \emph{statistical induction heads} that compute accurate next-token probabilities by leveraging the bigram frequencies observed in the prompt \cite{ref50}. These heads exhibit a multi-phase learning trajectory: after an initial uniform prediction stage, they first learn to use unigram statistics suboptimally, and then undergo a rapid phase transition to the correct bigram solution, mirroring the sudden emergence of ICL observed in larger models.

The formation of induction heads is a complex process that relies on a set of interacting subcircuits, providing the mechanistic substrate for the emergent algorithmic capabilities described earlier. Using a causal, optogenetics-inspired framework that allows clamping of activations during training, researchers have delineated three underlying subcircuits that must ``go right'' for an induction head to crystallize \cite{ref51}. These subcircuits involve (1) the ability to copy information from the previous token position, (2) the generation of a match signal when the current token matches a token earlier in the sequence, and (3) the integration of this match signal to shift the attention distribution toward the token that follows the matched occurrence. The development of these subcircuits is sequential and nested, with each building on the previous one, and their interaction creates the nonlinearities that give rise to the abrupt phase change in loss and ICL capability. This sequential learning of nested logits, driven by an intrinsic curriculum, explains why induction heads appear suddenly rather than gradually \cite{ref35}. The same study also showed that the competition between ICL and in-weights learning can be understood through the lens of induction head formation: once the head emerges, it competes with in-weights circuits, and the eventual dominance of one strategy can lead to transient ICL behavior, where in-context performance peaks and then decays as the model shifts to memorizing solutions.

The emergence of induction heads is intimately tied to the distributional properties of the pretraining data, revealing how the training environment shapes the model's algorithmic repertoire. Properties such as burstiness, long-tail token frequency distributions, and the need to incorporate long-range context have been shown to promote the formation of induction heads. Empirical analysis of pretraining data subsets that support ICL revealed that these subsets contain a higher proportion of rare, long-tail tokens and challenging examples where the information gain from long-range context is below average, forcing the model to develop induction mechanisms to predict the next token accurately \cite{ref36}. Conversely, synthetic data experiments demonstrated that specific distributional properties---like bursty co-occurrence of tokens---can drive the abrupt emergence of induction heads, while a lack of such structure leads to a dominance of in-weights learning \cite{ref35}. These findings highlight that induction heads are not an inevitable consequence of scaling but rather a functional response to the statistical demands of the training data.

Causal evidence linking induction heads to ICL performance is robust at multiple scales, grounding the algorithmic interpretations in concrete architectural components. In small attention-only transformers, targeted ablation of induction heads severely degrades ICL ability, while their presence correlates with the sudden drop in loss that characterizes the onset of ICL \cite{ref49}. The causal framework developed in \cite{ref51} further showed that by clamping the outputs of the identified subcircuits, one can prevent or accelerate the phase change, confirming that these subcircuits are necessary and sufficient for induction head formation. In larger models with MLPs, the evidence is correlational but consistent: induction heads appear at precisely the same point as the sharp increase in in-context learning performance, and their attention patterns match the expected match-copy behavior \cite{ref49}. Beyond simple bigram induction, higher-order \emph{n-gram heads} that extend the induction mechanism to longer contexts have been identified as critical for more complex in-context language learning tasks, such as recognizing and generating strings from regular languages \cite{ref52}. Hard-wiring these heads into recurrent and convolutional models improves their ICL performance and even yields better perplexity on natural language, underscoring the generality of the induction principle.

The dynamics of induction head formation also illuminate the competition between in-context and in-weights learning circuits, a tension that directly mirrors the broader interplay between implicit algorithm execution and memorization. The transient nature of ICL observed in some training regimes can be attributed to the eventual dominance of in-weights solutions that overfit the training data, causing the induction head to be suppressed or bypassed \cite{ref53}. This competition is modulated by regularization: L2 regularization can help preserve the induction head and prevent the decay of ICL, suggesting that the circuit's persistence is sensitive to the training dynamics. Moreover, the interaction between the induction head and other components, such as copy suppression heads that prevent naive repetition, reveals a more intricate ecosystem of attention heads that collectively calibrate the model's predictions \cite{ref54}. In sum, induction heads constitute the core circuit-level implementation of ICL, bridging the gap between the abstract algorithmic capabilities---gradient descent, Bayesian inference, and algorithm selection---and their physical instantiation in transformer parameters. Their formation is a multi-stage process driven by data properties and intrinsic subcircuit curricula, and their emergence dynamics are characterized by abrupt phase changes that can be causally manipulated. These insights form a mechanistic foundation for understanding how transformers learn to learn from context, complementing the functional perspective with a concrete account of the underlying neural operations.

\subsection{In-Context Learning as Implicit Algorithm Execution}

Complementing the circuit-level account of induction heads, a broader functional perspective posits that the forward pass of a transformer implicitly executes a learning algorithm, using the context as training data to construct a predictor without weight updates. This implicit algorithm execution view unifies theoretical constructions, empirical probing, and the observation that transformers can emulate gradient descent, Newton-like methods, ridge regression, and even Bayesian model averaging. By framing ICL as the execution of a learning procedure, this perspective demystifies how frozen models adapt to new tasks and reveals that the model can flexibly select among different algorithms depending on the input statistics---the very algorithmic flexibility that induction heads help implement mechanistically.

The simplest and most widely studied instantiation is gradient descent. Early work on linear self-attention layers established that a single attention head can be constructed to perform one step of gradient descent on a least-squares linear regression objective when the input includes labeled examples \cite{ref55}. In that setting, the global minimizer of the pre-training loss implements a single gradient step, and if the data distribution is non-isotropic, the learned algorithm becomes preconditioned gradient descent. This construction was extended to multi-layer transformers, where it was shown that critical points of the training objective can correspond to \(L\) iterations of preconditioned gradient descent, with the preconditioner adapting to both the input distribution and data inadequacy \cite{ref56}. Similarly, linear transformers are provably performing a variant of preconditioned gradient descent, maintaining an implicit linear model, and they can discover sophisticated strategies like momentum and adaptive rescaling when faced with heteroscedastic noise \cite{ref15}.

Beyond synthetic constructions, rigorous probing of trained transformers on linear regression tasks reveals that their internal dynamics align with gradient descent and ridge regression. Akyürek et al.~showed that trained in-context learners closely match the predictors computed by gradient descent, ridge regression, and exact least-squares regression, and that the late layers non-linearly encode weight vectors and moment matrices \cite{ref5}. The model transitions between these algorithms as depth and dataset noise vary, and for large width and depth, the behavior converges to a Bayesian estimator, indicating that the transformer can implicitly implement Bayesian model averaging. This Bayesian perspective is further supported by the observation that the transformer's outputs approximate the posterior predictive distribution conditioned on the prompt, and the learning rule can be seen as a natural gradient descent on candidate posteriors \cite{ref57}.

The expressiveness of the implicit algorithm goes far beyond first-order methods. Transformers have been shown to learn higher-order optimization, such as iterative Newton's method, which is exponentially faster than gradient descent for ill-conditioned problems. In the context of linear regression, experiments demonstrate that predictions from successive transformer layers closely match different iterations of Newton's method linearly, with each middle layer roughly computing three iterations \cite{ref6}. Theoretically, a transformer with \(k\) layers can implement \(k\) iterations of Newton's method for matrix inversion and logistic regression, achieving an \(\epsilon\) error with only logarithmically many layers. This finding reframes ICL as an implicit second-order optimizer, extending the gradient-descent analogy to more powerful algorithms.

The capability to handle non-linear functions further broadens the algorithmic repertoire. Transformers can implement functional gradient descent, learning non-linear functions in context by performing gradient descent in function space, with the optimal choice of activation depending on the target function class \cite{ref58}. In discrete settings, transformers can learn to implement multiple distinct algorithms for a single task and adaptively select the more sample-efficient one depending on the in-context sequence \cite{ref8}. This adaptive algorithm selection is a hallmark of sophisticated in-context learners.

A particularly compelling demonstration is the ability of a single transformer to perform in-context algorithm selection, akin to a statistician choosing the right estimator based on the data. Bai et al.~proved that transformers can implement base ICL algorithms such as least squares, ridge regression, Lasso, and gradient descent on two-layer networks, and then employ mechanisms like pre-ICL testing and post-ICL validation to select among them without external prompting \cite{ref7}. For example, a transformer can apply Bayes-optimal ICL on noisy linear models with mixed noise levels by internally validating candidate algorithms. This extends the implicit algorithm view from executing a fixed procedure to a meta-algorithmic decision process.

The mesa-optimization framework provides a unifying perspective: the forward pass constructs an internal learning objective and then solves it via optimization. Reverse-engineering of autoregressive transformers trained on sequence modeling uncovered gradient-based mesa-optimization algorithms that drive predictions, and the same learned optimization can be repurposed for few-shot supervised tasks, suggesting that mesa-optimization underlies ICL \cite{ref59}. The explicit design of mesa-layers that solve optimization problems specified in context further reinforces the algorithmic nature of the computation.

Architectural variants that emphasize iterative computation also align with this view. Looped transformers, which share parameters across iterative blocks, achieve performance comparable to standard transformers on data-fitting tasks while using fewer parameters, indicating that iterative structure naturally supports the emulation of iterative algorithms like gradient descent or Newton steps \cite{ref60}. The transient nature of ICL during training, where ICL can emerge and then be replaced by in-weights learning, has been linked to competition between circuits implementing these two different learning strategies \cite{ref53}, and understanding this competition is crucial for stabilizing ICL.

In summary, the interpretation of ICL as implicit algorithm execution is supported by a convergent body of theoretical and empirical evidence. Transformers can implement gradient descent, preconditioned gradient descent, Newton's method, ridge regression, functional gradient descent, and Bayesian posterior inference, and they can dynamically switch between these algorithms. This functional blueprint raises the next question: how does the transformer's internal computation compress the provided demonstrations into a representation that enables such algorithms? As we will see, a key pattern is the emergence of task vectors and a two-stage process in which lower layers extract representations and upper layers act as lightweight learners, bridging the algorithmic perspective with concrete internal mechanisms.

\subsection{Task Vectors and Representational Encoding}

Building on the insight that ICL corresponds to implicit algorithm execution, a key question is how the internal computations of transformers realize these learning algorithms. A prominent pattern is that transformers compress the provided in-context demonstrations into a compact representation that then modulates subsequent predictions. This idea is crystallized in the concept of a \emph{task vector}. Works such as \cite{ref61} demonstrate that for a wide range of models and tasks, the function learned through in-context learning can be expressed as the transformer acting on the query input together with a single vector derived from the entire training set. The model effectively maps the set of demonstration examples to a fixed-dimensional vector \(\theta(S)\) and then uses this vector to condition the forward pass. This finding bridges ICL with standard machine learning, where a training set is compressed into a set of parameters; here, the compression happens in a single forward computation without any weight updates, and the resulting task vector serves as a kind of implicit parameter that steers the model's behavior.

Closely related to the idea of task vectors is the observation that in-context learning often decomposes into two functionally distinct stages: representation learning in the lower layers and a linear or quasi-linear learning process in the upper layers. The study \cite{ref9} provides both theoretical and empirical evidence for this dissection. By constructing synthetic tasks with a compositional structure---where the label depends on an input through a fixed but possibly complex representation function composed with a per-instance linear function---the authors show that the optimal ICL strategy first transforms the raw inputs using the representation function and then applies linear ICL on top of the transformed dataset. They prove that transformers with moderate depth and size can approximately implement such algorithms, and trained models indeed exhibit the expected split: lower layers transform the demonstrations into a representation space, while upper layers execute a form of linear in-context learning. Probing experiments further reveal concrete copying behaviors on both inputs and their representations, and a post-ICL representation selection mechanism that copes with harder mixture settings. This layered decomposition aligns with the notion that the head of the transformer learns an implicit learning algorithm that operates on features extracted by the bottom, capturing the essence of ICL as representation extraction followed by inference.

When moving from synthetic settings with structured tokens to more realistic setups where inputs and outputs appear as separate tokens (unstructured data), the encoding mechanisms become more subtle. The paper \cite{ref62} investigates exactly this scenario in linear regression tasks. It identifies that a transformer with at least two layers of softmax self-attention and a look-ahead mask can successfully learn from prompts where the label token immediately follows its corresponding input token. Positional encoding is found to further improve performance, and multi-head attention with high embedding dimension is beneficial. These results suggest that the architecture itself, through the interplay of attention and positional information, can implicitly recover the pairing structure and build internal representations that support learning a predictor in context. The lower layers effectively encode the unstructured stream into structured internal signals that upper layers can then exploit for linear regression, mirroring the representation-to-linear-learning decomposition observed in compositional tasks.

The representational encoding of tasks is not limited to abstract task vectors; specific token types within natural language prompts can encode task-critical information. The analysis in \cite{ref63} shows that template tokens and stopwords are particularly prone to become task-encoding tokens, i.e., tokens whose representations heavily influence task performance. By ablating the representations of different token types, the study demonstrates that lexical meaning, repetition, and text formatting are key distinguishing characteristics of these tokens. While this work focuses on natural language prompts, it highlights that the compression of demonstrations into a usable internal representation may rely heavily on a small set of structurally salient tokens rather than on a uniform aggregation of all content. This complements the task vector perspective by showing that the model does not necessarily treat all demonstration tokens equally; instead, it leverages patterns in token type and format to build a robust internal encoding of the task.

Taken together, these lines of work paint a coherent picture: transformers implement in-context learning by first encoding the demonstration set into a compact representation---be it an explicit task vector or distributed patterns over token representations---and then using that encoding to modulate the forward computation in a manner that realizes a learning algorithm. The lower layers are predominantly responsible for extracting and structuring the relevant features from raw input-output pairs, while the upper layers act as lightweight learners that apply a learned hypothesis to the test query. This functional division is observed both in controlled settings with compositional representations and in more rudimentary linear regression tasks with unstructured tokenization. The emerging view is that ICL is not a monolithic ``black box'' but a structured process that mirrors classical two-stage predictors: feature extraction followed by a simple learning rule, all implemented within the layers of a single pre-trained transformer. This staged architecture provides a macroscopic functional skeleton for the implicit algorithms examined in the preceding section, and as the next section will show, its operational details can be directly probed through attention maps, query-key geometries, and the dynamics of information selection within the forward pass.

\subsection{Attention Patterns and Interpretability of Internal Mechanisms}

As foreshadowed by the macroscopic functional skeleton of ICL---where lower layers extract features and upper layers execute a learning rule---the operational details of this process can be directly probed through attention maps and the geometry of query-key embeddings. Interpreting how transformers implement ICL thus requires peeling back the layers of attention and feed-forward computations to reveal the internal learning dynamics. A primary window into this dynamics is the analysis of attention distributions over the context, which exposes how the model adapts its reliance on demonstration tokens without weight updates. Studies using gradient-based saliency maps have shown that the attention pattern over demonstration tokens strongly correlates with ICL performance. For instance, \cite{ref64} introduced sensitivity as an unsupervised measure that negatively correlates with accuracy: prompts that induce more peaked, less sensitive attention patterns yield better in-context performance. The underlying gradient-based saliency scores quantify the relevance of each input token, revealing that a well-chosen demonstration set focuses attention on task-relevant context, while poor prompts cause diffuse attention that weakens the in-context signal.

Further evidence that attention dynamically encodes task-specific relationships comes from contrastive demonstrations. \cite{ref44} used saliency maps to compare attention when ground-truth labels are flipped versus correct. The authors observed that label flipping significantly alters saliency maps, particularly in larger models, demonstrating that attention layers are tightly coupled to the input-label mapping. This sensitivity indicates that attention is not merely retrieving semantically similar tokens but is actively capturing the task rule inferred from demonstrations. Visualizing query-key interactions can then reveal how the model separates relevant from irrelevant context, effectively constructing an on-the-fly notion of task similarity.

The geometry of query-key embeddings is central to this adaptive selection. In a regression setting, \cite{ref20} showed that a single attention unit learns to implement a nearest-neighbors predictor whose attention window adapts to the properties of the pretraining functions. The window widens for smoother target functions (low Lipschitzness) and narrows under high label noise, directly reflecting the statistical structure of the task. This adaptivity arises from the softmax activation, which concentrates or spreads attention weights based on query-key similarity scores. The result is a non-parametric estimation mechanism that generalizes to new contexts; linear attention cannot replicate this behavior, underscoring the essential role of the softmax nonlinearity in ICL.

This view of attention as a learning algorithm is deepened by an optimization perspective. \cite{ref13} established a formal equivalence between self-attention and a gradient descent step in contrastive learning. In this framework, softmax-based attention weights correspond to updates from minimizing a contrastive loss without negative samples: the query representation is pulled toward keys of relevant demonstration tokens and pushed away from irrelevant ones. This dual interpretation explains how attention aggregates contextual information and provides a unifying lens for ICL as implicit gradient-based learning. Focused attention maps can then be seen as the convergence of this implicit optimization, maximizing alignment with correct demonstration patterns.

While attention maps reveal the retrieval and weighting mechanism, the feed-forward blocks complete the inference pipeline. They are thought to implement the learned transformation that maps the aggregated context representation to an output prediction. Sensitivity analyses of full models, as in \cite{ref64}, implicitly capture the contribution of feed-forward layers, and the observed performance--sensitivity link suggests that these layers are also modulated by prompt quality. Moreover, iterative forward tuning \cite{ref65} exploits a dual view between attention and gradient descent to perform multiple passes over demonstrations, accumulating meta-gradients in Key-Value matrices while feed-forward processing refines the internal task representation. Thus, as attention performs dynamic selection, feed-forward layers house the transformation rules that determine the final output.

Overall, the interpretability of ICL through attention patterns and associated mechanisms paints a coherent picture: softmax attention adaptively retrieves context information by implementing a near-neighbor or contrastive-gradient operation on extracted features, while feed-forward layers execute the learned mapping. Together, these components enable transformers to realize diverse learning algorithms in-context, from nearest-neighbor prediction to more complex regressions, all within a single forward pass and without any parameter updates. Understanding how such circuits emerge and compete during pretraining is the natural next step, and we now turn to the data and dynamical conditions that give rise to these ICL mechanisms.

\subsection{Data Properties, Phase Transitions, and Emergence Dynamics}

Building on the interpretability of learned attention patterns and the internal mechanisms that implement in-context learning (ICL) discussed in the previous subsection, a natural next question is what conditions in the pretraining pipeline cause these circuits to form in the first place. Understanding how ICL emerges and stabilizes requires a close examination of the pretraining data distribution and the internal learning dynamics of transformers. Rather than being a smooth, gradual acquisition, ICL often appears abruptly, undergoes phase transitions, and can even be transient, all of which are intimately tied to the statistical properties of the training corpus. This subsection synthesizes empirical and theoretical findings on how data-level characteristics---burstiness, task diversity, long-tail token frequency, and the difficulty of long-range dependencies---interact with model architecture to drive the emergence, competition, and eventual fate of ICL circuits.

A central empirical observation is that the distributional properties of language data, such as burstiness (the tendency for tokens to appear in clusters) and highly skewed rank-frequency distributions, act as crucial control knobs for the balance between in-context and in-weights learning. In a minimal attention-only network trained on a simplified dataset, it was shown that these properties directly govern whether ICL emerges and how abruptly it does so \cite{ref35}. Specifically, the emergence of ICL is driven by the formation of an induction head---a circuit that performs a match-and-copy operation, closely related to the attention-based nearest-neighbor mechanisms described earlier---which appears suddenly and then competes with in-weights learning strategies. The burstiness of the training data influences the timing and sharpness of this transition, and the induction head's abrupt emergence is traced to the sequential learning of three nested logits enabled by an intrinsic curriculum. This mechanistic account underscores that data distributions are not just a backdrop but an active ingredient that shapes the developmental trajectory of ICL abilities.

Complementing this view, the concept of a task diversity threshold has been introduced to explain when transformers can solve fundamentally new tasks in-context. When pretraining data contains a limited variety of tasks, the model behaves as a Bayesian estimator with the narrow pretraining task distribution as its prior, failing to generalize to unseen tasks. However, once the diversity of pretraining tasks exceeds a critical threshold, the transformer's behavior shifts qualitatively: it begins to act like ridge regression, corresponding to a Gaussian prior over a much broader space of tasks, and can solve unseen regression problems optimally \cite{ref31}. This threshold effect is a clear signature of a phase transition in the model's inductive bias, directly shaped by the composition of the pretraining data. The transition from a local Bayesian learner to a more general algorithm underscores how data diversity, alongside scale, dictates the breadth of ICL capabilities, mirroring the adaptivity of attention windows to function properties observed in the previous subsection.

The importance of long-tail tokens and the difficulty of long-range context is further highlighted by direct analyses of pretraining data that supports ICL. By identifying a small subset of pretraining data whose continued training significantly boosts ICL performance, researchers found that these supportive examples are characterized not by high domain relevance to downstream tasks, but by a higher concentration of rarely occurring, long-tail tokens and by being challenging examples where the information gain from long-range context is below average \cite{ref36}. This suggests that ICL is encouraged by data that forces the model to leverage distant dependencies and rare tokens, effectively training the circuits responsible for integrating information across a long context. Such data properties are aligned with the conditions that favor induction head formation: the need to match and copy information from earlier in the sequence, which directly reinforces the attention-based retrieval mechanisms that interpretability studies have revealed.

The abrupt emergence of ICL is often visible as a phase change in the loss, coinciding with the formation of induction heads. A detailed mechanistic study using a causal framework revealed that multiple subcircuits must come together for an induction head to form, and that these subcircuits develop in a sequence that yields a phase change \cite{ref51}. The timing of this phase change is dependent on data properties, and the study identified three interacting subcircuits that drive the transition. This work provides a granular view of the ``developmental stages'' of ICL circuits, showing that the sudden emergence is not a mystery but the result of a multi-step process where earlier subcircuits enable later ones, and the overall process is modulated by the statistics of the training data.

While ICL can appear abruptly, it is not necessarily a permanent acquisition. The phenomenon of transient ICL, where in-context capabilities first emerge, then disappear and give way to in-weights learning, has been observed across a range of model sizes and synthetic datasets \cite{ref53}. In these settings, the training loss decreases smoothly even as ICL accuracy rises and then falls, indicating an asymptotic preference for in-weights learning. The transient nature is attributed to competition between ICL and IWL circuits, and the relative strength of these circuits can be influenced by regularization: L2 regularization was found to promote more persistent ICL, avoiding the need for early stopping. This competition illustrates that the learning dynamics are not simply about the acquisition of ICL but about a continual tug-of-war between different learning strategies, with the final outcome dependent on data statistics, model size, and optimization.

The interplay between ICL and in-weights learning is also reflected in human-like curriculum effects observed in neural networks. In tasks with rule-like structure, ICL exhibits a blocking advantage---learning is better when related examples are grouped---whereas concurrent in-weight learning shows an interleaving advantage on tasks lacking such structure \cite{ref66}. These findings suggest that the two learning modes have complementary strengths and that their competition is sensitive to the structure of the task and the sequence of examples. In the context of emergence dynamics, this competition can lead to plateaus or shifts in performance as one mode dominates the other, depending on the current state of the training data distribution and the model's internal circuits.

In summary, the abrupt emergence, phase transitions, and transience of ICL are not arbitrary but are tightly coupled to pretraining data properties. Burstiness and long-tail token frequency encourage the formation of induction heads, task diversity determines whether the model generalizes beyond its pretraining prior, and the difficulty of long-range dependencies promotes the use of contextual information. These factors collectively shape the developmental trajectory of ICL circuits, leading to sharp phase changes and, in some cases, a transient reign of ICL before in-weights solutions take over. Understanding this competition between learning circuits is a fundamental prerequisite that future work must further unravel to design pretraining recipes that yield robust, persistent, and generalizable in-context learners---a theme that carries forward into the subsequent discussion of how these dynamics can be harnessed and optimized in practice.

\section{Methodologies for Enhancing In-Context Learning}

\subsection{Demonstration Selection and Retrieval}

The selection of informative demonstrations is a cornerstone of effective in-context learning (ICL), as the quality and composition of the prompt directly impact a language model's ability to infer the underlying task. Early ICL explorations often relied on fixed or randomly sampled examples, but subsequent research has shown that performance is highly sensitive to the chosen demonstrations, motivating a rich body of work on principled selection strategies. This subsection surveys key methodologies, spanning retrieval-based similarity measures, influence functions, coverage-driven set optimization, and approaches that leverage label ambiguity and comparability to counteract demonstration bias.

A dominant paradigm for demonstration selection is retrieval-based, where examples are chosen based on their similarity to the test input. Naive lexical overlap, as measured by BM25, has proven surprisingly effective. For instance, \cite{ref67} demonstrated that even simple word-overlap metrics like BM25 outperform random selection, and that retrieval-based ICL can be successfully extended to instruction-finetuned LLMs and Chain-of-Thought prompting. The same work further shows that training a task-specific dense retriever on the available demonstration pool yields additional gains. Beyond BM25, semantic similarity using dense embeddings offers a more nuanced matching. \cite{ref68} proposed InfoScore, a novel metric that leverages the language model's own feedback to evaluate individual example informativeness. Their LENS framework first filters the dataset using this metric and then applies a diversity-guided search to select examples that collectively depict the task, explicitly addressing the interdependence among examples.

A critical realization is that independently selecting the top-\(k\) most similar demonstrations can introduce redundancy while omitting salient information. The BERTScore-Recall (BSR) metric was identified in \cite{ref69} as a superior alternative that selects examples demonstrating more of the test input's reasoning patterns. The authors further developed set-level metrics that optimize coverage, such as Set-BSR, which casts selection as a subset optimization problem. On a suite of compositional tasks, Set-BSR outperformed independent ranking by up to 17 points on average, surpassing even task- or LLM-specific trained selectors while remaining entirely training-free. This coverage-driven perspective highlights the importance of considering the collective information content of the demonstration set rather than treating each example in isolation.

A related line of work models the interaction among demonstrations and the test input through determinantal point processes (DPPs). \cite{ref70} introduced CEIL, which uses DPPs to capture both the relevance of each example to the input and the diversity among selected examples. The DPP kernel is optimized through a contrastive learning objective that distills the LM's preferences, resulting in subsets that are both representative and complementary. CEIL achieved state-of-the-art performance across a wide range of NLP tasks, including classification and generation, and exhibited strong transferability across models and datasets.

Influence functions offer an alternative lens by quantifying how much each demonstration contributes to the model's prediction. \cite{ref71} proposed an influence-based selection method that computes in-context influences directly from few-shot ICL performance. By measuring the effect of including or excluding an example, the method can identify both highly positive and highly negative demonstrations. The analysis revealed up to a 16.3\% performance gap between the best and worst examples, underscoring the fragility of ICL to suboptimal selection. Furthermore, the framework was applied to study recency bias in example ordering, showing that influence scores capture positional effects.

Sequential selection approaches explicitly model the order in which examples are presented. \cite{ref72} formulated sequential example selection as a Markov decision process and trained a retriever using reinforcement learning. This allows the model to account for dependencies between examples and to construct sequences that guide the LM through a reasoning strategy. Evaluated on math word problems and scientific question answering, RetICL consistently matched or outperformed heuristic and learnable baselines, and case studies revealed that it implicitly learns representations of problem-solving strategies, demonstrating the benefit of treating selection as a sequential decision problem.

Addressing label ambiguity is another fruitful direction. \cite{ref73} hypothesized that an effective demonstration set should help resolve the inherent label ambiguity surrounding a test instance. By leveraging the LM's existing knowledge, the method prioritizes examples that the model previously misclassified and that lie on the decision boundary of the test input. This strategy not only uses semantic similarity but also ensures that the selected demonstrations challenge the model's uncertainties, leading to consistent improvements on text classification tasks. Similarly, \cite{ref46} tackled demonstration bias by constructing Comparable Demonstrations (CDs) through minimal text edits that flip labels. These CDs force the LM to attend to task-essential features by contrasting near-identical inputs with opposite labels, thus mitigating spurious correlations. CDs proved especially effective in out-of-distribution settings, where standard similarity-based selection often fails.

The relationship between the LM's parametric knowledge and demonstration selection is explored in \cite{ref74}. The authors distinguish `known' examples---those the LM answers correctly from its own knowledge---from `unknown' ones, and find that prompting with unknown examples can degrade performance by inducing hallucination. An optimal set blends both known and unknown information, guiding the model to retrieve relevant parametric knowledge rather than fabricating answers. This perspective bridges retrieval-based selection with the model's internal knowledge state.

For tasks with large label spaces, simply cramming all labels into the context window is infeasible. \cite{ref75} addresses this by using a pre-trained dense retrieval model to give the LLM only a partial view of the label space per inference. By selecting few-shot examples that cover a subset of labels, the approach sets new state-of-the-art results on intent classification and fine-grained sentiment analysis, and analyses confirm that both the similarity of the examples and the correct label correspondence are crucial components.

Finally, a complementary line of work focuses on curating stable demonstration subsets to reduce variance. \cite{ref76} introduced CondAcc and Datamodels, two methods that score training examples individually by their expected contribution to ICL accuracy when combined with random companions. Simply sampling from the resulting stable subsets improves average accuracy by up to 7.7\% over using the full training set, without any retrieval or calibration. Remarkably, these stable subsets are not characterized by content diversity or low perplexity, challenging conventional assumptions about what makes a good demonstration.

In summary, the field has moved from heuristic retrieval to sophisticated set-level optimization, influence-driven scoring, and ambiguity-aware construction. The common thread is a deeper recognition that ICL is not merely about surface-level similarity, but about crafting a demonstration set that communicates the task structure, resolves uncertainty, and complements the model's prior knowledge. The diversity of methods---ranging from training-free BSR-based coverage to reinforcement-learned sequential selection---reflects the multifaceted nature of the selection problem and points toward future integrated frameworks that combine these complementary strengths. Yet, even with a carefully chosen set, the order in which demonstrations appear and the relative weight assigned to each can dramatically alter ICL performance; the next subsection turns to strategies that optimize exactly these aspects, from curriculum ordering and permutation search to reweighting schemes and order-agnostic inference.

\subsection{Demonstration Ordering and Weighting}

Building on the demonstration selection strategies discussed above, the order in which those examples are presented and their relative influence are equally decisive for in-context learning (ICL). Even with a fixed set of carefully chosen demonstrations, the sequence can cause large swings in predictive accuracy, and not all examples contribute equally to the model's understanding of the task. This subsection examines methods that optimize the sequencing or weighting of demonstrations, aiming to extract maximal ICL performance from a given set of few-shot examples. Key strategies include curriculum-based ordering, per-instance permutation search, demonstration reweighting via self-prediction signals or unbiased fine-tuning, and order-agnostic inference frameworks that circumvent ordering sensitivity altogether.

Ordering demonstrations by difficulty is a natural strategy that mimics human learning. Curriculum learning for ICL, formalized as In-Context Curriculum Learning (ICCL), arranges examples so that the model is first exposed to simpler tasks or instances and gradually progresses to more complex ones within the prompt \cite{ref77}. The core idea is that an easy-to-hard sequence stabilizes the model's internal adaptation process, reducing the risk of inference-time confusion. ICCL requires a means to gauge the difficulty of each demonstration; in practice, difficulty can be estimated by the model's own loss on the example, the length of the output, or other signals. Experimental results indicate that ICCL is particularly effective for open-source LLMs that have undergone instruction tuning, as the ability to exploit a curriculum emerges during that tuning stage. Notably, the same work observes that while LLMs benefit from curriculum ordering, their capacity to assess difficulty lags behind human judgment, suggesting that better automated difficulty estimation could further boost performance.

Beyond fixed easy-to-hard sequences, a more ambitious approach is to treat ordering as an instance-specific optimization problem. The self-adaptive ICL framework \cite{ref78} introduces a select-then-rank paradigm: for each test query, the method first selects a subset of highly relevant demonstrations and then searches for the permutation that yields the correct prediction. The ranking component is guided by an information compression criterion that measures how well a particular ordering helps the model compress the task pattern. Empirically, this self-adaptive strategy achieved a 40\% relative improvement over random sampling on eight NLP benchmarks, underscoring the large headroom that remains in optimizing demonstration sequences. Although full permutation search is computationally prohibitive, efficient heuristics and amortized models can approximate the optimal order, making instance-aware ordering a promising direction.

When treating all demonstrations equally is suboptimal because some examples are noisy, redundant, or biased, reweighting provides a fine-grained lever. The Weighted In-Context Learning (WICL) method assigns continuous weights to each demonstration and incorporates these weights into the prompt or the attention computation \cite{ref79}. Since ground-truth weights are unavailable, WICL introduces a masked self-prediction (MSP) score: for a candidate weight configuration, the model is prompted with a masked version of the demonstrations and must predict the masked content; the prediction loss correlates strongly with the actual ICL performance. An efficient beam search over a discretized weight space then finds an approximately optimal weight vector. The resulting weighted demonstrations can be applied either through scaled token embeddings or by modifying the attention bias at specific positions, leading to large gains across text classification tasks.

Another reweighting perspective tackles the bias that arises when the demonstration set is not representative of the target distribution. Reweighted In-Context Learning (RICL) fine-tunes a language model to learn a weight for each input--output pair in the prompt, using an unbiased validation set as a signal \cite{ref80}. The training process optimizes the weights so that the model's predictions on the validation set match the desired unbiased behavior, effectively calibrating the influence of each demonstration. To reduce training cost, a linear approximation variant (LARICL) directly computes weights by solving a linear system derived from the influence of each example on the validation loss. Both RICL and LARICL provide convergence guarantees and substantially outperform unweighted ICL, particularly when the prompt contains imbalanced label distributions or spurious correlations.

A more radical departure from standard few-shot prompting is to eliminate order sensitivity entirely by aggregating information across demonstrations in parallel. Batch-ICL interprets the ICL process through the lens of meta-optimization, treating each demonstration as a source of a meta-gradient that encodes the task-solving direction \cite{ref81}. Instead of feeding all demonstrations in one sequence, the method performs \(N\) separate one-shot forward passes, each using a single demonstration alongside a dummy query to compute a meta-gradient via the self-attention mechanism. These meta-gradients are then summed and applied to the zero-shot forward pass for the real query, yielding a prediction that is order-agnostic because the aggregation operation is commutative. Empirical evaluation shows that Batch-ICL not only matches but often surpasses the best ordering achievable by standard sequential ICL, while being computationally more efficient due to parallelism. A multi-epoch extension further refines the meta-gradients by iteratively revisiting the demonstrations, emulating a form of internal optimization that implicitly explores many permutations without explicit enumeration. This line of work demonstrates that the ordering problem can be circumvented by rethinking the inference algorithm at the architectural level.

While the primary focus of reweighting and ordering is to maximize the utility of a fixed set of demonstrations, related techniques that internalize demonstrations into model parameters, such as Iterative Forward Tuning, can also avoid ordering issues by distilling the entire prompt into a set of meta-gradients that are applied once at test time \cite{ref65}. Although distinct from traditional ordering, these approaches share the overarching goal of rendering ICL robust to prompt arrangement.

In sum, effective demonstration ordering and weighting are crucial for reliable in-context learning. Curriculum ordering offers a simple, psychologically inspired baseline, while instance-level permutation learning and reweighting via self-prediction or unbiased fine-tuning provide substantial improvements. Order-agnostic methods like Batch-ICL point to a future where the sensitivity to demonstration sequence is eliminated altogether, enabling more robust and efficient few-shot adaptation. Combined with the selection strategies of the previous subsection, these techniques define the input construction pipeline, and we next explore how the phrasing of instructions and the structuring of label spaces can further elevate ICL.

\subsection{Prompt Engineering, Instruction Design, and Label Space Manipulation}

Building on the insights from demonstration ordering and weighting discussed above, prompt engineering, instruction design, and the structuring of label spaces form another critical axis for harnessing in-context learning (ICL). While the previous subsection addressed how to optimally arrange and reweight examples, the way those examples are framed, the verbal instructions that accompany them, and the very definition of the output vocabulary equally govern a model's ability to infer and execute a task correctly. This subsection analyzes how these formatting and structural choices collectively shape ICL performance, reliability, and robustness.

A foundational aspect of prompt design is the selection and phrasing of the instruction itself. The systematic evaluation suite InstructEval \cite{ref82} was developed to rigorously assess instruction selection methods across multiple models and tasks. A key finding from this analysis is that carefully curated, manually-written instructions or even simple, task-agnostic instructions frequently outperform automatically-induced instructions. This underscores a critical gap in the generalizability of current automatic instruction-induction techniques and highlights the value of human intuition and clarity in communicating task objectives. The role of instructions can be further amplified by hint-enhanced prompting. The Hint-enhanced In-Context Learning (HICL) paradigm \cite{ref83} explicitly tackles situations where a standard ICL setup causes the model to neglect query-relevant information within demonstrations. HICL first leverages the LLM's reasoning ability to extract pertinent knowledge from the demonstrations, then concatenates this extracted knowledge as an explicit ``hint'' back into the prompt, guiding the model more directly to the correct answer. This demonstrates that how information is synthesized and presented can be as crucial as the information itself.

A deeper understanding of how demonstrations influence the model's output requires decomposing their functional contributions. A comprehensive study reveals three key dimensions: label space, format, and discrimination \cite{ref45}. Counter-intuitively, this analysis finds that demonstrations have only a marginal impact on provoking the model's pre-existing \emph{discriminative} knowledge about the task. Instead, the primary efficacy of ICL lies in its ability to regulate the \emph{label space} and \emph{format} of the response. In essence, the demonstrations function as a sophisticated formatting and naming guide, showing the model which label words to use and how to structure its output, rather than fundamentally teaching it new discriminatory patterns. This insight positions ICL as a mechanism that acts similarly to detailed instructions, aligning the LLM's generation format and lexical choices with the task requirements. Building on this, retrieval, which selects the most semantically similar demonstrations, notably boosts the model's discriminative capability, suggesting a synergistic effect where formatting is learned from direct context and discrimination is sharpened by retrieving relevant knowledge.

The manipulation of the label space itself offers a powerful and efficient lever for improving ICL, especially in multimodal classification settings. For Vision-Language Models (VLMs), the ICL performance for classification often lags behind contrastive learning-based models. To mitigate this without increasing the number of in-context examples, which can exceed a model's context window and processing capacity, methods focus on increasing the knowledge density of each example by manipulating its label space \cite{ref84}. Two effective strategies are Label Distribution Enhancement and Visual Descriptions Enhancement. By enriching how labels are presented---for instance, by providing a distribution over labels or adding detailed visual descriptions of the target class shown in the demonstration image---a single in-context example can convey significantly more information. This approach has been shown to allow a 2-shot VLM to surpass a 4-shot baseline and even outperform the state-of-the-art CLIP model, demonstrating that the structure of the label space is a high-bandwidth channel for task specification.

The quality and sensitivity of a prompt are not static properties but can be empirically analyzed and even used to guide decoding. A comprehensive analysis of prompt sensitivity \cite{ref64} defines sensitivity as a measure of how much the model's output relies on the specific in-context examples. A crucial finding is that sensitivity serves as an unsupervised proxy for model performance, exhibiting a strong negative correlation with accuracy; prompts that induce lower sensitivity lead to more robust and correct predictions. This phenomenon was validated using gradient-based saliency scores to show how different prompts cause the model to distribute its attention differently across the input tokens. Based on this, a sensitivity-aware decoding strategy can be implemented, which penalizes sensitivity during the greedy decoding process. This approach is particularly beneficial when the input contains scarce information, making the model less prone to over-relying on superficial patterns from the demonstrations. Furthermore, the stability and generalization of ICL depend on the robustness of the prompt design. A holistic evaluation of design choices reveals instabilities caused by factors like the precise wording of an instruction and its interaction with the chosen demonstrations \cite{ref85}. While spurious input-label correlations are a known problem in traditional fine-tuning, they form only a minor issue for prompted models. Instead, the interaction between instructions, template formats, and example choices is the primary driver of variance, demanding careful and holistic design.

In cross-modal or structured tasks, prompt design must also bridge representational gaps. For vision-language models, a hierarchical approach to prompt tuning that leverages structured linguistic knowledge can be used to model the complex interconnections between entities and attributes associated with a category \cite{ref86}. By constructing a knowledge graph from category-related descriptions and using relationship-guided attention across different levels of the prompt (low-level for pairwise associations, high-level for global semantics), the model can handle more complex and long-term relationships, greatly improving generalization. Similarly, methodologies like label alignment for multi-modal prompting \cite{ref87} directly address the gap between a pre-trained VL model's label representation space and the specific labels of a downstream task. By dynamically adjusting category embeddings through a hierarchical loss that aligns parameter, feature, and logit spaces, these methods ensure that the prompt's instructions and target labels are coherently integrated with the model's internal representations, significantly boosting few-shot performance.

In sum, effective prompt engineering transcends simple template filling; it is a multifaceted design problem involving the careful orchestration of instructions, the strategic compression and enrichment of the label space, and the maintenance of representational alignment, all aimed at constructing a clear, robust, and information-dense channel for task communication. Importantly, these same design choices not only determine task performance but also interact intimately with the calibration and trustworthiness of ICL predictions---issues that the next subsection explores in depth.

\subsection{Calibration, Bias Mitigation, and Confidence Estimation}

Building directly on the design and formatting choices examined in the previous subsection---which shape how tasks are communicated via instructions, label spaces, and demonstration structures---the calibration and trustworthiness of ICL predictions emerge as equally critical concerns. A comprehensive empirical study of calibration in ICL \cite{ref88} revealed a non-monotonic relationship between the number of in-context examples and calibration quality. As the shot count increases, models initially suffer from heightened miscalibration before eventually improving; calibration errors are particularly pronounced in low-shot regimes. The same work further demonstrated that common usability enhancements, such as fine-tuning and chain-of-thought (CoT) prompting, can exacerbate miscalibration and produce overconfident predictions, along with unreliable natural language explanations. These findings underscore the need for explicit calibration mechanisms when applying ICL in safety-sensitive tasks. Among the post-hoc recalibration techniques examined, a scaling-binning calibrator proved effective in consistently reducing expected calibration error across a range of natural language understanding and reasoning benchmarks \cite{ref88}. This approach first applies a temperature-like scaling to logits and then performs adaptive binning of the scaled scores, adjusting the predicted probabilities to better match empirical accuracy.

Contextual biases arising from prompt formatting, verbalizer choice, and the arrangement of in-context examples can drastically degrade performance. To counteract these, Batch Calibration (BC) was introduced as a simple, inference-only method that controls the contextual bias introduced by the batch of inputs \cite{ref89}. BC operates by calibrating the model output with respect to a learned or estimated bias term derived from the full prompt context, unifying several prior calibration heuristics. It is inherently zero-shot and imposes negligible computational overhead; in the few-shot setting, BC can be extended to learn the contextual bias from a small set of labeled data, yielding state-of-the-art calibration on both natural language understanding and vision-language tasks \cite{ref89}. This approach effectively addresses brittleness to prompt perturbations and provides a unified framework for bias correction without requiring costly retraining or example re-ordering.

Beyond contextual bias, ICL is susceptible to more insidious label biases that cause the model to latch onto spurious correlations in the demonstrations. A systematic taxonomy of label biases in ICL for text classification distinguishes vanilla-label bias, context-label bias, and domain-label bias \cite{ref90}. Vanilla-label bias refers to the model's tendency to favor certain class labels regardless of input; context-label bias is induced by the composition of in-context examples; and domain-label bias---detected and formalized for the first time in that study---causes performance to plummet to near-random levels for many tasks, as the model is heavily biased by the overall label distribution of the target domain. The proposed domain-context calibration estimates the model's inherent label bias by feeding random in-domain words as prompts, then subtracts this bias from the predictions during inference. This simple procedure, which requires no additional training, dramatically improves ICL performance, with gains of up to 37\% Macro-F1 on severely domain-biased tasks for models like GPT-J and GPT-3 \cite{ref90}. The method generalizes across model scales and task instructions, highlighting the prevalence and impact of label biases in ICL.

An orthogonal strategy to improve calibration and fairness is the injection of controlled noise into model parameters. NoisyICL perturbs the parameters of a pretrained language model with small random Gaussian noise at inference time, without any gradient-based tuning \cite{ref91}. This perturbation reduces overconfident predictions and yields both higher accuracy and better-calibrated confidence estimates across a suite of 12 datasets. Further analysis revealed that NoisyICL leads to fairer predictions across different demographic groups and provides more faithful confidence scores that better reflect the model's true uncertainty, acting as an effective, zero-cost calibration method \cite{ref91}. The simplicity of the approach makes it an attractive baseline for combating the prior bias and miscalibration often observed in standard ICL.

Calibration and task performance frequently trade off, especially in low-resource scenarios, where overconfidence is pervasive across both ICL and supervised fine-tuning (SFT). To navigate this trade-off, self-ensembling techniques have been applied to ICL by aggregating predictions from multiple variations of the prompt, such as different example permutations or diverse instructions \cite{ref92}. The study demonstrated that self-ensembling over in-context example variants or prompt phrasing improves calibration significantly while maintaining or even boosting accuracy, without the need for additional training data. The benefits hold for both ICL and SFT regimes, suggesting that aggregating outputs across diverse instantiations of the same task is a robust way to reduce overconfidence and yield more reliable uncertainty estimates in few-shot settings \cite{ref92}.

The interplay between reasoning prompts and calibration adds another layer of complexity. Chain-of-thought prompting, while often improving task performance, tends to make models more overconfident and less calibrated; the generated rationales themselves can exhibit unreliable confidence \cite{ref88}. Consequently, deployers must be cautious when interpreting the confidence of CoT-augmented ICL, and recalibration methods should be applied after reasoning chains to restore faithful uncertainty estimates.

Beyond these dedicated methods, explanation-based calibration has also been explored to tie confidence to the presence of meaningful token attributions. For instance, the CME method leverages model explanations to detect non-inductive attributions and reduces the confidence of predictions that lack strong feature evidence, thereby improving calibration in both in-domain and out-of-domain settings \cite{ref93}. While not designed exclusively for ICL, such principles can be integrated with ICL pipelines to further refine confidence estimates.

In summary, the calibration of in-context learners is challenged by low-shot miscalibration, contextual biases, label artifacts, and the overconfidence induced by reasoning chains. Solutions ranging from scaling-binning and batch calibration to noise perturbation, domain-context bias correction, and self-ensembling provide a growing toolkit for building trustworthy ICL systems. Each technique addresses a distinct facet of the calibration problem, and their combination can yield robust, faithful, and fair predictions even in data-sparse regimes. When high-quality human demonstrations are themselves scarce, the calibration challenge intensifies, naturally motivating the data augmentation and synthetic demonstration generation strategies explored in the next subsection.

\subsection{Data Augmentation and Example Generation}

In the data-sparse regimes highlighted at the end of the previous section, the need for reliable and faithful predictions is frequently compounded by a scarcity of manually curated demonstrations. When high-quality, task-specific exemplars are unavailable, augmenting the prompt with generated examples has become a central strategy for strengthening in-context learning (ICL). This subsection surveys a spectrum of techniques that produce synthetic demonstrations, ranging from fully self-generated input--label pairs to automatic rationale generation and ensemble-based instruction synthesis. These methods relieve the dependency on manually curated exemplar pools while often maintaining or even improving task performance compared to zero-shot and few-shot baselines.

A direct line of work leverages the model's own generative capacity to create pseudo-demonstrations. The Self-ICL framework \cite{ref94} eliminates the need for external demonstration pools by bootstrapping a large language model (LLM) in a zero-shot manner. Starting from a test input, the model first generates a set of pseudo-inputs, then predicts pseudo-labels for them via zero-shot prompting, and finally uses these self-created input--label pairs as in-context demonstrations. Evaluations on challenging BIG-Bench Hard tasks show that Self-ICL consistently outperforms zero-shot baselines and, when combined with chain-of-thought, approaches the accuracy of using real human-annotated demonstrations. Similarly, SG-ICL \cite{ref95} generates demonstrations entirely from the pre-trained language model itself, minimizing reliance on external data. Across four text classification tasks, SG-ICL not only surpasses zero-shot learning but also achieves a performance boost roughly equivalent to 0.6 gold training samples, with notably lower variance than randomly selected real demonstrations. The SEC prompting strategy \cite{ref96} pushes the idea further by replacing hand-crafted examples with a ``self-contemplation'' process in which the model first creates its own demonstrations (and optionally rationales) before producing the final answer. Extensive experiments on arithmetic reasoning, commonsense reasoning, multilingual understanding, and code generation demonstrate that SEC matches or closely approaches few-shot ICL with human-crafted demonstrations, underscoring that contemporary LLMs can rely exclusively on their internal knowledge for decision making.

Beyond generic pseudo-examples, some methods tailor demonstrations to the specific query to bridge distribution gaps. Self-Demos \cite{ref97} addresses out-of-demonstration (OOD) queries by generating query-aware demos that interpolate between the available few-shot exemplars and the test input. This transforms an OOD query into an in-distribution one, significantly outperforming state-of-the-art baselines on a newly constructed tool-using benchmark and on public math reasoning tasks. By eliciting the LLM's inherent generalizability without external retrieval, Self-Demos reduces the fragility that arises when hand-crafted demonstrations fail to cover the query distribution.

When annotated data is scarce, enhancing demonstrations with rationales offers another avenue. For small language models (SLMs), where the benefits of chain-of-thought reasoning are harder to realize, Self-AMPLIFY \cite{ref98} introduces a three-step pipeline: it identifies promising samples using uncertainty or heuristics, generates natural-language rationales via post-hoc explanation methods applied to the SLM itself, and builds the final prompt by integrating these rationales into the demonstration context. On reasoning-intensive datasets, Self-AMPLIFY achieves strong improvements, marking the first method to exploit post-hoc explanations for SLM self-improvement in ICL. Complementary to this, ZARA \cite{ref99} automatically constructs pseudo-parallel rationale--answer pairs for small LMs by reframing the plausibility assessment of a rationale as a natural language inference problem. This zero-shot augmentation approach yields state-of-the-art results on the FEB benchmark for few-shot self-rationalization, validated both by automatic metrics and human evaluation. For larger models, SELF-EXPLAIN \cite{ref100} generates chain-of-thought exemplars without human crafting, inspired by the encoding-specificity principle from memory research. Prompting with these self-explanations makes LLMs more confident, better calibrated, and less biased on complex question-answering tasks, sometimes outperforming human-crafted CoT demonstrations.

Moving to the many-shot regime, where hundreds or thousands of examples can be placed in the extended context window, the need for automated example generation becomes even more acute. The Many-Shot ICL study \cite{ref101} proposes Reinforced ICL, which replaces human-written rationales with model-generated chain-of-thought explanations, and Unsupervised ICL, which strips rationales entirely and prompts only with domain-specific questions. Both variants prove highly effective on complex reasoning tasks, and importantly, many-shot learning with these generated contexts can override pre-training biases and learn high-dimensional numerical functions, capabilities that few-shot ICL struggles to achieve. This demonstrates that synthetic rationales not only alleviate annotation scarcity but also unlock qualitatively different learning behaviors.

Beyond single-demonstration augmentation, ensemble-based generation constructs diverse, high-quality instruction-tuning data that can later serve as in-context exemplars. Ensemble-Instruct \cite{ref102} addresses the challenge of applying Self-Instruct-style pipelines to smaller, permissively licensed LMs (10B--40B parameters). It simplifies ICL templates to ease prompt learning and crucially ensembles outputs from multiple heterogeneous LMs to select high-quality synthetic examples. Using separate pipelines for input-requiring and no-input instructions, Ensemble-Instruct yields instruction-tuning data that significantly boosts both vanilla and instruction-tuned models, and the resulting smaller instruction-tuned LMs can even outperform larger un-tuned counterparts. This line of work shows that a collaborative, multi-model generation process can produce demonstration sets that are richer and more reliable than those from a single model alone.

Additional self-improvement strategies that generate refined demonstrations or hints further enrich the augmentation landscape. AutoHint \cite{ref103} automatically derives hints from input-output pairs where the initial prompt yields incorrect predictions, and then summarizes these per-sample hints back into the original prompt. While aimed at prompt optimization, the enriched instructions effectively serve as augmented demonstrations when combined with few-shot exemplars. Moreover, iterative self-training and self-reflection loops, as explored in AlphaLLM \cite{ref104} and similar approaches, can generate improved reasoning traces that double as high-quality demonstration material.

Taken together, these techniques form a robust toolkit for data-augmented in-context learning. By having models produce their own examples, rationales, or even complete instruction sets, practitioners can overcome the bottleneck of scarce human-annotated demonstrations, reduce sensitivity to demonstration selection, and in many cases achieve performance on par with or superior to manually curated prompts. As these methods scale ICL to larger and richer demonstration sets, the memory and computational overhead of processing thousands of tokens with full self-attention becomes a pressing concern, directly motivating the architectural innovations for efficient inference covered in the next section. This active, model-driven generation paradigm thus represents a significant step toward more autonomous, universally applicable in-context learners.

\subsection{Architectural Innovations and Efficient Inference}

While data augmentation techniques enable the use of increasingly large and diverse demonstration sets, they bring the computational overhead of processing thousands of tokens with full self-attention into sharp focus. In-context learning (ICL) in large language models (LLMs) typically requires the model to attend to a prompt that concatenates a set of demonstrations with a test query, which introduces significant efficiency and scalability challenges. The self-attention mechanism incurs quadratic complexity with respect to the combined length of demonstrations, and the storage of key-value (KV) caches for all context tokens can quickly become prohibitive. Moreover, many standard architectures are limited by fixed positional embeddings that constrain the number of demonstrations, and the need to carry the entire prompt at every inference step magnifies computational cost. To address these bottlenecks, a range of architectural innovations have been proposed, spanning expanded context windows, structured attention variants, cross-attention caching, iterative internalization of demonstrations, and lightweight multimodal extensions. These methods collectively aim to maintain or even improve ICL performance while drastically reducing memory footprint, latency, and dependence on repeated demonstration reprocessing.

The advent of models with vastly expanded context windows has enabled the exploration of many-shot ICL, where hundreds or thousands of examples can be placed in the prompt. The study \cite{ref101} demonstrates that transitioning from few-shot to many-shot regimes yields significant performance gains across generative and discriminative tasks. The paper also introduces Reinforced and Unsupervised ICL, which use model-generated chain-of-thought rationales or entirely domain-specific questions without rationales, respectively, to circumvent the scarcity of human-generated examples. While promising, many-shot ICL exacerbates the quadratic self-attention cost, and as the number of demonstrations grows, the computational demands for both processing and KV caching become severe. This motivates the need for more efficient architectures that can scale to large demonstration sets without sacrificing speed or memory.

To alleviate the quadratic dependency on total context length, structured attention mechanisms have been developed to rewire the self-attention pattern specifically for ICL. The SAICL (Structured Attention for In-Context Learning) framework, described in \cite{ref105}, replaces the dense pairwise attention among all tokens with a structured design that removes unnecessary cross-demonstration dependencies. By allowing each demonstration to be processed independently of others while still attending to the query, SAICL makes the model invariant to the permutation of demonstrations and reduces the complexity to linear in the number of demonstrations. The approach can scale to hundreds of demonstrations, achieving continuous performance improvements with additional examples while delivering up to 3.4× inference speed-up compared to full attention. This structured design directly addresses the demonstration scaling limitation and provides a practical path to leveraging very large demonstration sets.

An orthogonal line of work decouples context conditioning from prompt concatenation entirely by employing cross-attention to cached reference representations. The XC-Cache model, introduced in \cite{ref106}, adapts pre-trained decoder-only LLMs by adding a small number of cross-attention layers that attend to cached states of the context. This eliminates the need to prepend the full prompt for every generated token, drastically reducing the memory footprint of KV caching---by approximately two orders of magnitude compared to standard self-attention caching. On question-answering benchmarks, XC-Cache not only outperforms standard ICL but also reaches performance comparable to fine-tuned models while maintaining the inference-time flexibility of ICL. By drawing inspiration from encoder-decoder architectures, this approach separates the processing of the context from the generation process, enabling memory-efficient conditional generation that is particularly valuable when the relevant context is not known in advance.

Another strategy to avoid storing and reprocessing demonstrations at every inference call is to internalize their information into the model's weights or attention states through an iterative optimization process. The method proposed in \cite{ref65} exploits the dual form between Transformer attention and gradient descent. It introduces a two-stage framework: a ``Deep-Thinking'' stage that iteratively performs forward optimization on the demonstrations, accumulating meta-gradients by manipulating the Key-Value matrices in the self-attention modules. The subsequent inference stage uses only the test query as input, applying the learned meta-gradients through attention without needing to concatenate any demonstrations. This iterative forward tuning effectively stores the knowledge from demonstrations within the modified attention states, eliminating the demonstration storage overhead at inference time. The approach yields substantially better accuracy than standard ICL while significantly reducing the computational cost per inference query, as demonstrations are no longer required as input.

In the multimodal domain, where ICL is often hampered by interleaved modalities and large input lengths, lightweight architectural add-ons have been developed to enhance ICL without incurring prohibitive resource costs. The MultiModal In-conteXt Tuning module, M²IXT, presented in \cite{ref107}, is a compact module that can be prepended to various multimodal unified models such as OFA, Unival, and LLaVA. M²IXT perceives an expandable context window to incorporate labeled examples of multiple modalities (text, image, coordinates) and is trained with a mixed-task strategy, enabling rapid few-shot adaptation across visual question answering, image captioning, visual grounding, and visual entailment. Remarkably, with only about 50K multimodal data, M²IXT boosts few-shot ICL performance by 18\% relative for OFA while remaining approximately 20× smaller than prior multimodal ICL systems like Flamingo or MMICL. Similarly, the Multi-Modal In-Context Tuning method MMICT, described in \cite{ref108}, introduces a Multi-Modal Hub (M-Hub) that captures various multi-modal features according to different inputs and objectives. MMICT enables the model to learn from in-context visual-guided textual features and subsequently generate outputs conditioned on textual-guided visual features, significantly outperforming traditional fine-tuning and vanilla in-context tuning. By designing a variety of in-context demonstrations within this hub, MMICT fully leverages ICL to boost multi-modal fine-tuning while keeping the architectural overhead manageable.

Collectively, these architectural innovations push the boundaries of ICL efficiency and scalability. They address the fundamental limitations of prompt concatenation by rethinking how demonstrations are represented, attended to, and stored. Many-shot ICL exploits larger contexts but demands more efficient attention; structured attention and cross-attention caching directly reduce memory and compute burdens. Iterative forward tuning offers a demonstration-free inference pathway by internalizing knowledge into model states. Lightweight multimodal modules provide a plug-and-play solution to equip existing multimodal models with strong ICL capabilities without massive parameter overhead. These advances are moving ICL from a computationally expensive, scale-limited paradigm toward a practical, scalable inference mechanism that can be deployed across modalities and with vast numbers of examples. However, while these architectural solutions make it possible to efficiently process large volumes of demonstrations, the very ability to learn from those demonstrations in the first place is determined by the data and training recipes employed during pretraining and fine-tuning, which are the focus of the next section.

\subsection{Pretraining and Fine-tuning Recipes for Stronger ICL}

While architectural innovations are crucial for the efficiency and scalability of in-context learning, the very ability to learn from demonstrations is fundamentally shaped by the data and training strategies employed during pretraining and fine-tuning. The development of robust in-context learning (ICL) capabilities is not solely a matter of model scale; the design of pretraining data and fine-tuning protocols plays a decisive role in shaping how well a model can learn from demonstrations. This subsection surveys a range of strategies that explicitly target the pretraining or fine-tuning stage to elicit stronger ICL, from concept-driven data construction to knowledge injection and meta-learning.

A foundational insight is that ICL can be enhanced by training on data that explicitly require analogical reasoning. The concept-aware training (CoAT) framework, introduced in \cite{ref109} and \cite{ref38}, operationalizes this idea by constructing training scenarios in which the model benefits from inferring latent concepts shared across input--output pairs. Instead of treating ICL as a byproduct of multitask scale, CoAT deliberately creates data where the correct response depends on understanding an abstract rule demonstrated in the context. This forces the model to learn to extract and apply analogical mappings from demonstrations. In practice, transformers pretrained with CoAT on only a handful of tasks can match the ICL performance of much larger models trained on hundreds or thousands of tasks, demonstrating that concept-aware data construction is a powerful axis for improving ICL from scratch. The resulting models are also more robust to functional deficiencies commonly observed in standard instruction-tuned models, underscoring the importance of data composition over mere volume.

While CoAT targets the \emph{content} of pretraining data, a complementary line of work addresses the stability of ICL by curating the training examples used during inference. The paper \cite{ref76} shows that carefully selecting a subset of a training set can dramatically reduce the variance of ICL performance without resorting to retrieval or calibration. The authors propose two scoring methods: CondAcc, which evaluates each training example by its average dev-set accuracy when paired with randomly chosen examples, and Datamodels, which learns a linear regressor to estimate how the presence of a specific example influences the model's output. When the top-scoring examples are used as the stable subset, average accuracy improves by 7.7\% (CondAcc) and 6.3\% (Datamodels) across five tasks and two LLMs. Notably, the curated subsets are not necessarily more diverse or lower in perplexity than random selections, challenging the common assumption that diversity and fluency are the primary drivers of ICL quality. This finding suggests that some training examples are inherently more ``supportive'' of the ICL mechanism, and that data curation can be a simple yet effective fine-tuning of the prompt corpus without altering the model itself.

The connection between specific pretraining instances and ICL ability was further explored in \cite{ref36}. Using an iterative, gradient-based approach, the authors identified a small subset of pretraining data that, when used for continued pretraining, can boost a model's ICL accuracy by up to 18\%. The supportive data are characterized by a higher mass of long-tail tokens and contain challenging examples where the information gain from long-range context is below average. This indicates that learning to incorporate difficult long-range dependencies is a key driver of ICL. These findings open the door to actively guiding the construction of pretraining corpora to include more such instances, thereby strengthening ICL from the earliest stages of training.

At the fine-tuning stage, instruction meta-learning has emerged as a powerful recipe for turning a general-purpose pretrained model into a strong in-context learner. The OPT-IML framework \cite{ref110} systematically investigates the decisions that govern instruction tuning. By creating a large benchmark of 2000 NLP tasks and evaluating generalization to held-out categories, tasks, and instances, the study reveals that scaling the diversity and number of instruction-tuning tasks, careful task sampling, and the inclusion of both demonstrations and specialized datasets all contribute to improved ICL. The resulting OPT-IML models (30B and 175B) demonstrate strong zero-shot and few-shot abilities across a wide range of benchmarks, highlighting that fine-tuning with a meta-learning objective over many tasks is an effective way to enhance ICL in large language models.

Another critical dimension is the role of factual knowledge. The Knowledgeable In-Context Tuning (KICT) framework \cite{ref111} tackles ICL from three angles: the model's inherent knowledge, the knowledge contained in the selected demonstrations, and the knowledge biases that affect output generation. KICT first injects factual knowledge into the model through continual self-supervised pretraining, then selects in-context examples with high knowledge relevance to the input, and finally calibrates the prediction using the model's prior knowledge. This holistic approach improves ICL performance by more than 13\% on text classification and 7\% on question answering, demonstrating that explicitly incorporating knowledge into the pretraining and fine-tuning pipeline can unlock substantial gains.

Beyond the methods discussed above, the generation of high-quality instruction-tuning data itself can be seen as a fine-tuning recipe. \cite{ref102} proposes to use smaller, permissively licensed language models to generate synthetic instruction-following examples. By categorizing and simplifying ICL templates and ensembling over multiple model outputs, this approach yields better instruction-tuning data than the original Self-Instruct method, and fine-tuning on such data improves both vanilla and instruction-tuned models. This illustrates that data-centric strategies at the fine-tuning stage can indirectly strengthen the in-context capabilities of the resulting model.

Collectively, these works demonstrate that neither ICL emergence nor its quality is purely a function of model scale. Pretraining data construction that promotes analogical reasoning, curation of stable example pools, identification and continued training on supportive pretraining instances, large-scale instruction meta-learning, and knowledge-aware tuning all offer concrete pathways to build more reliable and powerful in-context learners. As the community moves toward more data-efficient and controllable ICL, these recipes provide a foundation for designing models that learn from context as a first-class capability.

\section{In-Context Learning Across Modalities and Domains}

\subsection{Visual In-Context Learning}

Visual in-context learning (ICL) extends the idea of learning from a few demonstrations into the visual domain, enabling models to infer a desired image transformation directly from an input-output example pair without any parameter update. In pure computer vision tasks such as semantic segmentation, object detection, and image restoration, visual ICL presents a query image alongside a prompt consisting of one or more example images (the ``in-context pair'') that illustrate the required operation, and the model is expected to apply that operation to the query. This paradigm shifts the burden of adaptation from costly fine-tuning to the construction of a well-chosen prompt, and recent studies have shown that both the selection of the demonstration and the way its information is fused into the model's reasoning are decisive for performance.

Prompt selection and prompt fusion are two major factors that directly influence the inference performance of visual context learning, as systematically investigated in \cite{ref112}. The work shows that the most informative prompt should be identified using a pixel-level retrieval method that matches the query to candidate support pairs based on their visual similarity, ensuring that the demonstration is highly relevant to the task at hand. Once the prompt is chosen, the information from the in-context example must be combined with the query representation inside the large-scale vision model. Instead of relying on a single fusion mechanism, the authors propose a framework named prompt-SelF that harnesses knowledge from different positions within the model by applying multiple prompt fusion strategies in parallel. The aggregated predictions from these complementary fusion pathways are then ensembled to obtain the final output, which leads to consistent gains across single-object segmentation and detection benchmarks. This work underlines that effective prompt engineering in vision is not only about choosing the right example but also about designing how that example interacts with the internal representations of the pre-trained model.

Beyond selecting and fusing existing examples, another line of research explores augmenting the in-context pair itself with a learnable perturbation to amplify its guiding power. The method Instruct Me More (InMeMo), introduced in \cite{ref113}, injects a lightweight, trainable perturbation directly into the in-context image pair. This perturbation is optimized to encourage the model to produce the correct output for a given query, effectively adapting the prompt while keeping the foundational model frozen. On reference tasks such as foreground segmentation and single-object detection, InMeMo achieves substantial improvements over a baseline that uses only the raw in-context pair, raising the mean Intersection over Union (mIoU) by notable margins (7.35 and 15.13 points respectively). The approach exemplifies how a minimal amount of task-specific meta-training can unlock latent in-context capabilities of vision models, making visual ICL more robust to the variability in example quality.

A particularly promising direction for visual ICL is its integration with self-supervised objectives, notably masked image modeling (MIM). The framework SimICL \cite{ref114} combines the in-context learning idea with MIM pre-training to tackle medical image segmentation with very few annotations. In SimICL, the model is first trained to reconstruct masked image patches in a self-supervised manner, which builds rich visual representations. During inference, an in-context pair that consists of an input image and its corresponding segmentation mask is provided alongside the query image, and the model predicts the mask for the query without any gradient steps. This union of self-supervision and ICL proves to be highly effective in low-data regimes: on a wrist ultrasound dataset with only a limited number of manually annotated images, SimICL attained a Dice coefficient of 0.96 and a Jaccard index of 0.92, significantly outperforming state-of-the-art segmentation approaches and other visual ICL models. The success of SimICL highlights that MIM pre-training can supply the necessary inductive bias for visual ICL to generalize from a single demonstration, reducing the need for extensive labeled data.

The demand for training-free adaptation across diverse medical imaging modalities has also spurred the development of universal ICL frameworks. SegICL \cite{ref115} is designed to handle out-of-distribution (OOD) tasks without any model fine-tuning. It supports both text-guided segmentation and few-shot ICL with a small set of image-mask pairs, enabling the model to adapt to new modalities and segmentation targets on the fly. Experimental evidence shows that the segmentation performance improves steadily as the number of prompt samples increases, confirming that the model effectively extracts task context from the provided examples. Notably, SegICL achieves segmentation quality comparable to fully supervised models on both OOD and in-distribution data, demonstrating that a properly constructed in-context prompt can bridge the gap between the training distribution and unseen medical imaging scenarios. Such frameworks are especially valuable in clinical settings where acquiring large labeled datasets for every new modality is impractical.

The convergence of visual ICL and self-supervision, as well as the design of modular prompt adaptation techniques, points toward a future where a single vision model can be repurposed for a wide array of tasks using only a handful of illustrative examples. The current literature shows that the effectiveness of this paradigm hinges on multiple interconnected components: the retrieval of semantically relevant demonstrations, the multi-modal fusion of prompt and query features, the possibility of learning lightweight prompt perturbations, and the synergy with self-supervised pre-training objectives like masked image modeling. While substantial progress has been made in segmentation and detection, challenges remain in extending visual ICL to more dynamic tasks such as video understanding and in ensuring the robustness of prompts under distribution shifts. Nevertheless, the results in low-annotation medical imaging settings underscore the transformative potential of visual in-context learning for domains where expert labeling is prohibitively expensive. Building on these pure-vision foundations, the next subsection examines multimodal in-context learning, where models must jointly reason over interleaved images and text, introducing a further layer of cross-modal interaction challenges.

\subsection{Multimodal and Vision-Language In-Context Learning}

Multimodal in-context learning (ICL) extends the few-shot adaptation paradigm of large language models (LLMs) to vision-language models (VLMs) that process interleaved sequences of images and text, building directly on the pure-vision ICL foundations discussed in the previous subsection but adding the complexity of cross-modal interaction. In this setting, a VLM is given a prompt containing a handful of image--text demonstration pairs followed by a query image and must generate an answer or output text by leveraging patterns in the demonstrations without any parameter updates. Unlike pure language ICL or the visual-only ICL surveyed earlier, multimodal ICL requires the model to simultaneously interpret visual content, textual instructions, and the relationships across modalities within the context window. This subsection surveys the key findings, design strategies, lightweight adaptation modules, and evaluation benchmarks that characterize the current landscape of multimodal and vision-language in-context learning, and it bridges to the subsequent discussion of cross-lingual ICL by highlighting how alignment between heterogeneous information sources is a recurring theme across ICL paradigms.

A fundamental question centers on the relative importance of visual and linguistic information in multimodal demonstrations. An empirical investigation in \cite{ref116} revealed that ICL performance in VLMs is predominantly driven by the textual portion of the demonstrations; the visual component barely affects outcomes when the text is already informative. Through analysis of information flow and internal states, the authors demonstrated that the model attends to linguistic cues far more than to the accompanying images. Motivated by this, they proposed Mixed Modality In-Context Example Selection (MMICES), which selects demonstrations by jointly considering visual and language modalities, leading to improved ICL performance over purely text-based selection. This finding underscores a critical cross-modal imbalance that echoes the prompt selection challenges in visual ICL---where pixel-level retrieval proved essential---and foreshadows the alignment problems that also plague cross-lingual ICL when one modality or language dominates the context.

The design of effective multimodal in-context examples extends beyond simple selection heuristics. For image captioning, researchers have explored different strategies for image selection and caption assignment to configure in-context image--text pairs \cite{ref117}. Their experiments yielded counter-intuitive insights: in vision-language ICL, the synergy between modalities creates distinct dynamics compared to language-only settings, and optimal combinations can improve CIDEr scores by over 20 points on average. A more radical approach, In-Image Learning (I\(^2\)L) \cite{ref118}, condenses all demonstrations, visual cues, and chain-of-thought reasoning into a single aggregated image. By leveraging the image understanding capabilities of large multimodal models like GPT-4V, I\(^2\)L avoids the inaccuracies of textual descriptions of complex images and provides flexibility in positioning examples, while a hybrid variant automatically selects between I\(^2\)L and other ICL methods per instance. This trend toward fusing information into a unified representation parallels the prompt fusion strategies in visual ICL and the retrieval-augmented cross-lingual methods that source high-resource exemplars to bridge representational gaps.

To address the limited few-shot ICL capabilities of off-the-shelf VLMs, several lightweight tuning modules have been introduced. MultiModal In-conteXt Tuning (M\(^2\)IXT) \cite{ref107} is a compact, architecture-agnostic module that can be prepended to various multimodal unified models (OFA, Unival, LLaVA) and trained on a mixture of tasks. It perceives an expandable context window containing labeled examples of multiple modalities and, with as little as 50K training data, yields significant relative ICL gains (e.g., 18\% for OFA) while being about 20× smaller than Flamingo or MMICL. Multi-Modal In-Context Tuning (MMICT) \cite{ref108} introduces a Multi-Modal Hub (M-Hub) that unifies diverse multimodal features; the model learns from in-context visual-guided textual features and then generates outputs conditioned on textual-guided visual features. MMICT consistently outperforms standard fine-tuning and naïve concatenation baselines on a range of multimodal tasks, much as InMeMo's learnable perturbations amplify the guiding power of visual in-context pairs. Meanwhile, MMICL \cite{ref119} tackles the challenge of complex multi-modal prompts containing multiple images. By devising a novel context scheme and constructing the Multi-modal In-Context Learning (MIC) dataset, MMICL empowers VLMs to interpret interleaved multi-image text inputs, achieving state-of-the-art zero-shot performance on benchmarks like MME and MMBench while also alleviating language bias that can cause hallucinations under heavy textual context.

A further line of work directly trains VLMs to follow ICL instructions. \cite{ref120} found that even large-scale mixed-modality pretrained VLMs under-perform when prompted with few-shot demonstrations, attributing this to a lack of direct ICL instruction tuning. By extending the alignment framework with ICL support, a carefully designed curriculum, and effective data mixes, the approach achieved an average ICL performance boost of 11.3\% (up to 21.03\%) over strong baselines while also contributing new multimodal ICL evaluation benchmarks. This mirrors the lesson from cross-lingual ICL that instruction tuning in the target modality or language is critical for unlocking robust in-context capabilities, and it reinforces the broader theme that pre-training alone is insufficient without deliberate ICL-oriented alignment.

Cross-modal interaction bottlenecks remain a central challenge. When interleaved modalities are simply concatenated, the understanding process can convolute, as noted in the motivation for M\(^2\)IXT. Models may fail to extrapolate effectively from contextual examples due to limited cross-modal fusion. To foster deeper interaction, UNIMO-3 \cite{ref121} introduces multi-granularity interaction that simultaneously learns in-layer and cross-layer cross-modal connections, preventing interaction from collapsing into a narrow semantic space. Deeply Coupled Cross-Modal Prompt Learning (DCP) \cite{ref122} employs a Cross-Modal Prompt Attention mechanism that progressively enables mutual exchange of representations through multi-head attention. Although DCP is a prompt-tuning method rather than pure ICL, it illustrates a strategy for strengthening vision-language alignment that can benefit in-context adaptation, just as SimICL leverages masked image modeling to supply inductive bias for visual ICL. Moreover, a vision-language consistency analysis \cite{ref123} reveals that model performance across modalities can diverge on complex tasks, indicating that vision and language processing are not always synchronized; such inconsistency directly undermines the reliability of multimodal ICL and motivates the development of vision description prompting to improve difficult vision tasks. These cross-modal alignment challenges are conceptually analogous to the script and vocabulary disparities that hinder cross-lingual ICL, where deliberate alignment of semantic spaces is equally necessary.

Robust evaluation of multimodal ICL requires benchmarks that go beyond standard few-shot VQA or image captioning. The VL-ICL Bench \cite{ref124} provides a comprehensive suite spanning perception, reasoning, and long-context tasks with both image and text inputs and outputs. It exposes diverse strengths and weaknesses of state-of-the-art VLLMs, including GPT-4, and reveals that even the most advanced models find many ICL tasks challenging. In the specific setting of text-to-image ICL, the CoBSAT benchmark \cite{ref125} defines ten tasks to test whether multimodal LLMs can generate images from in-context textual prompts paired with example images. The study uncovers significant difficulties, with fine-tuning and chain-of-thought prompting offering partial mitigation, highlighting the infancy of T2I-ICL research. These evaluation efforts parallel the need for specialized benchmarks in cross-lingual ICL that probe alignment quality across typologically diverse language pairs.

Finally, moving toward a more unified treatment, one framework frames visual ICL as generative modeling over interleaved sequences of quantized text and vision tokens within a decoder-only sparse transformer, enabling multimodal outputs in a single pipeline \cite{ref126}. This unified representational space points toward a future where modality-specific gaps are minimized and ICL can seamlessly chain visual and textual reasoning, much as cross-lingual ICL aspires to a shared semantic space that renders language boundaries transparent. Overall, multimodal ICL has advanced rapidly through modality-aware example selection, lightweight parameter-efficient tuning modules, and specialized benchmarks, yet open problems persist in balancing cross-modal contributions, ensuring consistency, and scaling to complex reasoning tasks with interleaved context. These challenges resonate across the entire ICL landscape: the need to carefully select and align heterogeneous information sources, to design fusion mechanisms that prevent one channel from dominating, and to develop evaluation protocols that capture the full complexity of the adaptation setting. The next subsection turns to cross-lingual ICL, where analogous alignment gaps between high-resource and low-resource languages demand a parallel set of solutions---query alignment, structured prompt construction, and retrieval augmentation---that mirror the multimodal strategies discussed here.

\subsection{In-Context Learning for Low-Resource and Multilingual Settings}

In-context learning (ICL) extends its adaptation capabilities beyond multimodal integration into the equally critical context of low-resource and multilingual scenarios, where large language models (LLMs) can adapt to new tasks and languages through a few demonstrations without any parameter updates. This capability is particularly beneficial for languages that lack large annotated corpora, as it bypasses the expensive cycle of data collection and model fine-tuning. However, naive application of ICL in cross-lingual settings often falls short because random selection of input-label pairs for the prompt context leads to a misalignment between the high-resource source language and the low-resource target language in both input semantics and output label spaces \cite{ref127}. This subsection surveys strategies that explicitly address these alignment gaps, covering query-level over label-level alignment, specialized prompt construction, retrieval-augmented cross-lingual techniques, and the training of multilingual in-context learners, while emphasizing the pivotal role of instruction tuning in the target language and the treatment of script and vocabulary disparities.

The cornerstone of effective cross-lingual ICL lies in aligning the semantic content of the test instance with the demonstrations. A recent in-depth analysis \cite{ref128} revealed that the conventional approach of aligning labels via parallel examples is severely limited for low-resource languages. Instead, aligning on the query---i.e., selecting demonstrations that are semantically similar to the test input in the target language---proves substantially more effective. This query alignment strategy reduces the language gap by focusing on the shared meaning rather than exact lexical overlap, enabling LLMs to leverage their multilingual pretraining to better generalize. The study demonstrated that even a handful of query-aligned examples can dramatically improve low-resource understanding quality, underscoring the need to prioritize semantic relevance when constructing prompts across languages.

To further inject coherence and task-based alignment, the X-InSTA (Cross-lingual In-context Source-Target Alignment) method was proposed \cite{ref127}. X-InSTA builds prompts by ensuring that the selected input examples are not only semantically coherent but also aligned with the task across source and target languages. Through a structured alignment of the input examples and a deliberate pairing of labels, X-InSTA outperforms random prompt selection by large margins on text classification tasks over 44 cross-lingual pairs. This demonstrates that careful construction of the demonstration set can bridge the representation gap caused by language typological distance and disparate label mappings.

Retrieval-augmented methods offer a complementary pathway by sourcing informative demonstrations from a high-resource language corpus. The PARC (Prompts Augmented by Retrieval Crosslingually) pipeline \cite{ref129} retrieves semantically similar sentences from a high-resource language to augment the context for a low-resource test instance, achieving significant zero-shot improvements on sentiment classification, topic categorization, and natural language inference across 10 low-resource languages. The performance gains correlate positively with the similarity between the high- and low-resource languages and the amount of low-resource pretraining data, highlighting the importance of both language relatedness and model capacity. Similarly, CREA-ICL \cite{ref130} extracts semantically similar prompts from high-resource languages to improve zero-shot performance of multilingual pre-trained models. While it yields steady gains in classification tasks, generation tasks pose additional challenges, indicating that cross-lingual retrieval may need task-specific adaptation for open-ended outputs. These retrieval-augmented approaches collectively show how leveraging high-resource exemplars can effectively serve as a cross-lingual bridge.

Beyond inference-time prompt design, training strategies that explicitly target multilingual ICL capabilities have emerged. Efforts such as collecting and formatting task resources for Slavic languages \cite{ref131} revealed that ICL models fine-tuned on English instruction data can still learn from non-English contexts, but multilingual instruction fine-tuning consistently improves ICL ability across languages. Interestingly, this work found that single-task training in the target language can sometimes outperform massive multitask training, suggesting that specialized in-context learners attuned to the specific language and task requirements can be more effective. The broader lesson is that instruction tuning in the target language---or in a multilingual mix that includes the target language---is critical for unlocking robust cross-lingual ICL.

The impact of instruction tuning is further highlighted by cross-lingual instruction tuning methods such as CoIT \cite{ref132}, which tunes LLMs on translation and cross-lingual general task data to align semantic spaces across languages. While not exclusively designed for ICL, such alignment directly benefits ICL by reducing the representational distance between languages, enabling the model to better interpret in-context demonstrations in the target language. Similarly, CrossIn \cite{ref133} uses a mixed composition of cross-lingual instruction data to efficiently enhance task-solving and multilingual proficiency, providing a foundation on which improved ICL can be built.

Handling script and vocabulary disparities is a recurring obstacle. Languages with different writing systems (e.g., Cyrillic, Devanagari, Arabic scripts) or limited lexical overlap with the high-resource language can severely impair ICL when demonstrations are presented in the raw script. The alignment methods discussed---query alignment, cross-lingual prompt construction, and retrieval augmentation---circumvent surface-form mismatches by operating on semantic representations and, in the case of retrieval augmentation, by injecting high-resource language prompts that the model understands well. Moreover, multilingual instruction tuning that exposes the model to diverse scripts and transliterations further unifies the representation space, reducing the impact of orthographic variance. The integration of additional modalities, such as phonemic transcriptions \cite{ref134}, has shown promise in bridging script gaps in non-ICL transfer settings and could inspire future ICL-specific approaches.

In summary, just as multimodal ICL must carefully balance and align visual and linguistic signals, cross-lingual ICL demands deliberate alignment of input semantics, label spaces, and task formats across languages. The strategies of query alignment, aligned prompt construction (X-InSTA), retrieval augmentation (PARC, CREA-ICL), and careful multilingual instruction tuning collectively mitigate inherent input-output space misalignment and script disparities. By focusing on semantic relevance and task-level alignment, these methods significantly boost cross-lingual transfer without gradient updates, extending the reach of in-context learning to low-resource languages and complementing its applications in multimodal and domain-shift settings.

\subsection{Domain-Specific and Cross-Domain Adaptation with ICL}

\textbf{Mitigating distribution shifts with in-context demonstrations.}\\
In the same way that cross-lingual ICL aligns input--output spaces across languages without parameter updates, ICL can counteract distribution shifts by embedding target-domain examples directly in the prompt, enabling the model to infer the new data distribution at inference time. The effectiveness of this strategy, however, depends crucially on the representativeness of the few-shot examples. In cross-modality medical image segmentation, for instance, models trained on CT scans often produce near-zero Dice scores when applied to MRI images without any adaptation \cite{ref135}. ICL can alleviate this collapse by supplying a small number of annotated MRI slices, effectively guiding the model toward the target appearance and label semantics---much as query-aligned examples guide a multilingual LLM toward the target language meaning. The label-efficient UDA paradigm demonstrated in \cite{ref136} highlights that even a few source annotations, combined with self-ensembling consistency, can significantly boost cross-modality segmentation. ICL extends this principle to a training-free regime: carefully chosen target examples serve as anchors that recalibrate the model's internal representations, mirroring the recalibration strategies used in adversarial UDA where networks adapt by rejecting improbable patterns on a few target images \cite{ref137}.

\textbf{Vulnerability to domain and label shifts.}\\
Despite its promise, ICL is sensitive to the quality and composition of the prompt, echoing the observation in cross-lingual settings that random selection of input--label pairs leads to misalignment. If the few-shot demonstrations fail to reflect the full label distribution or contain significant noise, the model's predictions can remain biased toward the source domain. The concept of category-level misalignment, which motivates class-level alignment methods such as \cite{ref138}, also applies to ICL: when the target domain has a different class distribution or novel sub-types, standard ICL with random demonstration selection may not sufficiently correct for the shift. Moreover, the semantic prior learned during pretraining can dominate, causing the model to over-rely on learned prototypes instead of adapting to the novel domain. Explicit domain-agnostic prompting, as explored in \cite{ref139}, shows that aligning textual and visual prompts in a domain-invariant fashion improves cross-domain transfer. ICL frameworks can adopt a similar philosophy by structuring prompts to include domain-agnostic descriptions or by selecting examples that minimize style-specific cues, thereby reducing the pull of spurious correlations---just as cross-lingual ICL benefits from semantically aligned rather than surface-form aligned demonstrations.

\textbf{Cross-modal and cross-domain knowledge transfer for 3D semantic segmentation.}\\
Extending the multimodal integration challenges discussed in prior subsections, many specialized tasks such as 3D semantic segmentation from LiDAR point clouds and camera images require fusing information from multiple modalities. Conventional domain adaptation for this setting leans heavily on cross-modal consistency, where predictions from 2D and 3D branches are forced to agree on unlabeled target data \cite{ref140}. ICL can similarly exploit multi-modal demonstrations: a prompt could contain image-point-cloud pairs with segmentation masks, enabling the model to internalize the relationship between modalities and the task. The dynamic sparse-to-dense cross-modal learning described in \cite{ref141} underscores the need to align dense 2D features with sparse 3D features. When designing ICL prompts for such tasks, one might dynamically fuse modalities by concatenating visual tokens or by employing cross-attention modules within the architecture, mimicking the inter-domain cross-modal adversarial alignment that promotes high-level modality complementarity.

Medical image analysis further exemplifies cross-modal adaptation challenges. Brain tumor segmentation across MRI modalities, for example, benefits from edge-map guidance to bridge appearance gaps \cite{ref142}. ICL prompts in this context can incorporate contour information as an additional channel or textual hint, giving the model a low-level signal that is more invariant across modalities---analogous to how cross-lingual ICL uses semantic representations to circumvent script disparities. Similarly, the synergistic image-and-feature adaptation paradigm of SIFA \cite{ref143} transforms both the image style and the feature representation to align domains; in ICL, this can be achieved by selecting demonstration images that simulate the target style (e.g., through a style transfer preprocessing step) and by appending a textual prompt that instructs the model to be style-invariant.

\textbf{Few-shot specialization through meta-learning and hallucination.}\\
ICL can be viewed as a form of meta-learning where the model extracts task information from the prompt, a perspective that aligns naturally with the emerging ICL paradigms in reinforcement learning described in the following subsection. This connection is especially relevant when annotated data in the source or target domain is extremely scarce. The meta-hallucinator approach \cite{ref144} learns to generate diverse synthetic examples that help a segmentation model generalize across modalities. ICL can adopt a similar spirit by prompting the model with multiple hallucinated examples --- for instance, by using a generative model to produce plausible target-domain images given a few seed examples --- and constructing a prompt that includes both real and hallucinated demonstrations. This strategy can increase the effective number of shots and improve robustness to domain shift without any parameter updates, much as retrieval-augmented ICL enriches cross-lingual prompts with high-resource exemplars.

\textbf{Domain-agnostic prompts and calibration.}\\
A key direction for robust ICL-based domain adaptation is the design of prompts that are inherently domain-agnostic. The DAMP framework \cite{ref139} shows that learning prompts that simultaneously condition visual and textual branches in a domain-invariant manner greatly outperforms separate per-domain prompting. For ICL, this suggests that the language instruction itself should avoid domain-specific terminology and instead describe the task in general terms. Moreover, when visual demonstrations are present, they should be selected to be prototypical of the task rather than the domain. Calibration techniques such as batch calibration or in-context calibration can further mitigate the biases that remain when using few-shot target examples, ensuring that the model's confidence aligns with its accuracy even under severe label shifts---a concern that resonates with the evaluation and calibration challenges identified in reinforcement learning and speech ICL.

\textbf{Challenges and outlook.}\\
While ICL holds strong potential for domain-specific adaptation, several open questions remain that connect directly to the cross-modal alignment themes of the following subsection. The selection of faithful and representative demonstrations in a training-free manner is nontrivial, especially when the target domain is unlabeled. Techniques like cross-domain consistency losses \cite{ref145} and category-centric prototype alignment \cite{ref146} can inspire new in-context example scoring functions that favor examples predicted consistently across multiple views. The vulnerability to adversarial demonstration perturbations and the need for long-context handling when multiple modalities are involved also demand attention---challenges that are amplified in the dynamic, sequential settings of RL ICL. As foundation models continue to scale, cross-modal ICL that seamlessly integrates text, images, point clouds, and structured data will become increasingly crucial, making domain-agnostic, calibration-aware prompting a central pillar for deploying trustworthy in-context learners in specialized and safety-critical domains, whether adapting to a new language, a new imaging modality, or an entirely new sensory context.

\subsection{In-Context Learning in Reinforcement Learning, Speech, and Structured Data}

In-context learning (ICL) has demonstrated remarkable flexibility in adapting to new tasks using only a few demonstrations, a capability that extends well beyond the text and static vision domains discussed in prior sections. Building on the domain adaptation strategies and multimodal alignment challenges previously outlined, this subsection surveys emerging ICL paradigms in reinforcement learning (RL), speech processing, and reasoning over structured data. In these settings, contextual examples serve not only to illustrate desired outputs but also to specify dynamic policies, interpret acoustic commands, or leverage hierarchical knowledge structures. The unique representation, alignment, and evaluation challenges that arise in these modalities further motivate the cross-modal alignment techniques that the following section will examine in depth.

\textbf{In-Context Learning in Reinforcement Learning.}\\
RL agents typically learn policies through trial and error, but ICL offers a way to rapidly adapt behavior by conditioning on a few trajectory demonstrations or task descriptions provided in the prompt. Early work on instruction following demonstrated that a single neural network could jointly reason over linguistic and visual observations to directly output actions without intermediate representations, effectively using the instruction as a contextual cue \cite{ref147}. This approach, trained via contextual bandit RL, showed that raw visual and textual input could be mapped to correct actions, hinting that an agent can learn to interpret a new task ``in context'' from a single instruction. More explicitly, the use of step-by-step human demonstrations in the form of natural language instructions and action trajectories allows agents to generalize to unseen tasks in a zero-shot manner \cite{ref148}. By conditioning a high-level language generator and a low-level policy on the provided language context, the model learns to reuse low-level skills and compose them for novel goals, a clear instance of in-context policy specification.

Multimodal cues further enrich the context. In task-oriented visual dialog, a hierarchical reinforcement learning policy jointly learns a multimodal dialog state representation and a hierarchical dialog policy, integrating vision and language to improve task success \cite{ref149}. Here, the dialogue history and visual observations serve as an evolving context from which the agent infers the next action, analogous to ICL where the prompt history encodes the task. Interactive RL settings also benefit from multimodal in-context signals: audio-visual feedback (vocal commands and hand gestures) can provide robust instructions, and the integration of multiple sensory modalities with a confidence value mirrors the way a language model absorbs varied demonstrations to resolve ambiguous tasks \cite{ref150}. This work shows that the agent can leverage both visual and auditory context to accelerate learning, especially when the feedback is modulated by affordance knowledge.

The alignment of state and action representations with language semantics opens a direct path to language-based in-context task specification. MORE-3S transforms offline RL into a supervised learning task by embedding states and actions into a shared semantic space with language, using pretrained language models to encode goals and transitions \cite{ref151}. This alignment allows the model to interpret a language prompt describing the desired behavior and infer the appropriate policy without additional training, a core ICL capability. Similarly, hierarchical context can be exploited to decompose complex tasks: Context-Hierarchy Inverse Reinforcement Learning models the context as a directed acyclic graph, where each node corresponds to a subtask or a context factor, and the reward function is a modular network that leverages this hierarchical structure \cite{ref152}. By providing demonstrations along with the context hierarchy, the agent learns to infer the reward function and policy structure, effectively using the structured context as a prompt that organizes the learning process.

Key challenges in RL ICL include the dynamic and sequential nature of the task, the need to align temporal demonstrations with an agent's internal state representation, and the difficulty of generalizing across varying reward functions and environmental dynamics. The length of the context (number of demonstrations) is often limited by the model's capacity or the computational budget, and the ordering of demonstrations can significantly affect policy learning, just as in text-based ICL. Moreover, evaluation requires careful protocol design to distinguish between true in-context adaptation and memorization of training environments.

\textbf{In-Context Learning in Speech and Audio.}\\
While speech ICL is less mature, the principles of using a few acoustic examples to specify a task are beginning to emerge. In interactive RL, audio commands can be part of the multimodal prompt, enabling the agent to interpret spoken instructions without explicit speech recognition \cite{ref150}. The use of auditory events as intrinsic motivation also points to the potential of audio-based context: an agent can learn to predict auditory events caused by its actions and use prediction errors as intrinsic rewards, effectively learning from audio cues in the environment \cite{ref153}. Although this is not a direct instance of ICL, it demonstrates that audio signals can carry task-relevant information that can be exploited in a few-shot manner. Extending ICL to pure speech tasks---such as speech recognition, emotion detection, or speaker identification via acoustic prompts---remains an open frontier. The main obstacles include the variable-length and continuous nature of audio, the need for robust representations that align with semantic content, and the scarcity of large-scale pretrained models that can ingest raw audio as in-context demonstrations. Aligning speech with text or visual modalities in a shared context is another challenge, as the semantic gap between phonetic and linguistic levels is substantial.

\textbf{In-Context Learning for Structured Data.}\\
Structured domains such as graph reasoning and code generation can also benefit from in-context demonstrations. Graph neural networks, for instance, can be presented with a few input-output graph pairs that exemplify a transformation, and the model is expected to apply the same transformation to a new graph. The hierarchical context approach in inverse RL, where the task structure is encoded as a directed acyclic graph, provides a template for how structured priors can be injected into the prompt to guide learning \cite{ref152}. By decomposing a complex task into a graph of subtasks, the model can use the given context hierarchy to infer the appropriate reward modules and policy, analogous to providing a structured prompt that outlines the reasoning steps. In code generation, while not represented in the surveyed papers, a similar principle applies: a few examples of input-code-output triples can teach a model to generate new code for a novel problem. The challenges here include the representation of discrete structures as token sequences, the need for compositional generalization, and the design of prompts that capture the underlying algorithmic pattern without overfitting to superficial features.

\textbf{Shared Challenges and Future Directions.}\\
Across all these modalities, the core tension is between the need for rich, expressive context and the difficulty of aligning that context with the model's internal representations. In RL, the context consists of sequences of states, actions, and rewards, which must be converted into a format that the model can attend to effectively. In speech, raw audio must be serialized and aligned with task-specific labels. In structured data, the graphs or code must be linearized without losing structural information. These representation hurdles directly echo the cross-modal alignment challenge that the next subsection explores in depth: the need to map heterogeneous modalities into a shared space while preserving task-relevant signal. Evaluation also remains a challenge: standard RL benchmarks and speech tasks are not designed to measure in-context adaptation, requiring new protocols that isolate the model's ability to glean task information solely from demonstrations. Nevertheless, the convergence of ICL with these modalities promises agents that can quickly adapt to new tasks from minimal guidance, using language, images, audio, and structured knowledge interchangeably. Future work should focus on developing unified architectures that can handle heterogeneous context, robust pretraining strategies that expose models to diverse multimodal demonstrations, and calibration methods that ensure reliable behavior when the context is ambiguous or noisy---themes that will be further elaborated in the discussion of cross-modal alignment techniques.

\subsection{Cross-Modal Alignment and Unified Representation Learning for ICL}

Cross-modal alignment stands as a foundational prerequisite for achieving robust in-context learning (ICL) across modalities. While the preceding subsection surveyed how ICL extends into reinforcement learning, speech, and structured data---each presenting unique representation and alignment hurdles---multimodal ICL must bridge heterogeneity between visual, textual, and other modal tokens. This challenge directly impacts the fidelity of demonstration encoding and the model's ability to extract task-relevant patterns from interleaved context, much as RL agents must align state-action trajectories with language semantics and structured data must be linearized without loss of information. The central tension arises from the representation disparities between modalities, where pretrained unimodal encoders embed images and text into spaces that do not naturally align, and from the semantic gaps that emerge when visual scenes and linguistic descriptions capture complementary but non-isomorphic information. This subsection reviews the principal approaches to cross-modal alignment and unified representation learning that underpin multimodal ICL, analyzing how they address modality gaps, mitigate representation disparities, and shape the emergence of cross-modal in-context capabilities.

A dominant paradigm is the construction of unified sequence models that embed visual and textual tokens into a shared representational space and process them with a single decoder or encoder-decoder architecture. The framework introduced in \cite{ref126} illustrates this direction: visual and text prompts are quantized and embedded into a joint space structured as interleaved in-context sequences, and a decoder-only sparse transformer performs generative modeling over them, enabling multimodal output. This design directly aligns with the ICL mechanism by treating demonstration examples from any modality as members of a common token vocabulary, allowing attention to flow across modalities as if they were words---akin to how RL approaches like MORE-3S embed states, actions, and language goals into a shared semantic space. The challenge, however, lies in the quality of the shared embedding space. Simple projection layers may fail to capture fine-grained correspondences, necessitating more sophisticated alignment losses during pretraining or post-hoc tuning.

Contrastive learning has emerged as a powerful tool for reducing the modality gap and creating robust joint embedding spaces conducive to ICL. Methods like \cite{ref154} combine contrastive predictive coding and noise contrastive estimation to preserve both intra- and inter-class relationships, generating a task-agnostic joint embedding that significantly narrows the modality gap. Similarly, \cite{ref155} argues that standard cross-modal alignment alone may ignore intra-modal structure, leading to degraded representations. By adding intra-modal contrastive objectives alongside the cross-modal loss, the model learns to keep similar inputs close within each modality while still aligning paired image-text data, resulting in richer representational manifolds that can better support the retrieval and integration of multimodal demonstrations. In the context of ICL, such refined embeddings enable more accurate cross-modal analogy mapping, as the distances in the shared space more faithfully reflect both semantic similarity and inter-modal correspondence---directly addressing the alignment difficulties identified in the RL and speech settings of the previous subsection.

The benefits of hybrid contrastive strategies extend to tri-modal and fine-grained scenarios. The \cite{ref156} framework simultaneously performs intra- and inter-modal contrastive learning and introduces a semi-contrastive component to preserve inter-class relationships, explicitly addressing the modality gap by capturing cross-modal interactions at multiple levels. For tasks with limited semantic information, such as understanding dark scenes, \cite{ref157} shows that supervised cross-modal and intra-modal contrast encourages class-discriminative feature spaces where embeddings from different modalities cluster by category. When applied to ICL, these contrastive pretraining paradigms reduce the risk that a multimodal learner will misinterpret a visual demonstration because its projection lies far from semantically analogous textual examples, thus improving the robustness of in-context task inference---a concern that resonates with the hierarchical context decomposition challenges observed in structured data reasoning.

Beyond contrastive pretraining, lightweight modular tuning allows the injection of cross-modal alignment capabilities specifically for ICL without retraining the entire backbone. \cite{ref107} (M²IXT) introduces an expandable context window module that can be prepended to unified multimodal models such as OFA, Unival, or LLaVA, enabling rapid few-shot adaptation across multiple tasks. The module is trained with a mixed-task strategy on as few as 50K multimodal examples yet yields substantial ICL performance boosts by learning to better integrate interleaved multimodal examples. Similarly, \cite{ref108} proposes a Multi-Modal Hub (M-Hub) that captures various multimodal features and enables multi-modal LLMs to learn from in-context visual-guided textual features, effectively bridging the representation gap through a specialized in-context tuning paradigm. These approaches highlight that a small number of parameters, when specifically dedicated to cross-modal interaction, can compensate for the lack of direct multimodal ICL pretraining---offering a pragmatic path toward the unified architectures that can handle heterogeneous context, as called for in the previous subsection's future directions. Hierarchical prompt tuning, as exemplified by cross-modal bridges that learn layer-wise attention routing like \cite{ref158}, further refines the alignment by dynamically selecting attention modes that best fuse visual and textual information at different depths, mitigating the modality mismatch that otherwise impairs ICL.

Despite these advances, significant representation disparities persist, often causing ICL to rely disproportionately on a single modality. The empirical study in \cite{ref116} reveals that ICL in VLMs is predominantly driven by textual information, with visual demonstrations barely influencing performance. This indicates that the cross-modal alignment learned during pretraining may be insufficient to project visual examples into a space where they can effectively modulate the output distribution when presented in-context. The authors propose Mixed Modality In-Context Example Selection (MMICES) to explicitly consider both modalities during demonstration retrieval. Moreover, \cite{ref159} addresses representation disparities by introducing intent-oriented image summarization and demonstration composition, converting visual context into linguistically summarized tokens that alleviate cross-modal interaction bottlenecks. These findings underscore that alignment must not only map modalities into a common space but also ensure that the resulting representations carry equivalent task-relevant signal for the ICL mechanism---a requirement that echoes the difficulty in RL of ensuring that trajectory demonstrations and language instructions contribute balanced information to policy inference.

Semantic gaps further compound the difficulty. The pull of modality-specific priors, particularly in subjective domains, can overshadow the information provided by demonstrations. Research on large language models has shown a strong pull of prior knowledge that can ossify predictions and prevent the full integration of contradictory in-context examples \cite{ref42}. In multimodal ICL, such priors may originate from the dominant modality---often text---biasing the model to disregard visual cues even when they are critical. This effect highlights the need for alignment techniques that not only unify representations but also calibrate the influence of each modality during in-context inference, directly addressing the cross-modal alignment challenges that the following subsection will explore through calibration and evaluation frameworks. Iterative contrastive alignment frameworks, such as \cite{ref160}, use recurrent alignment layers and cross-modal contrastive losses at each encoding step to progressively refine the embedding space, directly optimizing for the coherences needed in conditional generation tasks that resemble in-context reasoning.

Finally, text-centric alignment strategies leverage the unique properties of language as a universal semantic medium. \cite{ref161} exploits LLMs with in-context learning and foundation models to handle unseen modality combinations by transforming all modalities into a unified textual semantic space. This approach circumvents direct cross-modal embedding learning by relying on the strong semantic understanding of language as a bridge, an idea that resonates with the finding that textual demonstrations dominate ICL and with the text-based task specification strategies in RL such as MORE-3S. By converting visual, audio, or structured data into language-based descriptions, one can create a homogeneous demonstration stream that aligns naturally with LLM processing, thus simplifying the cross-modal alignment challenge. However, this translation introduces information loss and requires careful handling of modality-specific nuances---limitations that mirror the serialization challenges faced when linearizing graphs and structured data for in-context demonstrations.

In summary, cross-modal alignment and unified representation learning for ICL are being pursued along multiple fronts: joint embedding via contrastive losses, modular tuning that injects alignment awareness, dynamic attention routing, and text-centric bridging. Each approach grapples with the tension between the need for a shared representational space and the preservation of modality-specific detail---a tension that directly parallels the representation, alignment, and evaluation challenges identified across RL, speech, and structured data in the previous subsection. As multimodal ICL matures, a critical direction will be the development of alignment objectives that explicitly optimize for the ability to transfer task structure from demonstrations across modalities, rather than merely aligning static distributions, thereby enabling robust and balanced cross-modal in-context reasoning. This imperative naturally leads into the subsequent discussion of calibration techniques and evaluation protocols that ensure reliable behavior when context is ambiguous or noisy.

\section{Applications, Robustness, and Future Directions}

\subsection{ICL for Diverse NLP and Structured Tasks}

In-context learning (ICL) has emerged as a versatile mechanism for adapting large language models (LLMs) to a wide array of natural language processing and structured reasoning tasks without parameter updates. This subsection surveys prominent ICL deployments in machine translation, text-to-SQL semantic parsing, and text simplification, highlighting task-specific strategies such as syntax-aware example selection, metric-driven demonstration curation, decomposition-based prompting, and automatic rationale generation. Collectively, these adaptations unlock substantial domain-specific gains and provide a blueprint for tailoring ICL to the unique demands of each task.

\textbf{Text-to-SQL: From Syntax-Driven Selection to Decompositional Workflows}\\
The text-to-SQL task, which translates natural language questions into executable SQL queries, has received intense focus in ICL research due to its structured output space and reliance on database schemas. Early investigations into prompt design revealed that careful demonstration selection is critical: leveraging the syntactic structure of the target SQL queries, such as using their parse trees or keyword composition, to retrieve demonstrations that are both diverse and similar to the test instance significantly outperforms random or purely semantic baselines. The work in \cite{ref162} demonstrated that a combination of diversity and similarity in demonstration selection, paired with database-related knowledge augmentations, improves execution accuracy on the Spider benchmark by 2.5 points over the previous state-of-the-art. This underscores the importance of aligning demonstrations with the underlying structural properties of the task. Similarly, \cite{ref163} introduced ODIS, a framework that synthesizes in-domain demonstrations from out-of-domain examples, enabling cross-domain generalization without relying on costly annotated target data. By blending out-of-domain and synthetic in-domain demonstrations, ODIS achieved gains of 1.1 and 11.8 points in execution accuracy across two datasets, showing that hybrid demonstration sources can effectively compensate for domain shifts.

Moving beyond demonstration selection, several works have enhanced ICL for text-to-SQL through advanced prompting paradigms that guide the model's reasoning. Chain-of-thought (CoT) prompting, which elicits intermediate reasoning steps, has been adapted to this structured setting. \cite{ref164} eliminates manual labeling by automatically generating CoT exemplars using schema-linking heuristics, thereby maintaining a cost-efficient single-API-call regime while achieving state-of-the-art performance on the Spider dev set. The approach extends naturally to multi-turn interactions, highlighting the flexibility of auto-CoT generation. \cite{ref165} takes a divide-and-conquer approach, decomposing the text-to-SQL process into subtasks (such as schema linking, condition extraction, and query assembly) and applying CoT to each subtask. This decomposition not only reduces reasoning complexity but also yields higher execution accuracy by preventing attention diffusion. Even more radical decomposition is pursued in \cite{ref166}, where a full workflow paradigm---including an information-determination module to prune redundant context, problem-classification to structure the prompt, self-correction, and active learning---expands the problem-solving scope of LLMs. This workflow approach achieved new state-of-the-art results on the Spider test set, with improvements of 2--3 percentage points over strong baselines, illustrating how modular decomposition can amplify ICL's reasoning capacity.

A persistent challenge in text-to-SQL ICL is the design of the overall prompt structure, especially when combining database schema, few-shot examples, and instructions. The comprehensive study in \cite{ref167} systematically examined the impact of varied prompt constructions and emphasized the lack of comparability across prior works. Their findings provide practical guidance: optimal prompt formatting is setting-dependent, and careful ordering of demonstration elements can influence the model's ability to correctly ground table and column names, thereby informing future standardized evaluation protocols.

\textbf{Machine Translation: Syntax and Coherency as Selection Principles}\\
Machine translation (MT) via ICL benefits from the same sensitivity to demonstration quality, but the nature of the task shifts the focus toward linguistic structure and discourse coherence. \cite{ref168} argues that superficial word-level similarity is insufficient; instead, examples should be selected based on syntactic similarity measured via polynomial distance between dependency trees. An ensemble that combines word-level and syntax-level criteria yields the highest COMET scores across 12 translation directions, confirming that deep syntactic alignment helps the model learn transferable translation patterns from prompts. This syntax-based selection is a natural fit for MT, where reordering and structural divergence are classic challenges.

Complementing syntactic similarity, \cite{ref169} proposes a coherency-based perspective, viewing ICL not merely as learning from examples but as the model's desire to maintain coherence with its prompt context. Their experiments on English\(\rightarrow\)\{Portuguese, German, French\} translation with GPT-Neo2.7B, Bloom3B, and XGLM2.9B show that translation quality improves when prompts are drawn from the same domain, and that long-term coherency between prompts and the test sentence---captured via a moving window of in-domain examples---correlates strongly with downstream performance. This suggests that on-the-fly ICL-MT adaptation can be effectively guided by sequential coherence, enabling practical deployment without prior knowledge of domain distributions.

\textbf{Text Simplification: Metric-Driven Demonstration Curation}\\
In text simplification, the output style and quality are multifaceted, making it difficult to hand-pick examples that reliably teach the desired transformation. \cite{ref170} introduced Metric-Based In-context Learning (MBL), which uses established simplification metrics---SARI, compression ratio, and BERT-Precision---to select demonstrations. The study showed that for larger models like GPT-175B, selecting examples with the highest SARI scores yields the best simplification quality, while for smaller models (GPT-13B, GPT-6.7B), compression ratio is a more effective selector. Notably, MBL is robust to example ordering and generalizes well to out-of-domain test sets, surpassing fine-tuned state-of-the-art models on TurkCorpus and ASSET. This metric-driven approach not only automates demonstration selection but also provides a control knob: by prioritizing different metrics, practitioners can implicitly steer the model's behavior toward more conservative versus more aggressive simplifications.

\textbf{Scaling and Structured Reasoning: Expanding the ICL Horizon}\\
Beyond these task-specific adaptations, a general challenge for structured tasks is the need to incorporate many examples, as complex reasoning often benefits from richer context. Conventional ICL is bounded by LLM context windows, limiting supervision to a handful of shots. \cite{ref171} breaks this barrier by encoding each demonstration separately with dedicated position embeddings and then jointly attending to them via a rescaled attention mechanism, achieving linear complexity in the number of exemplars. When applied to tasks including natural language inference and arithmetic reasoning, structured prompting improves accuracy and reduces variance as the example count grows, making it a powerful enabler for data-hungry structured applications like code generation or multi-step reasoning.

Finally, the leap from atomic skills to complex, compositional reasoning---common in structured tasks such as math word problems---has been investigated through the lens of ICL. \cite{ref172} probes whether models can spontaneously combine atomic skills (e.g., arithmetic and unit conversion) during ICL, finding that they generally cannot without explicit training signals. By proposing a hierarchical curriculum learning strategy that gradually exposes models to compositional tasks, the work induces generalization to complex reasoning scenarios, including cross-dataset and cross-domain transfers. While not a pure ICL solution, this study highlights that the inductive biases learned during pretraining strongly affect how effectively ICL leverages demonstrations for structured, multi-skill reasoning, reinforcing the interplay between pretraining recipes and downstream ICL performance.

In summary, ICL's deployment across NLP and structured tasks is far from one-size-fits-all. Syntax-sensitive selection, coherency metrics, task-specific evaluation metrics, decomposition prompting, and automatic example generation all serve to align demonstrations with the underlying task structure, significantly boosting performance. As new tasks are tackled, these adaptation principles provide a roadmap for transforming generic ICL into a precision instrument for domain-specific reasoning.

However, the very same dependence on careful prompt engineering and demonstration curation that enables these gains also renders ICL remarkably brittle. Even minor perturbations in example ordering, wording, or label distribution can cause large and often unpredictable output swings, exposing vulnerabilities that can be exploited by malicious actors. This fragility and its implications for adversarial manipulation are the subject of the next subsection.

\subsection{Robustness, Sensitivity, and Adversarial Vulnerabilities}

As the previous subsection underscored, the very reliance on prompt engineering and demonstration curation that enables ICL's adaptability also renders it remarkably brittle. In-context learning (ICL) exhibits a pronounced fragility to subtle variations in prompt construction, demonstration ordering, and label presentation, which can lead to unpredictable and often degraded performance. Even benign changes such as shuffling the sequence of exemplars, minor rewording, or alterations in label distribution can cause substantial output variance, revealing a sensitivity that undermines reliability in practice. This brittleness is not merely a nuisance; it opens the door to deliberate adversarial manipulation. As LLMs are increasingly deployed with ICL-based prompts, understanding these vulnerabilities becomes critical for safety and trustworthiness.

These vulnerabilities crystallize along several threat vectors. A particularly insidious threat is data poisoning, where an adversary contaminates the pool of candidate demonstrations to sabotage ICL. The ICLPoison framework \cite{ref173} demonstrates that discrete text perturbations---such as replacing a few tokens in otherwise correct examples---can strategically influence the hidden states of LLMs during the ICL process. By crafting poisoned examples that appear benign but systematically mislead the model's internal representations, an attacker can significantly degrade performance on downstream tasks. Experiments across diverse models, including the robust GPT-4, confirm that such attacks severely compromise ICL efficacy, highlighting the urgent need for defense mechanisms that can sanitize demonstration corpora before injection.

While data poisoning corrupts the content of demonstrations, attackers can also hijack model outputs without modifying the original prompt content, simply by appending imperceptible adversarial suffixes to in-context examples. The LLM hijacking attack \cite{ref174} uses a gradient-based prompt search to learn short, human-imperceptible suffix strings that, when concatenated to demonstrations, divert the model's attention and cause it to generate a targeted malicious response. This method requires no alteration of the original prompt content, making it particularly stealthy; the adversarial suffix merely needs to be added to one or more examples. The attack exploits the autoregressive nature of transformers, showing that ICL's context-driven computation can be coerced into following adversarial directives even when the surface demonstration seems unchanged.

Jailbreaking---eliciting disallowed or harmful content---extends this threat landscape by weaponizing the very mechanism that enables few-shot adaptation. The observation that few-shot demonstrations can both degrade and reinforce safety alignment is central to the In-Context Attack (ICA) and In-Context Defense (ICD) schemes \cite{ref175}. By providing several crafted examples of the model answering malicious queries, ICA can raise the probability of a successful jailbreak without any fine-tuning. Conversely, ICD leverages demonstrations of safe refusals to bolster the model's resistance to harmful prompts. This dual-use nature underscores how the same few-shot mechanism that enables rapid adaptation can also be weaponized to bypass safety filters. Advanced jailbreak techniques further harness ICL for evasion: DrAttack \cite{ref176} decomposes a harmful query into seemingly innocuous sub-prompts and then employs ICL to implicitly reconstruct and execute the original malicious intent. By presenting semantically similar but harmless reassembling demonstrations, DrAttack obfuscates its true purpose, achieving high attack success rates even against GPT-4 while using remarkably few queries. Such attacks reveal that ICL can be exploited not only through direct example contamination but also by restructuring the entire prompt-response paradigm to hide adversarial goals.

Beyond these explicit adversarial manipulations, ICL's fragility also manifests in its sensitivity to label biases and prompt formatting errors. ICL is known to be susceptible to majority label bias, where the distribution of labels in the demonstrations strongly influences predictions, sometimes overriding the correct input-label mapping. Similarly, subtle prompt perturbations---like adding extra spaces or changing delimiters---can produce large output swings. These sensitivities make ICL-based systems vulnerable to adversarial inputs crafted solely to flip predictions without appearing malicious. For instance, an adversary might inject a carefully chosen sequence of tokens into a prompt to skew sentiment classification or to force a specific answer. Because ICL lacks an explicit training phase, the model cannot learn to ignore such nuisance variations; instead, each prompt stands as a complete specification of the learning task, with no separation between ``training'' and ``test'' that would allow intrinsic robustness.

A promising direction to mitigate these risks is sensitivity-aware selective prediction, where the system assesses the reliability of its own outputs and abstains when confidence is low or when the prompt appears adversarially constructed. Defenses based on perplexity filtering, as explored in \cite{ref177}, show that adversarial suffixes often produce sequences with abnormally high perplexity, which can be flagged by a lightweight classifier. However, plain perplexity alone generates many false positives on legitimate long-form prompts, necessitating additional features. More sophisticated defense strategies, such as SafeDecoding \cite{ref178}, intervene during token generation by amplifying the probabilities of safety-related tokens while suppressing harmful sequences. Although originally designed for general jailbreak threats, such safety-aware decoding can also blunt the impact of adversarial ICL by preventing the model from completing malicious instructions even when the context is poisoned.

Ultimately, the robustness of ICL is challenged on multiple fronts: vulnerability to demonstration ordering and label biases, susceptibility to data poisoning and adversarial suffix hijacking, and the ability of few-shot examples to manipulate safety guardrails. The arms race between attack and defense requires that future ICL deployments incorporate multi-layered safeguards. These could include demonstration sanitization using influence-based filtering, runtime perplexity checks, safety-aware decoding, and perhaps training-free selective prediction modules that estimate prompt integrity. Only by acknowledging and systematically hardening ICL against these vulnerabilities can the paradigm be trusted in high-stakes applications.

\subsection{ICL vs.~Fine-tuning: Trade-offs, Complementarity, and Hybrid Approaches}

In light of the robustness challenges detailed in the previous section---where adversarial demonstrations, data poisoning, and prompt perturbations can derail in-context learners---it is essential to weigh ICL against alternative adaptation paradigms that may offer more stability. In-context learning (ICL) and parameter-update-based methods (including full fine-tuning, supervised instruction tuning (SIT), and parameter-efficient fine-tuning (PEFT) such as adapters, prompt tuning, and low-rank adaptation) represent two fundamentally different routes for specializing large language models. ICL consumes a small set of input-output demonstrations prepended to the query, without any gradient-based optimization, whereas fine-tuning updates a subset (or all) of the model parameters using task-specific training data. The practical choice between these paradigms involves navigating multidimensional trade-offs in performance, computational and storage cost, robustness, calibration, and ease of deployment---trade-offs that directly intersect with the safety and generalizability concerns that will be examined further in the next section. This subsection systematically contrasts ICL with fine-tuning approaches, highlights their inherent complementarities, and surveys hybrid methods that aim to fuse the flexibility of in-context learning with the stability and efficiency of parameter updates, thereby pointing toward more resilient adaptation strategies.

\textbf{Performance and Computational Trade-offs}\\
Empirical comparisons consistently show that few-shot PEFT can outperform ICL on unseen tasks while requiring dramatically lower computational overhead per inference. For example, the (IA)³ method together with the T-Few recipe demonstrated that training only a small set of scaling vectors yields higher accuracy than ICL across multiple held-out benchmarks, and does so without storing or processing in-context examples at prediction time \cite{ref179}. ICL incurs substantial memory and latency costs because all demonstration examples must be encoded in every forward pass; the longer the context, the higher the quadratic complexity of self-attention. In contrast, PEFT stores only a small set of adapted parameters, making inference lightweight and suitable for high-throughput production settings. Yet ICL retains a critical advantage when training data is extremely scarce or when no model update can be performed---e.g., due to access restrictions or privacy concerns---since it can adapt purely through prompting. Thus, the performance--cost frontier shifts with the number of available examples: in the very-low-resource regime, ICL is often the only viable choice, while even a handful of training instances can make PEFT a more efficient solution. This tension between instant adaptability and computational economy feeds into the open challenge of long-context efficiency discussed later.

\textbf{Calibration, Bias, and Stability}\\
Both paradigms suffer from miscalibration and overconfidence, especially in low-data scenarios---a problem that compounds the reliability concerns raised in the previous section. A thorough study revealed that simultaneous improvements in task accuracy and calibration are difficult to achieve across ICL, supervised fine-tuning, and instruction tuning, and that low-shot settings exacerbate miscalibration in all approaches \cite{ref92}. For ICL specifically, increasing the number of demonstrations can initially lead to worse calibration before improving, and common usability enhancements such as chain-of-thought prompting or fine-tuning can introduce additional miscalibration \cite{ref88}. ICL is notoriously sensitive to the selection and ordering of in-context examples; random sampling leads to high performance variance. However, this brittleness can be largely mitigated without altering the ICL algorithm itself: data curation methods such as CondAcc and Datamodels select stable subsets of training examples that greatly reduce variance and boost average accuracy by over 6--7\% \cite{ref76}. This finding is particularly relevant given the adversarial ordering attacks outlined earlier; curation-based stabilization can serve as a first line of defense. In contrast, fine-tuning exhibits lower sensitivity to example ordering but can still overfit to spurious correlations in limited data.

Bias is another axis of divergence. Standard ICL inherits the biases present in the few-shot prompts, and unbalanced or biased demonstrations can skew predictions---a vulnerability that adversarial attacks can readily exploit. Fine-tuning with an unbiased validation set can approximate unbiased ICL by learning an optimal example weighting. The reweighted in-context learning (RICL) algorithm fine-tunes a language model to determine example-specific weights that counteract prompt biases, while its low-cost variant LARICL provides a linear approximation that requires minimal training \cite{ref80}. This demonstrates that fine-tuning on a held-out set can serve as a corrective mechanism for ICL's prompt-sensitivity, effectively bridging the two paradigms and offering a blueprint for robustness-enhancing hybrids.

\textbf{Expressiveness and Task Novelty}\\
Theoretical and empirical analyses indicate that ICL and prompt/prefix-tuning are inherently limited in their ability to learn fundamentally new attention patterns. While ICL can simulate known algorithms like gradient descent within the transformer's forward pass, context-based fine-tuning methods cannot alter the relative attention distribution over the input content; they can only bias layer outputs in a fixed direction \cite{ref180}. This means ICL excels at eliciting skills already present in the pretrained model but may struggle with tasks requiring entirely novel attention behaviours, a weakness that becomes critical when confronting out-of-distribution scenarios (as explored in the next section). Full fine-tuning, in contrast, can reshape internal representations more profoundly, which explains why it often yields higher ultimate accuracy when sufficient data is available. Nonetheless, PEFT methods that adapt small parts of the network can still outperform ICL while being more expressive than pure prompt-tuning, striking a balance between learnability and efficiency that is essential for safe and consistent multi-task reasoning.

\textbf{Complementarity and the Emergence of Hybrid Methods}\\
ICL and fine-tuning are not purely competitive; they can be synergistically combined to address each other's shortcomings. The observation that ICL acts as implicit instruction tuning---where demonstrations shift hidden states in the same directions as parameter updates during supervised instruction tuning---lays a conceptual foundation for hybrid designs \cite{ref181}. Building on this, instruction prompt tuning (IPT) concatenates a natural-language demonstration with learned soft-prompt embeddings, merging the runtime flexibility of ICL with the trained adaptation of prompt tuning. Experiments show that IPT reduces the high variance of prompt tuning and benefits from demonstrations that are semantically similar to the test input, while also enabling positive transfer when in-context examples come from a different target task \cite{ref182}.

More integrative fusions appear in the FIAT paradigm, which leverages prompt-engineered instructions and chain-of-thought reasoning from very large models while simultaneously applying parameter-efficient updates to a smaller, locally tunable model. FIAT outperforms both ICL and fine-tuning across multilingual tasks over a wide range of training set sizes (100--10,000 examples), effectively harnessing the strengths of each paradigm without a hard commitment \cite{ref183}.

Other hybridization strategies include using fine-tuning to improve ICL directly: for instance, fine-tuning on supportive pretraining data or on high-quality curated examples can enhance a model's in-context capabilities, as seen in instruction meta-learning frameworks like OPT-IML \cite{ref110}. Moreover, the reweighted ICL family (RICL, LARICL) repurposes fine-tuning as a debiasing pre-processor, making ICL more robust to prompt engineering \cite{ref80}. These hybrids address the fundamental trade-off: ICL provides zero-training-cost adaptivity but suffers from prompt-sensitivity and high inference cost; fine-tuning offers stable, efficient inference at the price of a training phase and potential overfitting. By combining learned components with in-context demonstrations, the emerging hybrid methods enable practitioners to balance adaptability, computational efficiency, and performance, while directly mitigating several of the robustness concerns---such as prompt-order fragility and adversarial label biases---that were highlighted earlier.

\textbf{Concluding Remarks}\\
The choice between ICL and fine-tuning is not binary; the two paradigms inhabit overlapping regions of the performance-efficiency landscape and are increasingly fused into unified workflows. ICL's appeal lies in its instant adaptability and minimal infrastructure, while fine-tuning---especially parameter-efficient variants---delivers superior performance-per-inference-FLOP and stability. Hybrid approaches that embed task-specific parameters into prompts or pre-correct ICL's biases through light-weight training represent a promising path toward models that combine the best of both worlds: the rapid, demonstration-driven generalization of ICL with the robustness and economy of learned adaptations. This convergence directly informs the open challenges that follow, as future work on calibration-aware demonstration design, cross-modal reliability, and safety alignment can be grounded in the complementary strengths and ongoing convergence of ICL and fine-tuning.

\subsection{Open Problems and Future Directions}

Building on the convergent paradigm described in the previous section---where hybrids of in-context learning and parameter-efficient fine-tuning offer new resilience---this section maps the fundamental challenges that remain for deploying ICL in safety-critical and rapidly shifting environments. We delineate open problems and future directions spanning robust out-of-distribution generalization, consistent multi-task reasoning, safety alignment against unwanted task re-learning, long-context efficiency, cross-modal reliability, and the design of training-free, calibration-aware demonstration strategies that together could yield trustworthy, generalizable in-context learners.

\textbf{Out-of-Distribution Generalization.}\\
ICL often fails when test distributions diverge from pretraining or in-context demonstrations, as demonstrated empirically in both text and multimodal settings. In syntax-sensitive tasks, performance varies dramatically across LMs and is more strongly influenced by pretraining corpus composition---code-trained models generalize better---than by model scale \cite{ref184}. Similarly, multimodal LMs suffer from mapping deficiency when tested on synthetic or domain-specific images, and although ICL can improve zero-shot performance, the approach remains highly vulnerable to domain shifts, label shifts, and spurious correlations between demonstrations and test data \cite{ref185}. A key open problem is the development of ICL procedures that can explicitly detect and compensate for distributional mismatches, perhaps by combining instance-level hardness measures like in-context PVI \cite{ref186} with uncertainty-aware demonstration selection, thereby enabling models to reject or reframe prompts when the context is unreliable.

\textbf{Consistent Multi-Task Reasoning.}\\
Current ICL methods struggle with maintaining a coherent ability across diverse tasks, especially when the same context requires integrating multiple reasoning patterns. Multi-task training curricula have been shown to improve convergence and data efficiency \cite{ref187}, yet it remains unclear how to construct demonstrations that simultaneously evoke arithmetic, logical, and commonsense reasoning without catastrophic interference. While concept-aware training data construction encourages analogical reasoning \cite{ref109}, moving to a future where a single in-context prompt can reliably switch between disparate task families---without the model falling back on shallow heuristics or semantic priors---will require a deeper understanding of how to modularize internal representations and when to prefer in-context plasticity over prior knowledge.

\textbf{Safety Alignment and Forbidden Task Re-Learning.}\\
ICL can inadvertently override safety alignment when demonstrations illustrate prohibited behaviors or when the model's strong semantic priors dominate the provided input-label mappings. The phenomenon of ``demonstration shortcut'', where LLMs rely on pre-trained semantic associations rather than learning from the given examples, is pervasive and can be partly rectified by calibration methods such as In-Context Calibration \cite{ref39}. Even more concerning, in affective tasks the pull of prior knowledge grows with model size, ossifying predictions and preventing the model from adapting to genuinely new labeling conventions \cite{ref42}. This suggests that larger models may be more susceptible to ignoring demonstrations that conflict with their internal world model. A critical future direction is to design demonstration-aware alignment techniques that preserve safety while allowing genuine adaptation, possibly by explicitly modeling the trade-off between prior fidelity and in-context flexibility and developing demonstration selection strategies that can detect when the prior is harmful.

\textbf{Long-Context Efficiency and Calibration.}\\
As the number of demonstrations grows, ICL faces practical limitations due to finite context windows and computational cost. The calibration behavior is also non-monotonic: models often become more miscalibrated in low-shot regimes before improving, and chain-of-thought prompting can exacerbate overconfidence \cite{ref88}. While recalibration methods like scaling-binning can reduce errors, a more fundamental solution is needed. One promising direction is to decouple inference from context length entirely, as demonstrated by \(k\)NN Prompting, which leverages training data representations to achieve calibration-free scaling up to thousands of shots without being bounded by the context window \cite{ref188}. However, integrating such retrieval-based paradigms with standard ICL pipelines without sacrificing the ease of a single forward pass remains an open challenge. Future work should explore hybrid architectures that dynamically combine internal in-weight knowledge, in-context activations, and external memory to achieve both long-range dependency and reliable uncertainty estimates.

\textbf{Cross-Modal ICL Reliability.}\\
Extending ICL to vision-language and other multimodal contexts introduces new failure modes rooted in representation disparities and information bottlenecks between modalities. Visual demonstrations often contain spurious correlations that mislead the model, and cross-modal interactions can degrade performance even when individual modalities are informative \cite{ref185}. Recent visual ICL methods that employ intent-oriented image summarization and language-based demonstration composition \cite{ref159} represent steps forward, but robust generalization across visual domain shifts and modality-specific adversarial perturbations remains unsolved. A crucial direction is to build unified multimodal ICL frameworks with built-in cross-modal alignment checks, perhaps via mutual information barriers that prevent a dominant modality from overriding the other when it is less trustworthy.

\textbf{Training-Free, Calibration-Aware Demonstration Strategies.}\\
A recurring limitation is the high sensitivity of ICL to the choice and ordering of demonstrations, which can cause large performance variance. Data curation techniques like CondAcc and Datamodels can significantly stabilize ICL by preselecting stable example subsets \cite{ref76}, but these methods still require labeled development data and are not fully training-free. Iterative misconfidence-based selection \cite{ref189} and incremental utility labeling \cite{ref190} offer promising paths toward adaptive selection that does not rely on external metrics, yet they depend on repeated LLM calls. The ultimate goal is a lightweight, on-the-fly demonstration construction algorithm that simultaneously accounts for calibration, task priors, and distribution shift, much like a real-time meta-learner. Integrating in-context PVI to estimate instance hardness without fine-tuning \cite{ref186} with calibration-aware weighting could produce demonstration sets that not only maximize task accuracy but also yield trustworthy confidence estimates.

In summary, bridging the gap between the current ad-hoc usage of ICL and a principled, safe, and universally applicable paradigm requires a concerted effort across these fronts. Advances in each area will likely inform the others: calibration methods can improve OOD detection, long-context scalability enables richer multi-task demonstrations, and robust demonstration design can mitigate safety failures. Only through such integrated progress will ICL evolve into a reliable and trustworthy foundation for next-generation AI systems.


\begin{thebibliography}{190}
\addcontentsline{toc}{section}{References}
\bibitem{ref1} What Can Transformers Learn In-Context A Case Study of Simple Function Classes. arXiv:2208.01066v3, 2022. \url{https://arxiv.org/abs/2208.01066v3}
\bibitem{ref2} In-Context Learning through the Bayesian Prism. arXiv:2306.04891v2, 2023. \url{https://arxiv.org/abs/2306.04891v2}
\bibitem{ref3} General-Purpose In-Context Learning by Meta-Learning Transformers. arXiv:2212.04458v2, 2022. \url{https://arxiv.org/abs/2212.04458v2}
\bibitem{ref4} Transformers as Algorithms Generalization and Stability in In-context Learning. arXiv:2301.07067v2, 2023. \url{https://arxiv.org/abs/2301.07067v2}
\bibitem{ref5} What learning algorithm is in-context learning Investigations with linear models. arXiv:2211.15661v3, 2022. \url{https://arxiv.org/abs/2211.15661v3}
\bibitem{ref6} Transformers Learn Higher-Order Optimization Methods for In-Context Learning A Study with Linear Models. arXiv:2310.17086v1, 2023. \url{https://arxiv.org/abs/2310.17086v1}
\bibitem{ref7} Transformers as Statisticians Provable In-Context Learning with In-Context Algorithm Selection. arXiv:2306.04637v2, 2023. \url{https://arxiv.org/abs/2306.04637v2}
\bibitem{ref8} Understanding In-Context Learning in Transformers and LLMs by Learning to Learn Discrete Functions. arXiv:2310.03016v1, 2023. \url{https://arxiv.org/abs/2310.03016v1}
\bibitem{ref9} How Do Transformers Learn In-Context Beyond Simple Functions A Case Study on Learning with Representations. arXiv:2310.10616v1, 2023. \url{https://arxiv.org/abs/2310.10616v1}
\bibitem{ref10} What and How does In-Context Learning Learn Bayesian Model Averaging, Parameterization, and Generalization. arXiv:2305.19420v2, 2023. \url{https://arxiv.org/abs/2305.19420v2}
\bibitem{ref11} Large Language Models Are Latent Variable Models Explaining and Finding Good Demonstrations for In-Context Learning. arXiv:2301.11916v4, 2023. \url{https://arxiv.org/abs/2301.11916v4}
\bibitem{ref12} Why Can GPT Learn In-Context Language Models Implicitly Perform Gradient Descent as Meta-Optimizers. arXiv:2212.10559v3, 2022. \url{https://arxiv.org/abs/2212.10559v3}
\bibitem{ref13} In-context Learning with Transformer Is Really Equivalent to a Contrastive Learning Pattern. arXiv:2310.13220v1, 2023. \url{https://arxiv.org/abs/2310.13220v1}
\bibitem{ref14} Trained Transformers Learn Linear Models In-Context. arXiv:2306.09927v3, 2023. \url{https://arxiv.org/abs/2306.09927v3}
\bibitem{ref15} Linear Transformers are Versatile In-Context Learners. arXiv:2402.14180v1, 2024. \url{https://arxiv.org/abs/2402.14180v1}
\bibitem{ref16} Transformers Learn Nonlinear Features In Context Nonconvex Mean-field Dynamics on the Attention Landscape. arXiv:2402.01258v1, 2024. \url{https://arxiv.org/abs/2402.01258v1}
\bibitem{ref17} Implicit Bias and Fast Convergence Rates for Self-attention. arXiv:2402.05738v1, 2024. \url{https://arxiv.org/abs/2402.05738v1}
\bibitem{ref18} In-Context Convergence of Transformers. arXiv:2310.05249v1, 2023. \url{https://arxiv.org/abs/2310.05249v1}
\bibitem{ref19} How Transformers Learn Causal Structure with Gradient Descent. arXiv:2402.14735v1, 2024. \url{https://arxiv.org/abs/2402.14735v1}
\bibitem{ref20} In-Context Learning with Transformers Softmax Attention Adapts to Function Lipschitzness. arXiv:2402.11639v1, 2024. \url{https://arxiv.org/abs/2402.11639v1}
\bibitem{ref21} The Learnability of In-Context Learning. arXiv:2303.07895v1, 2023. \url{https://arxiv.org/abs/2303.07895v1}
\bibitem{ref22} A General framework for PAC-Bayes Bounds for Meta-Learning. arXiv:2206.05454v1, 2022. \url{https://arxiv.org/abs/2206.05454v1}
\bibitem{ref23} Generalization Bounds for Meta-Learning via PAC-Bayes and Uniform Stability. arXiv:2102.06589v3, 2021. \url{https://arxiv.org/abs/2102.06589v3}
\bibitem{ref24} PAC-Bayes Bounds for Meta-learning with Data-Dependent Prior. arXiv:2102.03748v1, 2021. \url{https://arxiv.org/abs/2102.03748v1}
\bibitem{ref25} PAC-Bayes meta-learning with implicit task-specific posteriors. arXiv:2003.02455v3, 2020. \url{https://arxiv.org/abs/2003.02455v3}
\bibitem{ref26} How Many Pretraining Tasks Are Needed for In-Context Learning of Linear Regression. arXiv:2310.08391v2, 2023. \url{https://arxiv.org/abs/2310.08391v2}
\bibitem{ref27} An Information-Theoretic Analysis of the Impact of Task Similarity on Meta-Learning. arXiv:2101.08390v3, 2021. \url{https://arxiv.org/abs/2101.08390v3}
\bibitem{ref28} Transfer Meta-Learning Information-Theoretic Bounds and Information Meta-Risk Minimization. arXiv:2011.02872v2, 2020. \url{https://arxiv.org/abs/2011.02872v2}
\bibitem{ref29} Model-agnostic Measure of Generalization Difficulty. arXiv:2305.01034v2, 2023. \url{https://arxiv.org/abs/2305.01034v2}
\bibitem{ref30} An analytic theory of generalization dynamics and transfer learning in deep linear networks. arXiv:1809.10374v2, 2018. \url{https://arxiv.org/abs/1809.10374v2}
\bibitem{ref31} Pretraining task diversity and the emergence of non-Bayesian in-context learning for regression. arXiv:2306.15063v2, 2023. \url{https://arxiv.org/abs/2306.15063v2}
\bibitem{ref32} Beyond Scale the Diversity Coefficient as a Data Quality Metric Demonstrates LLMs are Pre-trained on Formally Diverse Data. arXiv:2306.13840v2, 2023. \url{https://arxiv.org/abs/2306.13840v2}
\bibitem{ref33} Parallel Structures in Pre-training Data Yield In-Context Learning. arXiv:2402.12530v1, 2024. \url{https://arxiv.org/abs/2402.12530v1}
\bibitem{ref34} A Theory of Emergent In-Context Learning as Implicit Structure Induction. arXiv:2303.07971v1, 2023. \url{https://arxiv.org/abs/2303.07971v1}
\bibitem{ref35} The mechanistic basis of data dependence and abrupt learning in an in-context classification task. arXiv:2312.03002v1, 2023. \url{https://arxiv.org/abs/2312.03002v1}
\bibitem{ref36} Understanding In-Context Learning via Supportive Pretraining Data. arXiv:2306.15091v1, 2023. \url{https://arxiv.org/abs/2306.15091v1}
\bibitem{ref37} On the Effect of Pretraining Corpora on In-context Learning by a Large-scale Language Model. arXiv:2204.13509v2, 2022. \url{https://arxiv.org/abs/2204.13509v2}
\bibitem{ref38} Concept-aware Training Improves In-context Learning Ability of Language Models. arXiv:2305.13775v1, 2023. \url{https://arxiv.org/abs/2305.13775v1}
\bibitem{ref39} Rectifying Demonstration Shortcut in In-Context Learning. arXiv:2403.09488v3, 2024. \url{https://arxiv.org/abs/2403.09488v3}
\bibitem{ref40} Larger language models do in-context learning differently. arXiv:2303.03846v2, 2023. \url{https://arxiv.org/abs/2303.03846v2}
\bibitem{ref41} In-Context Learning Learns Label Relationships but Is Not Conventional Learning. arXiv:2307.12375v4, 2023. \url{https://arxiv.org/abs/2307.12375v4}
\bibitem{ref42} The Strong Pull of Prior Knowledge in Large Language Models and Its Impact on Emotion Recognition. arXiv:2403.17125v1, 2024. \url{https://arxiv.org/abs/2403.17125v1}
\bibitem{ref43} Ground-Truth Labels Matter A Deeper Look into Input-Label Demonstrations. arXiv:2205.12685v2, 2022. \url{https://arxiv.org/abs/2205.12685v2}
\bibitem{ref44} Towards Understanding In-Context Learning with Contrastive Demonstrations and Saliency Maps. arXiv:2307.05052v4, 2023. \url{https://arxiv.org/abs/2307.05052v4}
\bibitem{ref45} Decomposing Label Space, Format and Discrimination Rethinking How LLMs Respond and Solve Tasks via In-Context Learning. arXiv:2404.07546v1, 2024. \url{https://arxiv.org/abs/2404.07546v1}
\bibitem{ref46} Comparable Demonstrations are Important in In-Context Learning A Novel Perspective on Demonstration Selection. arXiv:2312.07476v2, 2023. \url{https://arxiv.org/abs/2312.07476v2}
\bibitem{ref47} Data. arXiv:1801.04992v2, 2017. \url{https://arxiv.org/abs/1801.04992v2}
\bibitem{ref48} B-Splines. arXiv:2108.06617v1, 2021. \url{https://arxiv.org/abs/2108.06617v1}
\bibitem{ref49} In-context Learning and Induction Heads. arXiv:2209.11895v1, 2022. \url{https://arxiv.org/abs/2209.11895v1}
\bibitem{ref50} The Evolution of Statistical Induction Heads In-Context Learning Markov Chains. arXiv:2402.11004v1, 2024. \url{https://arxiv.org/abs/2402.11004v1}
\bibitem{ref51} What needs to go right for an induction head A mechanistic study of in-context learning circuits and their formation. arXiv:2404.07129v1, 2024. \url{https://arxiv.org/abs/2404.07129v1}
\bibitem{ref52} In-Context Language Learning Architectures and Algorithms. arXiv:2401.12973v2, 2024. \url{https://arxiv.org/abs/2401.12973v2}
\bibitem{ref53} The Transient Nature of Emergent In-Context Learning in Transformers. arXiv:2311.08360v3, 2023. \url{https://arxiv.org/abs/2311.08360v3}
\bibitem{ref54} Copy Suppression Comprehensively Understanding an Attention Head. arXiv:2310.04625v1, 2023. \url{https://arxiv.org/abs/2310.04625v1}
\bibitem{ref55} One Step of Gradient Descent is Provably the Optimal In-Context Learner with One Layer of Linear Self-Attention. arXiv:2307.03576v1, 2023. \url{https://arxiv.org/abs/2307.03576v1}
\bibitem{ref56} Transformers learn to implement preconditioned gradient descent for in-context learning. arXiv:2306.00297v2, 2023. \url{https://arxiv.org/abs/2306.00297v2}
\bibitem{ref57} The Bayesian Learning Rule. arXiv:2107.04562v3, 2021. \url{https://arxiv.org/abs/2107.04562v3}
\bibitem{ref58} Transformers Implement Functional Gradient Descent to Learn Non-Linear Functions In Context. arXiv:2312.06528v5, 2023. \url{https://arxiv.org/abs/2312.06528v5}
\bibitem{ref59} Uncovering mesa-optimization algorithms in Transformers. arXiv:2309.05858v1, 2023. \url{https://arxiv.org/abs/2309.05858v1}
\bibitem{ref60} Looped Transformers are Better at Learning Learning Algorithms. arXiv:2311.12424v3, 2023. \url{https://arxiv.org/abs/2311.12424v3}
\bibitem{ref61} In-Context Learning Creates Task Vectors. arXiv:2310.15916v1, 2023. \url{https://arxiv.org/abs/2310.15916v1}
\bibitem{ref62} Benefits of Transformer In-Context Learning in Linear Regression Tasks with Unstructured Data. arXiv:2402.00743v1, 2024. \url{https://arxiv.org/abs/2402.00743v1}
\bibitem{ref63} Identifying and Analyzing Task-Encoding Tokens in Large Language Models. arXiv:2401.11323v2, 2024. \url{https://arxiv.org/abs/2401.11323v2}
\bibitem{ref64} How are Prompts Different in Terms of Sensitivity. arXiv:2311.07230v1, 2023. \url{https://arxiv.org/abs/2311.07230v1}
\bibitem{ref65} Iterative Forward Tuning Boosts In-context Learning in Language Models. arXiv:2305.13016v2, 2023. \url{https://arxiv.org/abs/2305.13016v2}
\bibitem{ref66} Human Curriculum Effects Emerge with In-Context Learning in Neural Networks. arXiv:2402.08674v1, 2024. \url{https://arxiv.org/abs/2402.08674v1}
\bibitem{ref67} Dr.ICL Demonstration-Retrieved In-context Learning. arXiv:2305.14128v1, 2023. \url{https://arxiv.org/abs/2305.14128v1}
\bibitem{ref68} Finding Support Examples for In-Context Learning. arXiv:2302.13539v3, 2023. \url{https://arxiv.org/abs/2302.13539v3}
\bibitem{ref69} Coverage-based Example Selection for In-Context Learning. arXiv:2305.14907v3, 2023. \url{https://arxiv.org/abs/2305.14907v3}
\bibitem{ref70} Compositional Exemplars for In-context Learning. arXiv:2302.05698v3, 2023. \url{https://arxiv.org/abs/2302.05698v3}
\bibitem{ref71} In-context Example Selection with Influences. arXiv:2302.11042v2, 2023. \url{https://arxiv.org/abs/2302.11042v2}
\bibitem{ref72} RetICL Sequential Retrieval of In-Context Examples with Reinforcement Learning. arXiv:2305.14502v2, 2023. \url{https://arxiv.org/abs/2305.14502v2}
\bibitem{ref73} Ambiguity-Aware In-Context Learning with Large Language Models. arXiv:2309.07900v2, 2023. \url{https://arxiv.org/abs/2309.07900v2}
\bibitem{ref74} Crafting In-context Examples according to LMs' Parametric Knowledge. arXiv:2311.09579v2, 2023. \url{https://arxiv.org/abs/2311.09579v2}
\bibitem{ref75} In-Context Learning for Text Classification with Many Labels. arXiv:2309.10954v2, 2023. \url{https://arxiv.org/abs/2309.10954v2}
\bibitem{ref76} Data Curation Alone Can Stabilize In-context Learning. arXiv:2212.10378v2, 2022. \url{https://arxiv.org/abs/2212.10378v2}
\bibitem{ref77} Let's Learn Step by Step Enhancing In-Context Learning Ability with Curriculum Learning. arXiv:2402.10738v1, 2024. \url{https://arxiv.org/abs/2402.10738v1}
\bibitem{ref78} Self-Adaptive In-Context Learning An Information Compression Perspective for In-Context Example Selection and Ordering. arXiv:2212.10375v2, 2022. \url{https://arxiv.org/abs/2212.10375v2}
\bibitem{ref79} Not All Demonstration Examples are Equally Beneficial Reweighting Demonstration Examples for In-Context Learning. arXiv:2310.08309v1, 2023. \url{https://arxiv.org/abs/2310.08309v1}
\bibitem{ref80} Fine-tune Language Models to Approximate Unbiased In-context Learning. arXiv:2310.03331v1, 2023. \url{https://arxiv.org/abs/2310.03331v1}
\bibitem{ref81} Batch-ICL Effective, Efficient, and Order-Agnostic In-Context Learning. arXiv:2401.06469v2, 2024. \url{https://arxiv.org/abs/2401.06469v2}
\bibitem{ref82} InstructEval Systematic Evaluation of Instruction Selection Methods. arXiv:2307.00259v2, 2023. \url{https://arxiv.org/abs/2307.00259v2}
\bibitem{ref83} Hint-enhanced In-Context Learning wakes Large Language Models up for knowledge-intensive tasks. arXiv:2311.01949v2, 2023. \url{https://arxiv.org/abs/2311.01949v2}
\bibitem{ref84} Manipulating the Label Space for In-Context Classification. arXiv:2312.00351v2, 2023. \url{https://arxiv.org/abs/2312.00351v2}
\bibitem{ref85} Mind the instructions a holistic evaluation of consistency and interactions in prompt-based learning. arXiv:2310.13486v1, 2023. \url{https://arxiv.org/abs/2310.13486v1}
\bibitem{ref86} Learning Hierarchical Prompt with Structured Linguistic Knowledge for Vision-Language Models. arXiv:2312.06323v1, 2023. \url{https://arxiv.org/abs/2312.06323v1}
\bibitem{ref87} LAMM Label Alignment for Multi-Modal Prompt Learning. arXiv:2312.08212v1, 2023. \url{https://arxiv.org/abs/2312.08212v1}
\bibitem{ref88} A Study on the Calibration of In-context Learning. arXiv:2312.04021v4, 2023. \url{https://arxiv.org/abs/2312.04021v4}
\bibitem{ref89} Batch Calibration Rethinking Calibration for In-Context Learning and Prompt Engineering. arXiv:2309.17249v2, 2023. \url{https://arxiv.org/abs/2309.17249v2}
\bibitem{ref90} Mitigating Label Biases for In-context Learning. arXiv:2305.19148v3, 2023. \url{https://arxiv.org/abs/2305.19148v3}
\bibitem{ref91} NoisyICL A Little Noise in Model Parameters Calibrates In-context Learning. arXiv:2402.05515v2, 2024. \url{https://arxiv.org/abs/2402.05515v2}
\bibitem{ref92} On Task Performance and Model Calibration with Supervised and Self-Ensembled In-Context Learning. arXiv:2312.13772v2, 2023. \url{https://arxiv.org/abs/2312.13772v2}
\bibitem{ref93} Calibration Meets Explanation A Simple and Effective Approach for Model Confidence Estimates. arXiv:2211.03041v1, 2022. \url{https://arxiv.org/abs/2211.03041v1}
\bibitem{ref94} Self-ICL Zero-Shot In-Context Learning with Self-Generated Demonstrations. arXiv:2305.15035v2, 2023. \url{https://arxiv.org/abs/2305.15035v2}
\bibitem{ref95} Self-Generated In-Context Learning Leveraging Auto-regressive Language Models as a Demonstration Generator. arXiv:2206.08082v1, 2022. \url{https://arxiv.org/abs/2206.08082v1}
\bibitem{ref96} Are Human-generated Demonstrations Necessary for In-context Learning. arXiv:2309.14681v4, 2023. \url{https://arxiv.org/abs/2309.14681v4}
\bibitem{ref97} Self-Demos Eliciting Out-of-Demonstration Generalizability in Large Language Models. arXiv:2404.00884v1, 2024. \url{https://arxiv.org/abs/2404.00884v1}
\bibitem{ref98} Self-AMPLIFY Improving Small Language Models with Self Post Hoc Explanations. arXiv:2402.12038v2, 2024. \url{https://arxiv.org/abs/2402.12038v2}
\bibitem{ref99} ZARA Improving Few-Shot Self-Rationalization for Small Language Models. arXiv:2305.07355v2, 2023. \url{https://arxiv.org/abs/2305.07355v2}
\bibitem{ref100} SELF-EXPLAIN Teaching Large Language Models to Reason Complex Questions by Themselves. arXiv:2311.06985v1, 2023. \url{https://arxiv.org/abs/2311.06985v1}
\bibitem{ref101} Many-Shot In-Context Learning. arXiv:2404.11018v1, 2024. \url{https://arxiv.org/abs/2404.11018v1}
\bibitem{ref102} Ensemble-Instruct Generating Instruction-Tuning Data with a Heterogeneous Mixture of LMs. arXiv:2310.13961v1, 2023. \url{https://arxiv.org/abs/2310.13961v1}
\bibitem{ref103} AutoHint Automatic Prompt Optimization with Hint Generation. arXiv:2307.07415v2, 2023. \url{https://arxiv.org/abs/2307.07415v2}
\bibitem{ref104} Toward Self-Improvement of LLMs via Imagination, Searching, and Criticizing. arXiv:2404.12253v1, 2024. \url{https://arxiv.org/abs/2404.12253v1}
\bibitem{ref105} Scaling In-Context Demonstrations with Structured Attention. arXiv:2307.02690v1, 2023. \url{https://arxiv.org/abs/2307.02690v1}
\bibitem{ref106} XC-Cache Cross-Attending to Cached Context for Efficient LLM Inference. arXiv:2404.15420v1, 2024. \url{https://arxiv.org/abs/2404.15420v1}
\bibitem{ref107} Lightweight In-Context Tuning for Multimodal Unified Models. arXiv:2310.05109v1, 2023. \url{https://arxiv.org/abs/2310.05109v1}
\bibitem{ref108} MMICT Boosting Multi-Modal Fine-Tuning with In-Context Examples. arXiv:2312.06363v2, 2023. \url{https://arxiv.org/abs/2312.06363v2}
\bibitem{ref109} Concept-aware Data Construction Improves In-context Learning of Language Models. arXiv:2403.09703v1, 2024. \url{https://arxiv.org/abs/2403.09703v1}
\bibitem{ref110} OPT-IML Scaling Language Model Instruction Meta Learning through the Lens of Generalization. arXiv:2212.12017v3, 2022. \url{https://arxiv.org/abs/2212.12017v3}
\bibitem{ref111} Knowledgeable In-Context Tuning Exploring and Exploiting Factual Knowledge for In-Context Learning. arXiv:2309.14771v2, 2023. \url{https://arxiv.org/abs/2309.14771v2}
\bibitem{ref112} Exploring Effective Factors for Improving Visual In-Context Learning. arXiv:2304.04748v1, 2023. \url{https://arxiv.org/abs/2304.04748v1}
\bibitem{ref113} Instruct Me More! Random Prompting for Visual In-Context Learning. arXiv:2311.03648v1, 2023. \url{https://arxiv.org/abs/2311.03648v1}
\bibitem{ref114} A Simple Framework Uniting Visual In-context Learning with Masked Image Modeling to Improve Ultrasound Segmentation. arXiv:2402.14300v3, 2024. \url{https://arxiv.org/abs/2402.14300v3}
\bibitem{ref115} SegICL A Universal In-context Learning Framework for Enhanced Segmentation in Medical Imaging. arXiv:2403.16578v2, 2024. \url{https://arxiv.org/abs/2403.16578v2}
\bibitem{ref116} Understanding and Improving In-Context Learning on Vision-language Models. arXiv:2311.18021v1, 2023. \url{https://arxiv.org/abs/2311.18021v1}
\bibitem{ref117} Exploring Diverse In-Context Configurations for Image Captioning. arXiv:2305.14800v6, 2023. \url{https://arxiv.org/abs/2305.14800v6}
\bibitem{ref118} All in an Aggregated Image for In-Image Learning. arXiv:2402.17971v2, 2024. \url{https://arxiv.org/abs/2402.17971v2}
\bibitem{ref119} MMICL Empowering Vision-language Model with Multi-Modal In-Context Learning. arXiv:2309.07915v3, 2023. \url{https://arxiv.org/abs/2309.07915v3}
\bibitem{ref120} Towards Multimodal In-Context Learning for Vision \& Language Models. arXiv:2403.12736v1, 2024. \url{https://arxiv.org/abs/2403.12736v1}
\bibitem{ref121} UNIMO-3 Multi-granularity Interaction for Vision-Language Representation Learning. arXiv:2305.13697v1, 2023. \url{https://arxiv.org/abs/2305.13697v1}
\bibitem{ref122} Deeply Coupled Cross-Modal Prompt Learning. arXiv:2305.17903v3, 2023. \url{https://arxiv.org/abs/2305.17903v3}
\bibitem{ref123} Lost in Translation When GPT-4V(ision) Can't See Eye to Eye with Text. A Vision-Language-Consistency Analysis of VLLMs and Beyond. arXiv:2310.12520v1, 2023. \url{https://arxiv.org/abs/2310.12520v1}
\bibitem{ref124} VL-ICL Bench The Devil in the Details of Benchmarking Multimodal In-Context Learning. arXiv:2403.13164v1, 2024. \url{https://arxiv.org/abs/2403.13164v1}
\bibitem{ref125} Can MLLMs Perform Text-to-Image In-Context Learning. arXiv:2402.01293v2, 2024. \url{https://arxiv.org/abs/2402.01293v2}
\bibitem{ref126} Towards More Unified In-context Visual Understanding. arXiv:2312.02520v2, 2023. \url{https://arxiv.org/abs/2312.02520v2}
\bibitem{ref127} Multilingual LLMs are Better Cross-lingual In-context Learners with Alignment. arXiv:2305.05940v3, 2023. \url{https://arxiv.org/abs/2305.05940v3}
\bibitem{ref128} LLMs Are Few-Shot In-Context Low-Resource Language Learners. arXiv:2403.16512v2, 2024. \url{https://arxiv.org/abs/2403.16512v2}
\bibitem{ref129} Cross-Lingual Retrieval Augmented Prompt for Low-Resource Languages. arXiv:2212.09651v4, 2022. \url{https://arxiv.org/abs/2212.09651v4}
\bibitem{ref130} From Classification to Generation Insights into Crosslingual Retrieval Augmented ICL. arXiv:2311.06595v3, 2023. \url{https://arxiv.org/abs/2311.06595v3}
\bibitem{ref131} Resources and Few-shot Learners for In-context Learning in Slavic Languages. arXiv:2304.01922v1, 2023. \url{https://arxiv.org/abs/2304.01922v1}
\bibitem{ref132} Extrapolating Large Language Models to Non-English by Aligning Languages. arXiv:2308.04948v2, 2023. \url{https://arxiv.org/abs/2308.04948v2}
\bibitem{ref133} CrossIn An Efficient Instruction Tuning Approach for Cross-Lingual Knowledge Alignment. arXiv:2404.11932v1, 2024. \url{https://arxiv.org/abs/2404.11932v1}
\bibitem{ref134} Enhancing Cross-lingual Transfer via Phonemic Transcription Integration. arXiv:2307.04361v1, 2023. \url{https://arxiv.org/abs/2307.04361v1}
\bibitem{ref135} ICMSC Intra- and Cross-modality Semantic Consistency for Unsupervised Domain Adaptation on Hip Joint Bone Segmentation. arXiv:2012.12570v1, 2020. \url{https://arxiv.org/abs/2012.12570v1}
\bibitem{ref136} LE-UDA Label-efficient unsupervised domain adaptation for medical image segmentation. arXiv:2212.02078v1, 2022. \url{https://arxiv.org/abs/2212.02078v1}
\bibitem{ref137} Domain Adaptive Medical Image Segmentation via Adversarial Learning of Disease-Specific Spatial Patterns. arXiv:2001.09313v3, 2020. \url{https://arxiv.org/abs/2001.09313v3}
\bibitem{ref138} Cross-Region Domain Adaptation for Class-level Alignment. arXiv:2109.06422v2, 2021. \url{https://arxiv.org/abs/2109.06422v2}
\bibitem{ref139} Domain-Agnostic Mutual Prompting for Unsupervised Domain Adaptation. arXiv:2403.02899v1, 2024. \url{https://arxiv.org/abs/2403.02899v1}
\bibitem{ref140} Cross-modal Learning for Domain Adaptation in 3D Semantic Segmentation. arXiv:2101.07253v2, 2021. \url{https://arxiv.org/abs/2101.07253v2}
\bibitem{ref141} Sparse-to-dense Feature Matching Intra and Inter domain Cross-modal Learning in Domain Adaptation for 3D Semantic Segmentation. arXiv:2107.14724v5, 2021. \url{https://arxiv.org/abs/2107.14724v5}
\bibitem{ref142} Self-semantic contour adaptation for cross modality brain tumor segmentation. arXiv:2201.05022v1, 2022. \url{https://arxiv.org/abs/2201.05022v1}
\bibitem{ref143} Synergistic Image and Feature Adaptation Towards Cross-Modality Domain Adaptation for Medical Image Segmentation. arXiv:1901.08211v4, 2019. \url{https://arxiv.org/abs/1901.08211v4}
\bibitem{ref144} Meta-hallucinator Towards Few-Shot Cross-Modality Cardiac Image Segmentation. arXiv:2305.06978v1, 2023. \url{https://arxiv.org/abs/2305.06978v1}
\bibitem{ref145} CrDoCo Pixel-level Domain Transfer with Cross-Domain Consistency. arXiv:2001.03182v1, 2020. \url{https://arxiv.org/abs/2001.03182v1}
\bibitem{ref146} Unsupervised Domain Adaptation Network with Category-Centric Prototype Aligner for Biomedical Image Segmentation. arXiv:2103.02220v1, 2021. \url{https://arxiv.org/abs/2103.02220v1}
\bibitem{ref147} Mapping Instructions and Visual Observations to Actions with Reinforcement Learning. arXiv:1704.08795v2, 2017. \url{https://arxiv.org/abs/1704.08795v2}
\bibitem{ref148} Ask Your Humans Using Human Instructions to Improve Generalization in Reinforcement Learning. arXiv:2011.00517v3, 2020. \url{https://arxiv.org/abs/2011.00517v3}
\bibitem{ref149} Multimodal Hierarchical Reinforcement Learning Policy for Task-Oriented Visual Dialog. arXiv:1805.03257v1, 2018. \url{https://arxiv.org/abs/1805.03257v1}
\bibitem{ref150} Multi-modal Feedback for Affordance-driven Interactive Reinforcement Learning. arXiv:1807.09991v1, 2018. \url{https://arxiv.org/abs/1807.09991v1}
\bibitem{ref151} MORE-3S Multimodal-based Offline Reinforcement Learning with Shared Semantic Spaces. arXiv:2402.12845v1, 2024. \url{https://arxiv.org/abs/2402.12845v1}
\bibitem{ref152} Context-Hierarchy Inverse Reinforcement Learning. arXiv:2202.12597v1, 2022. \url{https://arxiv.org/abs/2202.12597v1}
\bibitem{ref153} Noisy Agents Self-supervised Exploration by Predicting Auditory Events. arXiv:2007.13729v1, 2020. \url{https://arxiv.org/abs/2007.13729v1}
\bibitem{ref154} COBRA Contrastive Bi-Modal Representation Algorithm. arXiv:2005.03687v2, 2020. \url{https://arxiv.org/abs/2005.03687v2}
\bibitem{ref155} Vision-Language Pre-Training with Triple Contrastive Learning. arXiv:2202.10401v4, 2022. \url{https://arxiv.org/abs/2202.10401v4}
\bibitem{ref156} Hybrid Contrastive Learning of Tri-Modal Representation for Multimodal Sentiment Analysis. arXiv:2109.01797v1, 2021. \url{https://arxiv.org/abs/2109.01797v1}
\bibitem{ref157} Understanding Dark Scenes by Contrasting Multi-Modal Observations. arXiv:2308.12320v2, 2023. \url{https://arxiv.org/abs/2308.12320v2}
\bibitem{ref158} Switch-BERT Learning to Model Multimodal Interactions by Switching Attention and Input. arXiv:2306.14182v1, 2023. \url{https://arxiv.org/abs/2306.14182v1}
\bibitem{ref159} Visual In-Context Learning for Large Vision-Language Models. arXiv:2402.11574v1, 2024. \url{https://arxiv.org/abs/2402.11574v1}
\bibitem{ref160} ICAF Iterative Contrastive Alignment Framework for Multimodal Abstractive Summarization. arXiv:2108.05123v3, 2021. \url{https://arxiv.org/abs/2108.05123v3}
\bibitem{ref161} Text-centric Alignment for Multi-Modality Learning. arXiv:2402.08086v1, 2024. \url{https://arxiv.org/abs/2402.08086v1}
\bibitem{ref162} Enhancing Few-shot Text-to-SQL Capabilities of Large Language Models A Study on Prompt Design Strategies. arXiv:2305.12586v1, 2023. \url{https://arxiv.org/abs/2305.12586v1}
\bibitem{ref163} Selective Demonstrations for Cross-domain Text-to-SQL. arXiv:2310.06302v1, 2023. \url{https://arxiv.org/abs/2310.06302v1}
\bibitem{ref164} ACT-SQL In-Context Learning for Text-to-SQL with Automatically-Generated Chain-of-Thought. arXiv:2310.17342v1, 2023. \url{https://arxiv.org/abs/2310.17342v1}
\bibitem{ref165} Divide and Prompt Chain of Thought Prompting for Text-to-SQL. arXiv:2304.11556v1, 2023. \url{https://arxiv.org/abs/2304.11556v1}
\bibitem{ref166} Decomposition for Enhancing Attention Improving LLM-based Text-to-SQL through Workflow Paradigm. arXiv:2402.10671v1, 2024. \url{https://arxiv.org/abs/2402.10671v1}
\bibitem{ref167} How to Prompt LLMs for Text-to-SQL A Study in Zero-shot, Single-domain, and Cross-domain Settings. arXiv:2305.11853v3, 2023. \url{https://arxiv.org/abs/2305.11853v3}
\bibitem{ref168} Going Beyond Word Matching Syntax Improves In-context Example Selection for Machine Translation. arXiv:2403.19285v1, 2024. \url{https://arxiv.org/abs/2403.19285v1}
\bibitem{ref169} In-context Learning as Maintaining Coherency A Study of On-the-fly Machine Translation Using Large Language Models. arXiv:2305.03573v1, 2023. \url{https://arxiv.org/abs/2305.03573v1}
\bibitem{ref170} Metric-Based In-context Learning A Case Study in Text Simplification. arXiv:2307.14632v1, 2023. \url{https://arxiv.org/abs/2307.14632v1}
\bibitem{ref171} Structured Prompting Scaling In-Context Learning to 1,000 Examples. arXiv:2212.06713v1, 2022. \url{https://arxiv.org/abs/2212.06713v1}
\bibitem{ref172} Laying the Foundation First Investigating the Generalization from Atomic Skills to Complex Reasoning Tasks. arXiv:2403.09479v1, 2024. \url{https://arxiv.org/abs/2403.09479v1}
\bibitem{ref173} Data Poisoning for In-context Learning. arXiv:2402.02160v2, 2024. \url{https://arxiv.org/abs/2402.02160v2}
\bibitem{ref174} Hijacking Large Language Models via Adversarial In-Context Learning. arXiv:2311.09948v1, 2023. \url{https://arxiv.org/abs/2311.09948v1}
\bibitem{ref175} Jailbreak and Guard Aligned Language Models with Only Few In-Context Demonstrations. arXiv:2310.06387v1, 2023. \url{https://arxiv.org/abs/2310.06387v1}
\bibitem{ref176} DrAttack Prompt Decomposition and Reconstruction Makes Powerful LLM Jailbreakers. arXiv:2402.16914v2, 2024. \url{https://arxiv.org/abs/2402.16914v2}
\bibitem{ref177} Detecting Language Model Attacks with Perplexity. arXiv:2308.14132v3, 2023. \url{https://arxiv.org/abs/2308.14132v3}
\bibitem{ref178} SafeDecoding Defending against Jailbreak Attacks via Safety-Aware Decoding. arXiv:2402.08983v2, 2024. \url{https://arxiv.org/abs/2402.08983v2}
\bibitem{ref179} Few-Shot Parameter-Efficient Fine-Tuning is Better and Cheaper than In-Context Learning. arXiv:2205.05638v2, 2022. \url{https://arxiv.org/abs/2205.05638v2}
\bibitem{ref180} When Do Prompting and Prefix-Tuning Work A Theory of Capabilities and Limitations. arXiv:2310.19698v2, 2023. \url{https://arxiv.org/abs/2310.19698v2}
\bibitem{ref181} Exploring the Relationship between In-Context Learning and Instruction Tuning. arXiv:2311.10367v1, 2023. \url{https://arxiv.org/abs/2311.10367v1}
\bibitem{ref182} How Does In-Context Learning Help Prompt Tuning. arXiv:2302.11521v1, 2023. \url{https://arxiv.org/abs/2302.11521v1}
\bibitem{ref183} FIAT Fusing learning paradigms with Instruction-Accelerated Tuning. arXiv:2309.04663v2, 2023. \url{https://arxiv.org/abs/2309.04663v2}
\bibitem{ref184} In-context Learning Generalizes, But Not Always Robustly The Case of Syntax. arXiv:2311.07811v2, 2023. \url{https://arxiv.org/abs/2311.07811v2}
\bibitem{ref185} On the Out-Of-Distribution Generalization of Multimodal Large Language Models. arXiv:2402.06599v1, 2024. \url{https://arxiv.org/abs/2402.06599v1}
\bibitem{ref186} Measuring Pointwise \$\textbackslash\{\}mathcal\{V\}\$-Usable Information In-Context-ly. arXiv:2310.12300v2, 2023. \url{https://arxiv.org/abs/2310.12300v2}
\bibitem{ref187} How does Multi-Task Training Affect Transformer In-Context Capabilities Investigations with Function Classes. arXiv:2404.03558v1, 2024. \url{https://arxiv.org/abs/2404.03558v1}
\bibitem{ref188} \$k\$NN Prompting Beyond-Context Learning with Calibration-Free Nearest Neighbor Inference. arXiv:2303.13824v1, 2023. \url{https://arxiv.org/abs/2303.13824v1}
\bibitem{ref189} Misconfidence-based Demonstration Selection for LLM In-Context Learning. arXiv:2401.06301v1, 2024. \url{https://arxiv.org/abs/2401.06301v1}
\bibitem{ref190} Take One Step at a Time to Know Incremental Utility of Demonstration An Analysis on Reranking for Few-Shot In-Context Learning. arXiv:2311.09619v2, 2023. \url{https://arxiv.org/abs/2311.09619v2}
\end{thebibliography}

\end{document}
